\documentclass{article}
\usepackage{enumitem} % For custom lists
\usepackage{geometry} % For better margins
\geometry{a4paper, margin=1in}
\usepackage{appendix}
\usepackage{longtable}
\usepackage{pdflscape}
\usepackage{graphicx}
\usepackage{lscape}  % Corrected landscape package
\usepackage{tabularx}
\usepackage{hyperref}
\usepackage{graphicx}
\usepackage{pgffor} % For looping through images
\usepackage{caption}
\usepackage{booktabs}
\usepackage{siunitx}
\usepackage{geometry}
\geometry{margin=1in}






\begin{document}

	% Title Page
	
	\title{Exploring Household Consumption Patterns and Socioeconomic Disparities: Insights from the HCES Survey (2023-24)}
	\author{Ayan Bashir Sheikh \\ M.Sc Statistics,\\ Department of Statistics, Savitribai Phule Pune University \\Ganeshkhind Pune, Maharashtra}
	\date{May 24, 2025}	
	\maketitle	
	\newpage
		\tableofcontents
		
\newpage
	% Abstract
	\section{Abstract}
	The Household Consumption Expenditure Survey (HCES) provides a comprehensive dataset on household spending patterns across demographic and socio-economic segments. This study leverages HCES 2023-24 data to analyze trends in Monthly Per Capita Expenditure (MPCE), regional disparities, and the effects of inflation, income distribution, and policy interventions. Findings from this study contribute to economic forecasting and policy-making by providing deeper insights into consumption behavior and welfare analysis.
	\newpage

\begin{center}
	

\section*{ Acknowledgements}
\end{center}


I extend my heartfelt gratitude to all those who supported and guided me in the successful completion of this project titled \textit{“Exploring Household Consumption Patterns and Socioeconomic Disparities: Insights from the HCES Survey (2023–24)”}.

I am thankful to the \textbf{Ministry of Statistics and Programme Implementation (MoSPI)}, Government of India, for offering me the opportunity to undertake this internship under the National Internship in Official Statistics (NIOS) 2025 programme. This platform enabled me to engage deeply with real-world official statistics and empirical research using the HCES 2023–24 dataset.

\bigskip
\begin{flushright} 
	\textbf{Ayan Bashir Sheikh} \\
	M.Sc. Statistics \\
	Savitribai Phule Pune University, Pune
\end{flushright}




\newpage
\begin{center}
	\section*{Certificate}
\end{center}

This is to certify that the internship report titled \textbf{“Exploring Household Consumption Patterns and Socioeconomic Disparities: Insights from the HCES Survey (2023–24)”} has been successfully completed and submitted by \textbf{Ayan Bashir Sheikh}, a student of M.Sc. Statistics at Savitribai Phule Pune University, under the National Internship in Official Statistics (NIOS) 2025, organized by the Ministry of Statistics and Programme Implementation (MoSPI), Government of India.

The study was carried out during the internship period from March 24, 2025 to May 24, 2025. This report is an original piece of research based on empirical analysis using the Household Consumption Expenditure Survey (HCES) 2023–24 dataset.

This certificate is issued as a formal recognition of the successful completion of the internship and submission of the project report.

\bigskip

\begin{flushleft}
	\textbf{(Signature)} \\
	\textbf{Authorized Signatory} \\
	Ministry of Statistics and Programme Implementation (MoSPI) \\
	Government of India
\end{flushleft}


\newpage

	% Introduction
	\section{Introduction}
	The Household Consumption Expenditure Survey (HCES) is a statistical initiative designed to collect data on household consumption patterns and expenditure on goods and services. It serves as a key source for analyzing economic well-being, updating the consumer goods basket, and determining Consumer Price Index (CPI) weights. Additionally, HCES plays a crucial role in assessing welfare, poverty, and inequality, enabling policymakers to formulate and monitor economic and social policies tailored to different population segments.
	
	The data collected in HCES provides insights into household living standards, changes in consumption behavior over time, and disparities among various socio-economic groups. Food consumption data derived from HCES is instrumental in food security and nutrition analysis, making it a valuable resource for public health and economic planning.
	
	Following the recommendations of the National Statistical Commission (NSC), the survey was originally scheduled for 2020-21 but was postponed due to the COVID-19 pandemic. To accommodate the need for updated macroeconomic indicators such as Gross Domestic Product (GDP) and CPI, two consecutive surveys were conducted between August 2022 – July 2023 and August 2023 – July 2024. The first survey's findings were published in February 2024, followed by the release of detailed reports and unit-level data in June 2024.
	
	
	
	% Literature Review (Optional, if required)
	\section{Literature Review}
	Household Consumer Expenditure Surveys (HCES) play a crucial role in assessing economic conditions, informing policies, and understanding consumption patterns. Studies highlight their importance in analyzing income distribution, poverty, and household well-being.
	
	Research has focused on improving data collection techniques, enhancing accuracy, and ensuring representativeness. The integration of technology has strengthened survey reliability, while comparisons across regions and years provide insights into economic trends and policy impacts.
	
	HCES aligns with economic theories like Engel's Law, exploring the relationship between household expenditures and macroeconomic factors. They also contribute significantly to global development initiatives, particularly in tracking Sustainable Development Goals (SDGs).
	
	Overall, HCES remains essential for economic research and policy-making, with continuous improvements ensuring its relevance in addressing contemporary economic challenges.
	
	
	
	% Data and Methodology
	\section{Data and Methodology}
	\subsection{Data Source}
	The study utilizes unit-level data from the HCES 2023-24 survey, covering diverse demographic, economic, and consumption variables.
	
	\subsection{Statistical Methods Used}
	The analysis employs descriptive statistics, regression models, and data visualization techniques to explore expenditure patterns and trends. Advanced methods such as clustering and predictive modeling are applied to forecast future consumption behaviors.
	
	

\section{Focus of the Study}

This research delves into key dimensions of household consumption in India, exploring socioeconomic patterns, regional disparities, and the evolving structure of expenditures in response to policy, income, and demographic transitions.

\subsection*{1. Disparities in Household Consumption}

Analyzes Monthly Per Capita Expenditure (MPCE) differences between rural and urban households, highlighting how cost of living, access to services, and lifestyle variation influence spending behavior.

\subsection*{2. Impact of Rising Inflation}

Explores the burden of inflation on household budgets, with emphasis on low-income groups. Examines changes in both essential and discretionary spending as households adapt to price volatility.

\subsection*{3. Income-Based Consumption Patterns}

Studies how different income groups allocate resources across food, housing, education, health, and discretionary items. Highlights contrasts in priorities between low, middle, and high income households.

\subsection*{4. Regional Disparities in Household Spending}

Compares MPCE across Indian states to uncover regional inequalities. Links observed patterns to differences in infrastructure, income levels, and local economic conditions.

\subsection*{5. Socioeconomic Determinants of Expenditure}

Examines how social identity (caste/community), employment status, and household type shape consumption levels. Provides insight into the intersection of social structure and economic capacity.

\subsection*{6. Digital Payments and Consumption Behaviour}

Assesses the rise of digital payment systems in both rural and urban settings, and their influence on how households transact and consume. Reflects broader behavioral shifts linked to financial technology.

\subsection*{7. Welfare Programs and Financial Resilience}
Evaluates the impact of key welfare schemes—including the Public Distribution System (PDS), mid-day meals, and health coverage—on household consumption and financial security.
Emphasis is placed on ration card coverage patterns, including the growing “No Card” population, urban–rural disparities, and state-level anomalies .

\subsection*{8. Forecasting Household Expenditures}

Applies machine learning techniques to predict future MPCE trends. Supports long term planning by identifying emerging patterns in consumption under varied economic scenarios.

\subsection*{9. Food vs. Non-Food Expenditure Trends}

Investigates the shift in consumption priorities from food to non-food items, influenced by rising income, changing aspirations, and urbanization. Tracks evolving household demand structures.

\subsection*{10. Prioritization Within the Consumption Basket}

Analyzes how households balance essential and non-essential expenditures across income groups. Offers a deeper understanding of spending flexibility and constraints.


\subsection*{11. Urbanization and the Changing Food Share}

Explores how urbanization is reshaping food expenditure patterns, particularly among rural households. Captures shifts in dietary preferences and nutritional priorities.

\subsection*{12. Clustering States by Food and Cereal Spending}

Uses cluster analysis to group states based on their food consumption profiles, particularly in cereals and processed foods. Highlights regional dietary trends with implications for food security and agricultural policy.

	
	% Analysis and Findings
	\section{Analysis and Findings}
	
		\begin{longtable}{|l|c|c|c|c|c|c|c|}
			\hline
			\textbf{State/UT} & \multicolumn{2}{c|}{\textbf{FSUs Allotted}} & \multicolumn{2}{c|}{\textbf{FSUs Surveyed}} & \multicolumn{3}{c|}{\textbf{Sample Households}} \\ \cline{2-8}
			& \textbf{Rural} & \textbf{Urban} & \textbf{Rural} & \textbf{Urban} & \textbf{Rural} & \textbf{Urban} & \textbf{Total} \\ \hline
			\endhead
			
			Andhra Pradesh & 361 & 240 & 360 & 239 & 6,306 & 4,159 & 10,465 \\ \hline
			Arunachal Pradesh & 160 & 80 & 160 & 79 & 2,653 & 1,420 & 4,073 \\ \hline
			Assam & 341 & 140 & 337 & 140 & 6,053 & 2,519 & 8,572 \\ \hline
			Bihar & 762 & 200 & 759 & 200 & 13,656 & 3,599 & 17,255 \\ \hline
			Chhattisgarh & 160 & 120 & 160 & 120 & 2,856 & 2,112 & 4,968 \\ \hline
			Delhi & 20 & 200 & 20 & 198 & 310 & 3,049 & 3,359 \\ \hline
			Goa & 20 & 20 & 20 & 20 & 360 & 360 & 720 \\ \hline
			Gujarat & 320 & 320 & 320 & 319 & 5,699 & 5,572 & 11,271 \\ \hline
			Haryana & 160 & 140 & 155 & 139 & 2,785 & 2,478 & 5,263 \\ \hline
			Himachal Pradesh & 82 & 60 & 80 & 60 & 1,355 & 1,051 & 2,406 \\ \hline
			Jharkhand & 221 & 140 & 219 & 140 & 3,866 & 2,465 & 6,331 \\ \hline
			Karnataka & 381 & 340 & 380 & 340 & 6,728 & 5,805 & 12,533 \\ \hline
			Kerala & 220 & 200 & 218 & 200 & 3,883 & 3,517 & 7,400 \\ \hline
			Madhya Pradesh & 482 & 320 & 472 & 319 & 8,426 & 5,679 & 14,105 \\ \hline
			Maharashtra & 660 & 660 & 659 & 652 & 11,629 & 11,117 & 22,746 \\ \hline
			Manipur & 180 & 160 & 91 & 145 & 1,637 & 2,607 & 4,244 \\ \hline
			Meghalaya & 120 & 60 & 120 & 59 & 2,154 & 1,062 & 3,216 \\ \hline
			Mizoram & 80 & 120 & 80 & 120 & 1,440 & 2,160 & 3,600 \\ \hline
			Nagaland & 120 & 60 & 109 & 60 & 1,962 & 1,080 & 3,042 \\ \hline
			Odisha & 381 & 140 & 379 & 140 & 6,768 & 2,445 & 9,213 \\ \hline
			Punjab & 180 & 160 & 178 & 160 & 3,161 & 2,810 & 5,971 \\ \hline
			Rajasthan & 501 & 260 & 500 & 259 & 8,765 & 4,378 & 13,143 \\ \hline
			Sikkim & 80 & 40 & 80 & 40 & 1,439 & 719 & 2,158 \\ \hline
			Tamil Nadu & 420 & 400 & 419 & 397 & 7,484 & 7,008 & 14,492 \\ \hline
			Telangana & 200 & 180 & 199 & 179 & 3,571 & 3,215 & 6,786 \\ \hline
			Tripura & 180 & 100 & 180 & 100 & 3,240 & 1,800 & 5,040 \\ \hline
			Uttar Pradesh & 1,100 & 600 & 1,100 & 600 & 19,663 & 10,644 & 30,307 \\ \hline
			Uttarakhand & 100 & 60 & 100 & 60 & 1,693 & 1,056 & 2,749 \\ \hline
			West Bengal & 600 & 420 & 597 & 419 & 10,706 & 7,414 & 18,120 \\ \hline
			Andaman Nicobar Islands & 40 & 20 & 38 & 20 & 641 & 359 & 1,000 \\ \hline
			Chandigarh & 20 & 20 & 20 & 20 & 359 & 359 & 718 \\ \hline
			Jammu  Kashmir & 100 & 100 & 99 & 100 & 1,782 & 1,800 & 3,582 \\ \hline
			Ladakh & 20 & 20 & 20 & 20 & 360 & 360 & 720 \\ \hline
			Lakshadweep & 14 & 20 & 14 & 20 & 249 & 356 & 605 \\ \hline
			Puducherry & 20 & 40 & 20 & 40 & 359 & 718 & 1,077 \\ \hline
			
			\caption{Number of FSUs allotted and surveyed, and number of households surveyed}
			\label{table:fsu_survey}
		\end{longtable}
		
	
	
	
	\begin{figure}[h]
		\centering
		\includegraphics[width=0.8\textwidth]{D:/internship/Competition/Figure_1.png} % Adjust path & format if needed
		\caption{Household Surveyed Rural vs Urban }
		\label{fig:figure1}
	\end{figure}

\begin{itemize}
	\item \textbf{Highest Surveyed State}: Uttar Pradesh stands out with the largest number of surveyed households (around 30,000), predominantly from rural areas.
	\item \textbf{Balanced Surveys}: States such as Maharashtra and Tamil Nadu show a balanced distribution between rural and urban surveys, suggesting a comprehensive approach to capturing demographic data.
	\item \textbf{Smaller Regions}: Union territories like Lakshadweep and Goa feature the smallest surveyed household counts (below 1,000), reflecting their limited population sizes.
	\item \textbf{Trends in Urban Surveys}: Urban areas in states like Delhi demonstrate significant representation compared to rural areas, highlighting urban-centric demographic and development patterns.
\end{itemize}

\newpage
	
	\begin{longtable}{|l|c|c|c|}
	
		\hline
		\textbf{State/UT} & \textbf{Estimated Rural HH} & \textbf{Estimated Urban HH} & \textbf{Estimated Total HH} \\ \hline
		\endhead
		
		Jammu  Kashmir        & 1,799,259   & 566,668   & 2,365,927 \\ \hline
		Bihar                 & 20,151,145  & 2,167,306  & 22,318,451 \\ \hline
		Sikkim                & 120,725     & 57,498     & 178,223 \\ \hline
		Arunachal Pradesh     & 217,653     & 62,305     & 279,958 \\ \hline
		Nagaland              & 276,945     & 120,338    & 397,283 \\ \hline
		Manipur               & 499,367     & 183,691    & 683,058 \\ \hline
		Mizoram               & 112,840     & 102,746    & 215,586 \\ \hline
		Tripura               & 804,896     & 201,873    & 1,006,768 \\ \hline
		Meghalaya             & 563,295     & 110,882    & 674,177 \\ \hline
		Assam                 & 6,248,139   & 981,145    & 7,229,284 \\ \hline
		West Bengal           & 16,976,256  & 8,004,246  & 24,980,502 \\ \hline
		Himachal Pradesh      & 1,557,010   & 243,984    & 1,800,994 \\ \hline
		Jharkhand             & 5,709,801   & 1,615,380  & 7,325,181 \\ \hline
		Odisha                & 9,384,147   & 2,193,468  & 11,577,615 \\ \hline
		Chhattisgarh          & 4,591,188   & 1,332,890  & 5,924,078 \\ \hline
		Madhya Pradesh        & 11,266,164  & 4,484,483  & 15,750,647 \\ \hline
		Gujarat               & 7,071,053   & 6,551,886  & 13,622,939 \\ \hline
		Dadra  Nagar Haveli  & 54,438      & 131,873    & 186,311 \\ \hline
		Maharashtra           & 14,768,694  & 12,562,136 & 27,330,830 \\ \hline
		Andhra Pradesh        & 9,468,709   & 4,515,617  & 13,984,326 \\ \hline
		Karnataka             & 9,300,573   & 6,880,675  & 16,181,247 \\ \hline
		Punjab                & 3,804,113   & 2,744,353  & 6,548,466 \\ \hline
		Goa                   & 134,511     & 223,872    & 358,383 \\ \hline
		Lakshadweep (U.T.)    & 3,017       & 9,344      & 12,361 \\ \hline
		Kerala                & 4,886,446   & 4,925,091  & 9,811,537 \\ \hline
		Tamil Nadu            & 11,176,960  & 10,686,204 & 21,863,164 \\ \hline
		Puducherry (U.T.)     & 106,244     & 240,645    & 346,889 \\ \hline
		A and N Islands (U.T.) & 76,438     & 50,536     & 126,974 \\ \hline
		Telangana             & 5,689,698   & 5,407,536  & 11,097,234 \\ \hline
		Ladakh (U.T.)         & 43,519      & 9,736      & 53,256 \\ \hline
		Chandigarh (U.T.)     & 11,313      & 242,733    & 254,046 \\ \hline
		Uttarakhand           & 1,856,434   & 754,955    & 2,611,389 \\ \hline
		Haryana               & 3,627,657   & 2,554,196  & 6,181,853 \\ \hline
		Delhi                 & 89,607      & 3,643,667  & 3,733,274 \\ \hline
		Rajasthan             & 10,936,974  & 4,247,935  & 15,184,908 \\ \hline
		Uttar Pradesh         & 30,398,615  & 9,075,344  & 39,473,958 \\ \hline
		All-India             & 193,783,843 & 97,887,235 & 291,671,078 \\ \hline
		
		\caption{Estimated Household Data}
		\label{table:estimated_households}
	\end{longtable}
	

\begin{figure}[h]
	\centering
	\includegraphics[width=1.6\textwidth]{D:/internship/Competition/top5.png} % Adjust path & format if neded
	\caption{Proportion of Rural Vs Urban Households Surveyed}
	\label{fig:\detokenize{\thefigure}}

	\end{figure}

\begin{figure}[h]
	\centering
	\includegraphics[width=1.6\textwidth]{D:/internship/Competition/bottom_5.png} % Adjust path & format if neded
	\caption{Proportion of Rural Vs Urban Households Surveyed}
	\label{fig:\detokenize{\thefigure}}
	
\end{figure}
 States with the highest estimated number of households, emphasizing regions of demographic density.

\begin{itemize}
	\item \textbf{Top Highlighted States}: Uttar Pradesh, Bihar, West Bengal, Maharashtra, and Madhya Pradesh stand out due to their large household estimates, indicating their significant share in the national population.
	\item \textbf{Color Gradient Representation}: A gradient from lighter to darker shades represents the variation in household numbers across regions, with the highlighted states falling into the higher bracket.
	\item \textbf{Range of Values}: The estimates range significantly, with the highlighted states surpassing a certain threshold, underlining their demographic prominence.
\end{itemize}






	\begin{itemize}
	\item \textbf{Lowest Household State}: Sikkim stands out as the state with the lowest estimated number of households, which likely correlates with it having the smallest population in the country.
	\item \textbf{Highlighted States}: Other northeastern states, such as Arunachal Pradesh, Nagaland, Mizoram, and Meghalaya, are also identified with minimal household estimates, reflecting lower population densities.	

	
	

		\item \textbf{Color Gradient Representation}: The map employs a gradient, transitioning from green (representing higher household counts) to red (indicating the lowest estimates). Sikkim falls at the extreme lower end of the gradient.
		\item \textbf{Implications}: These insights suggest regions with lower demographic density, which can help prioritize resource distribution and infrastructure development in sparsely populated areas.
	\end{itemize}
	

\newpage
%\newpage
	\begin{figure}[h]
	\centering
	\includegraphics[width=0.5\textwidth]{D:/internship/Competition/Figure_3.png} % Adjust path & format if needed
	\caption{Proportion of Rural Vs Urban Households Surveyed}
	\label{fig:\detokenize{\thefigure}}
	
	\end{figure}
	\begin{itemize}
		\item \textbf{Rural Majority}: The chart shows that 58.9\% of the surveyed households are from rural areas, indicated by the blue section. This reflects a higher focus on rural demographics in the survey.
		\item \textbf{Urban Minority}: The orange section, representing urban households, accounts for 41.1\%. While smaller in proportion, urban surveys still constitute a significant share.
		\item \textbf{Demographic Insight}: The chart highlights the predominance of rural households in the survey, suggesting that rural areas are prioritized or have higher household counts compared to urban regions.
	\end{itemize}



	\begin{figure}[h]
	\centering
	\includegraphics[width=0.8\textwidth]{D:/internship/Competition/Figure_4.png} % Adjust path & format if needed
	\caption{Surveyed FSUs: rural Vs. Urban Across States}
\label{fig:\detokenize{\thefigure}}

\end{figure}
	\begin{itemize}
		\item \textbf{Rural Dominance}: The graph indicates that in the majority of states, more rural FSUs were surveyed compared to urban FSUs. This highlights a stronger emphasis on rural areas in the survey.
		\item \textbf{State-wise Peaks}: Notable states with significantly high rural FSU surveys include Bihar, Maharashtra, Uttar Pradesh, and West Bengal, reflecting their large rural population bases.
		\item \textbf{Urban Representation}: While urban FSUs are fewer in most states, their consistent inclusion ensures the survey captures urban demographics for a balanced perspective.
		\item \textbf{Regional Variation}: The variation in surveyed FSUs across states emphasizes demographic differences and the tailored approach of the survey to capture diverse data.
	\end{itemize}
	
	
\section*{ Survey Completion Rates}

\begin{itemize}
	\item \textbf{High Completion Rates Nationwide}: Most states and union territories exhibit high survey completion rates for both rural and urban areas, reflecting robust data collection efforts.
	\item \textbf{Variations Noted}: Nagaland and Manipur stand out, with noticeable disparities between rural and urban completion rates. Manipur shows a significant gap, with urban areas achieving higher completion rates than rural areas.

	\begin{figure}[h!]
	\centering
	\includegraphics[width=0.8\textwidth]{D:/internship/Competition/Figure_6.png} % Adjust path & format if needed
	\caption{Survey Completion Rates}
	\label{fig:\detokenize{\thefigure}}
	
\end{figure}
	
	\item \textbf{Uniform Patterns Elsewhere}: Other states such as Gujarat, Maharashtra, and Karnataka display almost identical completion rates in both rural and urban surveys, suggesting uniformity in survey execution.
	\end{itemize}
	
	
	\newpage
	\section*{ Distribution of Survey Completion Rates}
	
	\begin{itemize}
		\item \textbf{Median Completion Rates}: Both rural and urban groups demonstrate high median completion rates close to 100\%, indicating successful survey execution in most regions.
		\item \textbf{Outliers in Rural Group}: Rural survey completion rates feature outliers below 90\%, with a significant outlier close to 50\%. This suggests inconsistencies or challenges in completing surveys in some rural areas.
	
		\begin{figure}[h]
		\centering
		\includegraphics[width=0.8\textwidth]{D:/internship/Competition/Figure_7.png} % Adjust path & format if needed
		\caption{Distribution of Survey Completion Rates}
		\label{fig:\detokenize{\thefigure}}
		
	\end{figure}	

		\item \textbf{Uniformity in Urban Group}: Urban survey completion rates are tightly clustered around the median, with fewer and less extreme outliers compared to rural areas, reflecting more consistency in urban data collection.
	\item \textbf{Comparison}: The plot highlights that while both groups achieve high completion rates overall, rural areas exhibit greater variability and face unique survey completion challenges.
	\end{itemize}
	
	




\section*{Household Coverage Per FSU - Rural vs Urban}

\begin{itemize}
	\item \textbf{Rural Majority in Coverage:} Most states display higher household coverage per FSU in rural
	areas (blue bars) compared to urban areas (orange bars). This trend aligns with the rural household
	dominance across the country.
	\item \textbf{State-Level Insights}: Uttar Pradesh, Bihar, and Maharashtra have notably high rural coverage
	rates. On the urban side, states like Delhi and Goa showcase relatively balanced or elevated
	coverage.
	\item \textbf{Contrast Across Regions:} Northeastern states (e.g., Arunachal Pradesh, Mizoram) show relatively
	lower overall coverage, reflecting smaller population sizes and demographic concentration in
	rural zones.
	
		\begin{figure}[h]
		\centering
		\includegraphics[width=0.8\textwidth]{D:/internship/Competition/Figure_8.png} % Adjust path & format if needed
		\caption{Household Coverage Per FSU - Rural vs Urban}
		\label{fig:\detokenize{\thefigure}}
		
	\end{figure}	
	
	\item \textbf{Range of Coverage:} The household coverage per FSU spans from less than 5 in smaller states
	like Goa to more than 17 in states like Uttar Pradesh, indicating regional variations in survey
	intensity.
\end{itemize}



\section*{Relationship Between FSUs Surveyed and Households Surveyed}

\begin{itemize}
	

\item \textbf{Positive Correlation:} The scatter plot demonstrates a clear positive linear relationship between
the total FSUs (Field Survey Units) surveyed and the total households surveyed. As the number
of FSUs surveyed increases, the households surveyed increase proportionally.
\item   textbf{Clustered Distribution:} The points are well-clustered along the trend line, indicating consistency
in data collection efforts across various FSUs.
\item \textbf{Data Variability:} Some points deviate slightly from the main trend, suggesting variability in
household coverage per FSU in certain regions.



\item\textbf{ Categorical Indicators:} The data points are color-coded (e.g., red, blue, light blue), potentially
representing different states or territories. These color distinctions might correlate with unique
regional characteristics or demographics.

\end{itemize}

	\begin{figure}[h]
	\centering
	\includegraphics[width=0.8\textwidth]{D:/internship/Competition/Figure_5.png} % Adjust path & format if needed
	\caption{Relationship Between FSUs Surveyed and Households Surveyed}
	\label{fig:\detokenize{\thefigure}}
	
\end{figure}	


\section*{Elbow Method}

\begin{itemize}
	\item  \textbf{Sharp Decrease in WCSS:} The plot shows a rapid decline in the Within-Cluster Sum of Squares
	(WCSS) as the number of clusters (k) increases from 1 to 2, and then from 2 to 3, indicating a
	better-defined clustering structure with additional clusters.
		\begin{figure}[h]
		\centering
		\includegraphics[width=0.8\textwidth]{D:/internship/Competition/Figure_12.png} % Adjust path & format if needed
		\caption{Elbow Method}
		\label{fig:\detokenize{\thefigure}}
		
	\end{figure}	
	\item \textbf{Optimal Number of Clusters:} The ”elbow” is observed at around k = 3 or k = 4, where the
	rate of decrease in WCSS begins to slow down. This suggests that the ideal number of clusters is 3
	or 4, as adding more clusters beyond this point results in minimal gains in clustering performance.
	
	
	\item \textbf{Diminishing Returns:} Beyond the elbow point, further increases in k lead to diminishing returns,
	as the reduction in WCSS becomes negligible. This implies that additional clusters would add
	complexity without significant improvement in model efficiency.
	
\end{itemize}

\newpage
\section*{Cluster-wise Average FSUs \& Households Surveyed}


\begin{itemize}
	\item \textbf{Dominance of Rural Households:} Across all clusters, rural households surveyed consistently
	demonstrate higher counts compared to urban households. This underscores the survey’s focus on
	rural demographics.
	\item \textbf{Cluster 0 and Cluster 3:} Both clusters exhibit significantly high counts for surveyed rural and
	urban households, with rural numbers dominating. This indicates high survey activity in these
	clusters, possibly due to higher population density or prioritization.
	
	\item \textbf{Cluster 1 -} Minimal Survey Coverage: This cluster displays the lowest counts across all
	metrics, suggesting regions with limited survey activities.
	
	
	\begin{figure}[h]
		\centering
		\includegraphics[width=0.8\textwidth]{D:/internship/Competition/Figure_14.png} % Adjust path & format if needed
		\caption{Cluster-wise Average FSUs \& Households Surveyed}
		\label{fig:\detokenize{\thefigure}}
		
	\end{figure}
	
	\item \textbf{Moderate Survey Coverage in Cluster 2:} This cluster reflects moderate activity with noticeable
	disparities between rural and urban households surveyed, again emphasizing rural focus.
	
	\item \textbf{FSUs Patterns:} Similar trends are observed for FSUs surveyed, with rural FSUs showing consistently
	higher averages compared to urban FSUs across most clusters.
	
\end{itemize}

\newpage
\section*{Cluster-wise Survey Averages}

\begin{itemize}
	\item \textbf{Cluster 0 - High Averages}: Rural metrics (492.4 and 8780.6) significantly exceed urban metrics
	(321.9 and 5636.4), reflecting a predominant rural survey focus. Urban metrics are still notable
	but consistently lower, indicating demographic differences in surveyed averages.
	
	
	\item \textbf{Cluster 1 - Low Coverage:} Both rural (64.2 and 1123.9) and urban (63.6 and 1109.4) averages
	are the lowest among all clusters. This suggests minimal survey activity, potentially due to sparse
	demographics or logistical challenges.
	
		\begin{figure}[h!]
		\centering
		\includegraphics[width=0.8\textwidth]{D:/internship/Competition/Figure_15.png} % Adjust path & format if needed
		\caption{Cluster-wise Survey Averages}
		\label{fig:\detokenize{\thefigure}}
		
	\end{figure}
	
	\item \textbf{Cluster 2 - Moderate Metrics:} Rural averages (238.5 and 4248.9) are again higher than urban
	ones (155.7 and 2752.0), maintaining the trend of rural emphasis. Moderate survey values imply
	balanced focus but lower priority compared to Cluster 0 and Cluster 3.
	
	
	\item \textbf{Cluster 3 - Highest Survey Effort:} The highest averages for both rural (879.5 and 15646.0) and
	urban (626.0 and 10880.5) metrics, highlighting substantial survey activities. This cluster reflects
	regions with the densest population or strategic importance in the survey.
	\item \textbf{Overall Trend:} Across all clusters, rural averages consistently exceed urban averages, emphasizing
	rural demographics as a priority in the survey.
	
\end{itemize}

\newpage
\section*{Household Surveys Across Clusters}

\begin{itemize}
	

\item \textbf{Cluster with the Highest Median:} Cluster 3 exhibits the highest median number of households
surveyed, indicating a greater concentration of survey efforts in this cluster.
\item \textbf{Declining Median Values Across Clusters:} After Cluster 3, the median values decrease sequentially
for Cluster 0, Cluster 2, and Cluster 1, reflecting varying survey intensities across these
clusters.
\item \textbf{Cluster 1 - Lowest Activity:} Cluster 1 not only has the lowest median but also a smaller range
and interquartile spread, suggesting limited household survey coverage in this cluster.
\item \textbf{Cluster 3 - Wide Range and High Spread:} Cluster 3 features the widest range and interquartile
spread, highlighting greater variability in the number of households surveyed within
this cluster.
\item \textbf{Overall Distribution:} The plot demonstrates significant differences in survey intensities across
clusters, with Cluster 3 standing out as a focal point for survey activity and Cluster 1 being
relatively less prioritized.


\end{itemize}


	\begin{figure}[ht]
	\centering
	\includegraphics[width=0.8\textwidth]{D:/internship/Competition/Figure_16.png} % Adjust path & format if needed
	\caption{Household Surveys Across Clusters}
	\label{fig:\detokenize{\thefigure}}
	
\end{figure}

\newpage
\section*{Cluster-wise Survey Patterns (Radar Chart)}

\begin{itemize}
	
\item \textbf{Cluster 3 - Dominant Coverage:} The red section of the radar chart, representing Cluster 3,
spans the largest area across all axes. This indicates the highest values in surveyed rural households,
urban households, and urban FSUs, highlighting extensive survey coverage in this cluster.
\item \textbf{Cluster 0 - Moderate Coverage:} The blue area, representing Cluster 0, shows moderate values
across the axes. While it has a balanced presence, its coverage is smaller compared to Cluster 3,
particularly in urban FSUs.
\item \textbf{Cluster 1 - Minimal Survey Data:} The orange region, representing Cluster 1, is the smallest
across the axes, emphasizing the least survey activity in terms of rural households, urban households,
and urban FSUs.

\item \textbf{Cluster 2 - Mixed Patterns:} The green section, representing Cluster 2, reflects mixed patterns
with values higher than Cluster 1 but lower than Clusters 0 and 3. It indicates moderate survey
efforts.

\item \textbf{Overall Trends:} The radar chart visually compares survey patterns and highlights significant
differences in survey coverage across clusters. Cluster 3 emerges as the focal point with the highest
survey intensity, while Cluster 1 exhibits minimal efforts.
\end{itemize}

	\begin{figure}[ht]
	\centering
	\includegraphics[width=0.8\textwidth]{D:/internship/Competition/Figure_17.png} % Adjust path & format if needed
	\caption{Cluster-wise Survey Patterns (Radar Chart)}
	\label{fig:\detokenize{\thefigure}}
	
   \end{figure}







\newpage

	
\section{Monthly Per Capita Expenditure (MPCE)}









	
	\begin{longtable}{|c|l|c|c|}
		\hline
		\textbf{State Code} & \textbf{State Name} & \textbf{Rural MPCE} & \textbf{Urban MPCE} \\ \hline
		\endhead
		
		01  & Jammu  Kashmir              & 4774  & 6327 \\ \hline
		02  & Himachal Pradesh             & 5825  & 9223 \\ \hline
		03  & Punjab                       & 5817  & 7359 \\ \hline
		04  & Chandigarh                   & 8857  & 13425 \\ \hline
		05  & Uttarakhand                  & 5003  & 7486 \\ \hline
		06  & Haryana                      & 5377  & 8427 \\ \hline
		07  & Delhi                        & 7400  & 8534 \\ \hline
		08  & Rajasthan                    & 4510  & 6574 \\ \hline
		09  & Uttar Pradesh                & 3481  & 5395 \\ \hline
		10  & Bihar                        & 3670  & 5080 \\ \hline
		11  & Sikkim                       & 9377  & 13927 \\ \hline
		12  & Arunachal Pradesh            & 5995  & 9832 \\ \hline
		13  & Nagaland                     & 5155  & 8022 \\ \hline
		14  & Manipur                      & 4531  & 5945 \\ \hline
		15  & Mizoram                      & 5963  & 8709 \\ \hline
		16  & Tripura                      & 6259  & 8034 \\ \hline
		17  & Meghalaya                    & 3852  & 7839 \\ \hline
		18  & Assam                        & 3793  & 6794 \\ \hline
		19  & West Bengal                  & 3620  & 5775 \\ \hline
		20  & Jharkhand                    & 2946  & 5393 \\ \hline
		21  & Odisha                       & 3357  & 5825 \\ \hline
		22  & Chhattisgarh                 & 2739  & 4927 \\ \hline
		23  & Madhya Pradesh               & 3441  & 5538 \\ \hline
		24  & Gujarat                      & 4116  & 7175 \\ \hline
		25  & D  N. Haveli  Daman  Diu  & 4311  & 6837 \\ \hline
		26  & Maharashtra                  & 4145  & 7363 \\ \hline
		27  & Andhra Pradesh               & 5327  & 7182 \\ \hline
		28  & Karnataka                    & 4903  & 8076 \\ \hline
		29  & Goa                          & 8048  & 9726 \\ \hline
		30  & Lakshadweep                  & 6350  & 6377 \\ \hline
		31  & Kerala                       & 6611  & 7783 \\ \hline
		32  & Tamilnadu                    & 5701  & 8165 \\ \hline
		33  & Puduchery                    & 7598  & 8637 \\ \hline
		34  & Andaman N. Island          & 7771  & 10453 \\ \hline
		35  & Telangana                    & 5435  & 8978 \\ \hline
		36  & Ladakh                       & 5010  & 7533 \\ \hline
		
		\caption{Monthly Per Capita Expenditure (MPCE) across Indian States (Rural vs. Urban)}
		\label{table:mpce_states}
	\end{longtable}
	






\subsection{ MPCE Distribution in Rural vs. Urban Areas }





\begin{itemize}
	\item \textbf{Wider Distribution for Urban Areas}: The orange violin plot for urban areas exhibits a broader range, with MPCE extending beyond 12,000Rs. This indicates a more diverse expenditure pattern among urban households.
	\item \textbf{Higher Median in Urban Areas}: The box plot within the urban violin highlights a higher median MPCE compared to rural areas, suggesting that urban households, on average, spend more.


\begin{figure}[h]
	\centering
	\includegraphics[width=0.9\textwidth]{D:/internship/Competition/Figure_20.png} % Adjust path & format if needed
	\caption{ MPCE Distribution in Rural vs. Urban Areas (Violin Plot) }
	\label{fig:\detokenize{\thefigure}}
\end{figure}	

	\item \textbf{Tighter Distribution in Rural Areas}: The blue violin plot for rural areas shows a narrower range, with MPCE largely concentrated below 8,000Rs. This reflects a more consistent expenditure pattern among rural households.
\item \textbf{Overlap in Lower Ranges}: Both plots overlap in the lower expenditure ranges (0Rs–4,000Rs), indicating similarity in spending among the lowest-income groups in both demographics.
\item \textbf{Policy Implications}: The distinct patterns emphasize the need for tailored economic policies to address the disparities in expenditure between rural and urban households.
\end{itemize}




\newpage
\section{ Trends in Rural vs. Urban MPCE Across Fractile Classes}





	
	\begin{longtable}{|c|c|c|c|}
		\hline
		\textbf{Fractile} & \textbf{Total Expenditure} & \textbf{Household Size} & \textbf{MPCE} \\ \hline
		\endhead
		
		0-5     & 71,657,262,868   & 42,736,481   & 1677 \\ \hline
		5-10    & 90,833,408,181   & 42,726,577   & 2126 \\ \hline
		10-20   & 211,403,024,803  & 85,480,531   & 2473 \\ \hline
		20-30   & 242,163,622,787  & 85,472,417   & 2833 \\ \hline
		30-40   & 270,266,748,526  & 85,476,775   & 3162 \\ \hline
		40-50   & 298,921,699,477  & 85,467,313   & 3497 \\ \hline
		50-60   & 330,473,715,843  & 85,480,935   & 3866 \\ \hline
		60-70   & 367,858,056,668  & 85,472,525   & 4304 \\ \hline
		70-80   & 417,553,308,080  & 85,472,218   & 4885 \\ \hline
		80-90   & 492,590,169,238  & 85,475,974   & 5763 \\ \hline
		90-95   & 296,107,354,997  & 42,736,651   & 6929 \\ \hline
		95-100  & 433,204,423,242  & 42,737,029   & 10137 \\ \hline
		
		\caption{Fractile-wise Distribution of Total Expenditure, Household Size, and MPCE for Rural Sector}
		\label{table:fractile_rural}
	\end{longtable}
	


	
	\begin{longtable}{|c|c|c|c|}
		\hline
		\textbf{Fractile} & \textbf{Total Expenditure} & \textbf{Household Size} & \textbf{MPCE} \\ \hline
		\endhead
		
		0-5     & 43,107,834,397   & 18,144,642   & 2376 \\ \hline
		5-10    & 56,120,204,278   & 18,144,637   & 3093 \\ \hline
		10-20   & 133,811,935,916  & 36,289,896   & 3687 \\ \hline
		20-30   & 157,985,265,930  & 36,291,215   & 4353 \\ \hline
		30-40   & 180,682,200,884  & 36,291,854   & 4979 \\ \hline
		40-50   & 204,013,920,066  & 36,290,695   & 5622 \\ \hline
		50-60   & 229,856,488,092  & 36,290,184   & 6334 \\ \hline
		60-70   & 261,203,938,231  & 36,283,291   & 7199 \\ \hline
		70-80   & 303,117,438,687  & 36,292,326   & 8352 \\ \hline
		80-90   & 367,994,935,413  & 36,297,493   & 10138 \\ \hline
		90-95   & 232,543,354,883  & 18,144,072   & 12816 \\ \hline
		95-100  & 368,559,422,756  & 18,147,728   & 20309 \\ \hline
		
		\caption{Fractile-wise Distribution of Total Expenditure, Household Size, and MPCE (Urban)}
		\label{table:fractile_urban}
	\end{longtable}
	


\begin{itemize}
	\item \textbf{General Trend}: Both Rural and Urban MPCE show an increasing trend across the fractile classes. This indicates that as fractile classes increase (representing higher economic strata), the average Monthly Per Capita Expenditure also rises.
	\item \textbf{Urban MPCE Consistently Higher}: The orange line (Urban MPCE) remains consistently above the blue line (Rural MPCE) for all fractile classes. This reflects the higher cost of living and spending capacity in urban areas compared to rural areas.

\begin{figure}[h]
	\centering
	\includegraphics[width=0.9\textwidth]{D:/internship/Competition/Figure_19.png} % Adjust path & format if needed
	\caption{  Trends in Rural vs. Urban MPCE Across Fractile Classes}
	\label{fig:\detokenize{\thefigure}}
\end{figure}	


	\item \textbf{Wide Gap in Higher Fractile Classes}: The gap between Rural and Urban MPCE widens significantly in the higher fractile classes (e.g., 90–100\%), signifying greater inequality in expenditure levels among wealthier groups in rural and urban demographics.
\item \textbf{Policy Implications}: These trends suggest the need for targeted policies addressing the disparity in rural and urban expenditure patterns, especially for high-income groups. Additionally, efforts might focus on improving spending capacity in rural areas to narrow the gap.
\end{itemize}




\newpage

\subsection{ Rural vs. Urban MPCE }

\begin{itemize}
	\item \textbf{Higher Urban Expenditure}: The chart indicates that Urban MPCE (Monthly Per Capita Expenditure) exceeds Rural MPCE across most percentage ranges, with the orange dashed line extending further outward compared to the blue solid line in the radar chart.
	\item \textbf{Expenditure Patterns}: Urban areas demonstrate a consistently higher spending pattern, reflecting differences in cost of living, consumption habits, and availability of goods and services compared to rural areas.
	\item \textbf{Wide Gap in Higher Ranges}: The disparity becomes more prominent in higher percentage ranges (70\%-90\%), indicating that the urban population in these ranges spends substantially more than their rural counterparts.


\begin{figure}[h]
	\centering
	\includegraphics[width=\textwidth]{D:/internship/Competition/Figure_18.png} % Adjust path & format if needed
	\caption{ Rural vs. Urban MPCE (Radar Chart) }
	\label{fig:\detokenize{\thefigure}}
\end{figure}	

	\item \textbf{Narrower Gap in Lower Ranges}: At lower percentage ranges (0\%-10\%), the gap between Rural and Urban MPCE is relatively narrow, suggesting similar expenditure levels among lower-income groups in both regions.
\item \textbf{Policy Implications}: This comparison underscores the need for tailored policy interventions to address the unique expenditure patterns and cost of living disparities between rural and urban populations.
\end{itemize}




\subsection{ Share of Rural and Urban MPCE Across Fractile Classes}




\begin{itemize}
	\item \textbf{Increasing MPCE Across Fractile Classes}: The bar chart shows a steady increase in total Monthly Per Capita Expenditure (MPCE) as we move from lower fractile classes (0--5\%) to higher ones (95--100\%). This suggests that households in higher economic strata spend significantly more on average.
	\item \textbf{Urban MPCE Dominance}: For every fractile class, the urban MPCE (orange segments) surpasses the rural MPCE (blue segments), reflecting higher expenditure patterns in urban households due to factors like increased income, cost of living, and diverse consumption.
	
	
	
	\begin{figure}[h]
		\centering
		\includegraphics[width=0.9\textwidth]{D:/internship/Competition/Figure_21.png} % Adjust path & format if needed
		\caption{Share of Rural and Urban MPCE Across Fractile Classes }
		\label{fig:\detokenize{\thefigure}}
	\end{figure}	
	
	
	\item \textbf{Proportional Increase in MPCE}: Both rural and urban MPCE rise proportionally across fractile classes, but the gap between urban and rural MPCE widens as we move toward higher classes, emphasizing greater inequality in expenditure among wealthier households.
	\item \textbf{Fractile Representation}: Each fractile class is well-represented, providing insights into expenditure trends across various economic segments of the population.
	\item \textbf{Policy Implications}: The consistent gap between rural and urban MPCE suggests the need for targeted policies aimed at reducing expenditure disparity and improving economic opportunities in rural areas.
\end{itemize}



\section*{ MPCE Distribution Across Fractile Classes}



\begin{itemize}
	\item \textbf{Increasing Trend Across Fractile Classes}: The box plots indicate that the Monthly Per Capita Expenditure (MPCE) values rise steadily from lower fractile classes (0--5\%) to higher ones (95--100\%). This trend reflects increasing expenditure with higher income levels.
	\item \textbf{Urban MPCE Dominance}: Yellow box plots (urban sectors) consistently display higher MPCE values across all fractile classes compared to blue box plots (rural sectors). This highlights the disparity in spending between rural and urban households.


\begin{figure}[h]
	\centering
	\includegraphics[width=0.9\textwidth]{D:/internship/Competition/Figure_22.png} % Adjust path & format if needed
	\caption{ MPCE Distribution Across Fractile Classes}
	\label{fig:\detokenize{\thefigure}}
\end{figure}	



	\item \textbf{Greater Variability in Urban Areas}: Urban MPCE values exhibit wider boxes and whiskers, especially in higher fractile classes, suggesting greater variability in spending patterns among urban households compared to rural ones.
\item \textbf{Smaller Range in Rural Areas}: Rural MPCE values are more concentrated with narrower boxes and whiskers, reflecting less variability in expenditure across different income groups in rural areas.
\item \textbf{Policy Implications}: These findings emphasize the need to address expenditure disparities through policies aimed at improving economic conditions in rural areas and reducing inequality between rural and urban households.
\end{itemize}



\section*{Cumulative Distribution of MPCE for Rural vs. Urban Areas}

\begin{itemize}
	\item \textbf{Distinct Expenditure Patterns}: The blue dashed line (representing rural MPCE) and the orange solid line (representing urban MPCE) highlight distinct expenditure patterns between rural and urban areas.
	\item \textbf{Urban Areas Show Higher MPCE Values}: The urban CDF (orange line) consistently reaches higher MPCE values than the rural CDF (blue line), reflecting greater per capita spending in urban areas across most income levels.



\begin{figure}[h]
	\centering
	\includegraphics[width=\textwidth]{D:/internship/Competition/Figure_23.png} % Adjust path & format if needed
	\caption{Cumulative Distribution of MPCE for Rural vs. Urban Areas}
	\label{fig:\detokenize{\thefigure}}
\end{figure}	


	\item \textbf{Steeper Growth in Urban MPCE at Lower Values}: At MPCE values below 5,000Rs, the urban CDF exhibits steeper growth compared to the rural CDF, suggesting a larger proportion of urban households in this expenditure range.
\item \textbf{Converging Trends at Higher MPCE Values}: Both distributions converge as MPCE values approach 30,000Rs, indicating similar expenditure behaviors among the wealthiest households in both demographics.
\item \textbf{Policy Implications}: These findings underscore the need for policies targeting expenditure disparities, with a focus on uplifting rural household spending capacity while addressing income inequality.
\end{itemize}


\newpage
\section{ MPCE Growth Rate (Rural vs Urban) }

\begin{itemize}
	\item \textbf{Higher Growth in Rural Areas at Extremes}: Rural areas (blue bars) show a higher MPCE growth rate for the lowest income percentiles (0--5\%) and the highest income percentiles (95--100\%). This indicates targeted economic interventions or natural growth patterns benefiting these specific segments.
	\item \textbf{Urban Areas Dominate Middle Percentiles}: Urban areas (orange bars) exhibit higher MPCE growth rates in the middle-income percentiles (10--90\%), which may reflect increased opportunities and consumption trends in urban centers for this demographic.

\begin{figure}[h]
	\centering
	\includegraphics[width=\textwidth]{D:/internship/Competition/Figure_24.png} % Adjust path & format if needed
	\caption{MPCE Growth Rate (Rural vs Urban)}
	\label{fig:\detokenize{\thefigure}}
\end{figure}	



	\item \textbf{Broad Trends in Rural and Urban Growth}: Rural MPCE growth displays a more polarized distribution, with significant growth at the extremes, while urban MPCE growth is more consistent across the income percentiles.
	\item \textbf{Policy Implications}: The observed trends suggest the need for strategies aimed at sustaining rural growth across all income percentiles and managing urban growth disparities, especially for low-income and high-income households.
\end{itemize}




\subsection{High-Growth vs. Low-Growth Fractile Classes}

\begin{itemize}
	\item \textbf{Distinct Patterns Between Rural and Urban Growth}: Blue dots (Rural Growth) and orange dots (Urban Growth) exhibit distinct growth patterns across fractile classes. Urban Growth generally surpasses Rural Growth for mid-range fractile classes (10--90\%), reflecting disparities in growth dynamics.
	\item \textbf{High-Growth Fractile Classes}: Rural and Urban Growth rates are significantly higher in the lowest (0--5\%) and highest (95--100\%) fractile classes, possibly driven by targeted economic interventions or natural economic expansions at these extremes.



\begin{figure}[h]
	\centering
	\includegraphics[width=\textwidth]{D:/internship/Competition/Figure_25.png} % Adjust path & format if needed
	\caption{High-Growth vs. Low-Growth Fractile Classes}
	\label{fig:\detokenize{\thefigure}}
\end{figure}	



	\item \textbf{Mid-Range Fractile Classes}: Urban Growth maintains relatively steady and dominant growth rates in the middle fractile classes, whereas Rural Growth fluctuates, highlighting greater variability and challenges in sustaining consistent rural growth.
	\item \textbf{Policy Implications}: The analysis underscores the need for policies that address growth variability in rural areas while sustaining urban growth. Interventions should focus on improving economic opportunities in rural mid-range fractile classes to bridge the growth gap.
\end{itemize}







\subsection{ MPCE Growth Rate (Rural vs Urban)}

\begin{itemize}
	\item \textbf{Regional Differences in Growth}: States like Odisha, Rajasthan, and Chhattisgarh show exceptionally high MPCE growth rates in rural areas, reflecting economic improvements or increased spending capacity among rural households in these regions. Conversely, states like Punjab exhibit negative rural growth, suggesting challenges or reduced expenditure trends in rural sectors.
	\item \textbf{Urban Growth Dominance in Some States}: Urban growth rates surpass rural growth rates in states such as Andhra Pradesh, Gujarat, Maharashtra, and Haryana. This trend highlights urban economic dynamism and the potential impact of urban-focused policies.


\begin{figure}[h]
	\centering
	\includegraphics[width=\textwidth]{D:/internship/Competition/Figure_28.png} % Adjust path & format if needed
	\caption{ MPCE Growth Rate (Rural vs Urban)}
	\label{fig:\detokenize{\thefigure}}
\end{figure}	


	\item \textbf{All-India Trends}: At the "All-India" level, urban MPCE growth outpaces rural growth, indicating overall stronger expenditure trends in urban areas.
	\item \textbf{Policy Implications}: The wide disparities in rural and urban growth across states suggest the need for customized policy interventions. Regions with declining rural growth, such as Punjab, may require targeted support, while high-growth states might focus on sustaining momentum.
\end{itemize}





\subsection{ Urban-to-Rural MPCE Ratio Analysis}



\begin{itemize}
	\item \textbf{Kerala Leads in Disparity}: With a ratio of approximately 2.6, Kerala shows the most significant difference in Monthly Per Capita Expenditure (MPCE) between urban and rural areas. This indicates that urban households in Kerala spend considerably more than their rural counterparts.
	\item \textbf{Higher Ratios Across Developed States}: States like Maharashtra, Gujarat, and Tamil Nadu also exhibit higher ratios (above 2.0), reflecting similar urban-rural disparities, likely driven by economic opportunities and cost of living in urban areas.
	



\begin{figure}[h]
	\centering
	\includegraphics[width=\textwidth]{D:/internship/Competition/Figure_29.png} % Adjust path & format if needed
	\caption{ Urban-to-Rural MPCE Ratio Analysis}
	\label{fig:\detokenize{\thefigure}}
\end{figure}	



\item \textbf{Relatively Balanced Ratios in Some States}: States like Bihar, Chhattisgarh, and Uttar Pradesh show ratios closer to 1.5, suggesting more balanced expenditure patterns between rural and urban households.
\item \textbf{All-India Average}: The All-India Urban-to-Rural MPCE ratio lies around 2.0, indicating that, on average, urban households spend twice as much as rural ones across the country.
\item \textbf{Policy Implications}: The data highlights the need for targeted policies aimed at reducing urban-rural expenditure disparities, particularly in states with higher ratios like Kerala, while promoting rural development and economic opportunities.
\end{itemize}







\section*{ MPCE Comparison Heatmap (2023-24)}

\begin{itemize}
	\item \textbf{Highest Urban MPCE (Kerala)}: Kerala's urban MPCE value of Rs.9302 significantly exceeds its rural MPCE value of Rs.3611, showcasing the most pronounced urban-rural expenditure disparity across states.
	\item \textbf{Balanced States}: Bihar and Chhattisgarh display relatively balanced MPCE values, with rural MPCE values closer to urban ones compared to other states. Bihar's values are Rs.3426 (rural) and Rs.4792 (urban), while Chhattisgarh's are Rs.2796 (rural) and Rs.4781 (urban).
	\item \textbf{Higher Urban MPCE Nationwide}: Across all states, urban MPCE consistently surpasses rural MPCE, indicating higher expenditure trends in urban areas due to economic opportunities and cost of living differences.
	\item \textbf{All-India Average}: The All-India average for rural MPCE is Rs.4122, while urban MPCE stands at Rs.6996, highlighting a nationwide disparity between rural and urban expenditures.
	\item \textbf{Policy Implications}: These patterns suggest the need for targeted policies aimed at reducing urban-rural disparities and improving economic conditions in rural areas, particularly in states with significant gaps like Kerala and Haryana.
\end{itemize}

	
	\begin{figure}[h]
		\centering
		\includegraphics[width=\textwidth]{D:/internship/Competition/Figure_31.png} % Adjust path & format if needed
		\caption{Urban-to-Rural MPCE Ratio Analysis}
		\label{fig:\detokenize{\thefigure}}
	\end{figure}
	


\section*{ Predicted MPCE for 2024-25 (Rural vs Urban)}

\begin{itemize}
	\item \textbf{Higher Urban MPCE Nationwide}: The orange line, representing Urban MPCE, is consistently higher than the blue line for Rural MPCE across all states, indicating higher spending trends in urban areas due to factors such as cost of living and economic opportunities.
	\item \textbf{Kerala Leads in Urban MPCE}: Kerala shows the highest urban MPCE value among all states, reflecting higher living standards and expenditure patterns in urban households.
	\item \textbf{Narrower Gap in States like Bihar and Chhattisgarh}: States such as Bihar and Chhattisgarh exhibit relatively smaller gaps between rural and urban MPCE, indicating more balanced spending trends.
	\item \textbf{Consistent Growth Across States}: Both Rural and Urban MPCE display an increasing trend across states with higher expenditure levels observed in economically advanced states like Maharashtra, Gujarat, and Tamil Nadu.
	\item \textbf{All-India Trends}: The "All-India" data point reflects an average MPCE value where urban MPCE dominates rural MPCE, highlighting the nationwide urban-rural expenditure disparity.
	\item \textbf{Policy Implications}: These trends emphasize the need for targeted rural development policies to reduce disparities and promote balanced economic growth across states.
\end{itemize}

\begin{figure}[h]
	\centering
	\includegraphics[width=\textwidth]{D:/internship/Competition/Figure_32.png} % Adjust the file path if needed
	\caption{Predicted MPCE for 2024-25 (Rural vs Urban)}
	\label{fig:Predicted_MPCE}
\end{figure}




\newpage


\section*{Urban-to-Rural MPCE Ratio Trend Analysis}

\begin{itemize}
	\item \textbf{Peak in Nominal MPCE Ratio (2009-10)}: This marked the period of the highest urban-to-rural expenditure difference, illustrating a significant disparity in nominal spending.
	\item \textbf{Real MPCE Ratio Stability and Decline}: Initially steady, the Real MPCE Ratio displayed balanced inflation-adjusted spending patterns. The subsequent decline suggests narrowing economic disparities over time.
	

\begin{figure}[h]
	\centering
	\includegraphics[width=0.9\textwidth]{D:/internship/Competition/Figure_34.png} % Ensure the file path is accurate
	\caption{Urban-to-Rural MPCE Ratio Trends Over Time}
	\label{fig:Urban_to_Rural_MPCE}
\end{figure}

\item \textbf{Divergence Between Nominal and Real Ratios}: While the Nominal MPCE Ratio highlights temporal spending disparities, the Real MPCE Ratio points to longer-term economic convergence trends.
\item \textbf{Implications for Policy Development}: These trends indicate potential benefits from targeted rural development policies, further emphasizing the necessity for balanced economic growth across urban and rural sectors.
\end{itemize}





\section*{Projected Future MPCE Analysis Using Exponential Trend}

\begin{itemize}
	\item \textbf{Exponential Growth in MPCE}: Both rural and urban MPCE show exponential growth trends. By 2030, the projected MPCE for urban areas is expected to exceed Rs 12,000, while for rural areas, it is projected to cross Rs 8,000.
	\item \textbf{Persistent Urban-Rural Gap}: Despite growth in both sectors, urban MPCE consistently remains higher than rural MPCE, indicating a persistent disparity in expenditure levels.
	\item \textbf{Convergence or Stability?}: The proportional increase in rural MPCE appears to outpace urban MPCE slightly, suggesting a potential narrowing of relative disparities in the future.
	\item \textbf{Policy Implications}: These projections underscore the need for sustained policy focus on rural development to bridge expenditure gaps while supporting sustainable urban growth.
\end{itemize}

\begin{figure}[h]
	\centering
	\includegraphics[width=\textwidth]{D:/internship/Competition/Figure_35.png} % Update the file path with the correct image location.
	\caption{Projected Future MPCE Using Exponential Trend}
	\label{fig:Projected_MPCE}
\end{figure}





\newpage
\section{Rural vs Urban MPCE Spending Patterns}

\begin{itemize}
	\item \textbf{Significant Gap in Food Total}: Urban households exhibit much higher MPCE for the "Food Total" category, indicating diverse and extensive food consumption compared to rural areas.
	\item \textbf{Higher Urban Spending Across Most Categories}: Urban MPCE surpasses rural MPCE significantly in:
	\begin{itemize}
		\item Milk and Milk Products
		\item Vegetables
		\item Egg, Fish, and Meat
		\item Beverages and Processed Foods
	\end{itemize}
	\item \textbf{Narrower Gap in Basic Necessities}: Spending differences in categories like Cereals and Cereal Substitutes and Pulses and Their Products are smaller, reflecting similarities in essential food consumption.
	\item \textbf{Similar Spending in Pan, Tobacco, and Intoxicants}: Both rural and urban households display comparable expenditure in this category, suggesting similar consumption habits.
	\item \textbf{Higher Fuel and Light Costs for Urban Areas}: Urban households spend more on Fuel and Light, likely due to higher energy demands and costs associated with urban living.
\end{itemize}

\begin{figure}[h]
	\centering
	\includegraphics[width=1.2\textwidth]{D:/internship/Competition/Figure_36.png} % Replace with the correct file path
	\caption{Rural vs Urban MPCE Spending Patterns}
	\label{fig:MPCE_Spending}
\end{figure}


\newpage

\section{Rural vs Urban Expenditure Share Analysis}

\begin{itemize}
	\item \textbf{Rural Areas}: 
	\begin{itemize}
		\item Food Share: 57.0\% of total expenditure.
		\item Non-Food Share: 43.0\% of total expenditure.
	\end{itemize}
	\item \textbf{Urban Areas}: 
	\begin{itemize}
		\item Food Share: 47.6\% of total expenditure.
		\item Non-Food Share: 52.4\% of total expenditure.
	\end{itemize}
	\item \textbf{Urban-Rural Comparison}: 
	\begin{itemize}
		\item Rural households prioritize food, allocating a higher percentage of their expenditure to essentials.
		\item Urban households have a more balanced spending pattern, with a larger proportion dedicated to non-food items like housing and education.
	\end{itemize}
	\item \textbf{Policy Implications}: 
	\begin{itemize}
		\item Policies aimed at increasing rural discretionary income can help shift spending patterns and reduce dependency on food expenditures.
		\item Enhancing access to non-food resources in rural areas might bridge the urban-rural divide.
	\end{itemize}
\end{itemize}

\begin{figure}[h]
	\centering
	\includegraphics[width=\textwidth]{D:/internship/Competition/Figure_37.png} % Ensure the correct file path
	\caption{Rural vs Urban Expenditure Share}
	\label{fig:Rural_Urban_Share}
\end{figure}


\newpage
\section{Healthy vs Processed Food Consumption Analysis}



\begin{figure}[h]
	\centering
	\includegraphics[width=\textwidth]{D:/internship/Competition/Figure_38.png} % Ensure the correct file path
	\caption{Healthy vs Processed Food Consumption in Rural and Urban Areas}
	\label{fig:Food_Consumption}
\end{figure}




\begin{itemize}
	\item \textbf{Rural Areas}: 
	\begin{itemize}
		\item Healthy Food Share: 68.4\% of total food consumption.
		\item Processed Food Share: 31.6\% of total food consumption.
	\end{itemize}
	\item \textbf{Urban Areas}: 
	\begin{itemize}
		\item Healthy Food Share: 61.6\% of total food consumption.
		\item Processed Food Share: 38.4\% of total food consumption.
	\end{itemize}
	\item \textbf{Urban-Rural Comparison}: 
	\begin{itemize}
		\item Rural households consume a higher proportion of healthy food compared to urban households.
		\item Urban areas exhibit greater processed food consumption, potentially reflecting urban lifestyle and accessibility.
	\end{itemize}
	\item \textbf{Policy Implications}: 
	\begin{itemize}
		\item Encouraging urban populations to shift towards healthier options through awareness campaigns and incentives.
		\item Supporting rural access to diverse and balanced food choices while preserving traditional eating habits.
	\end{itemize}
\end{itemize}




\subsection{Comparison of Rural vs Urban MPCE in Non-Food Categories}


\begin{figure}[h]
	\centering
	\includegraphics[width=\textwidth]{D:/internship/Competition/Figure_39.png} % Replace with the actual file path of the chart
	\caption{Comparison of Rural vs Urban MPCE in Non-Food Categories}
	\label{fig:Rural_Urban_NonFood_MPCE}
\end{figure}


\begin{itemize}
	\item \textbf{General Observation}: Urban MPCE is consistently higher than Rural MPCE across all non-food categories.
	\item \textbf{Significant Gap in Non-Food Total}: Urban areas show the highest MPCE in the Non-Food Total category, highlighting broader spending capabilities in non-food sectors.
	\item \textbf{Key Categories with Notable Differences}:
	\begin{itemize}
		\item \textbf{Education}: Urban MPCE is significantly higher, indicating greater expenditure on educational services.
		\item \textbf{Medical}: Urban households spend noticeably more, reflecting either higher costs or more frequent access to healthcare services.
		\item \textbf{Conveyance}: Urban spending outpaces rural, aligning with increased reliance on transport due to urban living conditions.
		\item \textbf{Rent, Taxes, and Cesses}: Urban households allocate more funds in this category, consistent with higher costs of living in urban areas.
	\end{itemize}
	\item \textbf{Smaller Disparities in Other Categories}: 
	\begin{itemize}
		\item Categories such as Clothing, Bedding and Footwear and Durable Goods show smaller MPCE differences, indicating relatively closer spending patterns between rural and urban populations.
	\end{itemize}
	\item \textbf{Policy Implications}: 
	\begin{itemize}
		\item The data underscores the need for tailored policies to address disparities in education and healthcare accessibility in rural areas.
		\item Enhanced rural spending capacity and infrastructure development could help bridge urban-rural gaps in key non-food sectors.
	\end{itemize}
\end{itemize}







\subsection{Percentage Share of MPCE in Rural Areas}

\begin{figure}[h]
	\centering
	\includegraphics[width=1.2\textwidth]{D:/internship/Competition/Figure_40.png} % Replace with the actual file path.
	\caption{Percentage Share of MPCE in Rural Areas}
	\label{fig:Rural_MPCE_Share}
\end{figure}




\begin{itemize}
	\item \textbf{Major Share in Non-Food Total}: 55.2\% of rural expenditure is directed toward non-food items, reflecting significant allocation across various categories.
	\item \textbf{Durable Goods}: Constitutes 6.7\% of total expenditure, representing investments in essential long-lasting items.
	\item \textbf{Clothing, Bedding, and Footwear}: Accounts for 6.9$\%$, indicating spending on textiles and accessories.
	\item \textbf{Miscellaneous Goods and Entertainment}: Captures 6.5\% of total expenditure, showcasing some discretionary spending.
	\item \textbf{Essentials}: 
	\begin{itemize}
		\item Conveyance: 7.9\%
		\item Medical: 7.1\%
		\item Consumer Services: 5.5\%
		\item Education: 3.4\%
	\end{itemize}
	\item \textbf{Taxes and Cesses}: Minimal expenditure at 0.6\%, reflecting limited tax obligations in rural households.
\end{itemize}








\subsection{ Urban MPCE Spending Patterns}



\begin{itemize}
	\item \textbf{Major Share in Non-Food Total}: Non-food total accounts for 53.3\% of urban expenditure, highlighting significant spending beyond essential food items.
	\item \textbf{Key Expenditure Categories}:
	\begin{itemize}
		\item Conveyance: 7.5\%
		\item Education: 5.3\%
		\item Medical: 5.2\%
		\item Rent: 5.5\%
		\item Durable Goods: 5.8\%
		\item Taxes and Cesses: 5.3\%
	\end{itemize}
	\item \textbf{Miscellaneous Spending Patterns}:
	\begin{itemize}
		\item Consumer services (5.1$\%$) and miscellaneous goods and entertainment (6.1$\%$) reflect discretionary expenditures within urban households.
		\item Clothing, Bedding and Footwear: 6.1$\%$, indicating moderate spending in these areas.
	\end{itemize}
	\item \textbf{Policy Implications}: 
	\begin{itemize}
		\item These spending patterns underline the economic diversity of urban areas, with investments across essential and discretionary sectors.
		\item Policymakers may aim to optimize urban resource allocations by improving access to education, healthcare, and conveyance.
	\end{itemize}
\end{itemize}



\begin{figure}[h]
	\centering
	\includegraphics[width=\textwidth]{D:/internship/Competition/Figure_41.png} % Replace with the actual file path of the pie chart
	\caption{Percentage Share of MPCE in Urban Areas}
	\label{fig:Urban_MPCE_Share}
\end{figure}







\subsection{ Rural vs Urban MPCE in Non-Food Categories}



\begin{figure}[h]
	\centering
	\includegraphics[width=\textwidth]{D:/internship/Competition/Figure_42.png} % Replace with the actual file path of the heatmap
	\caption{Rural and Urban MPCE Comparison Across Non-Food Categories}
	\label{fig:Rural_Urban_Heatmap}
\end{figure}


\begin{itemize}
	\item \textbf{Education}: Urban MPCE (\textbf{418}) is significantly higher than Rural MPCE (\textbf{133}), reflecting greater investment in education in urban households.
	\item \textbf{Medical Expenditure}: Urban households spend \textbf{409}, almost 50\% more compared to rural households (\textbf{282}), highlighting higher access or cost of healthcare.
	\item \textbf{Conveyance}: Urban MPCE is nearly double (\textbf{592}) of rural areas (\textbf{313}), likely due to increased transportation needs in urban environments.
	\item \textbf{Non-Food Total}: The most significant disparity, with Urban MPCE at \textbf{4220} compared to Rural MPCE of \textbf{2183}, indicating broader urban expenditure capacities.
	\item \textbf{Smaller Disparities}:
	\begin{itemize}
		\item Categories such as \textbf{Clothing, Bedding, and Footwear}, and \textbf{Durable Goods} show relatively closer spending patterns, suggesting common necessities across both sectors.
		\item \textbf{Taxes and Cesses}: Minimal contributions, but urban households (\textbf{38}) still outpace rural (\textbf{10}) by a wide margin.
	\end{itemize}
	\item \textbf{Policy Implications}:
	\begin{itemize}
		\item Enhanced rural investment in education, healthcare, and transport infrastructure could help bridge these stark divides.
		\item Targeted subsidies in essential sectors (e.g., durable goods, clothing) may alleviate economic burdens in rural areas.
	\end{itemize}
\end{itemize}





\section{ Urbanization Impact on Food Share}

\begin{itemize}
	\item \textbf{Higher Food Share in Rural Areas}: 
	\begin{itemize}
		\item Food expenditure in rural areas consistently accounts for a larger percentage compared to urban counterparts across all states.
		\item Represents a focus on subsistence-based consumption and dependency on agricultural produce.
	\end{itemize}
	\item \textbf{Regional Trends}: 
	\begin{itemize}
		\item States like Kerala and Punjab exhibit narrower disparities, possibly due to economic and cultural factors.
		\item Larger gaps are observed in states like Bihar and Madhya Pradesh, where rural food share is significantly higher, indicating economic constraints.
	\end{itemize}
	\item \textbf{Economic Implications}: 
	\begin{itemize}
		\item Lower food share in urban areas reflects diversified expenditure patterns, allocating more resources to non-food items such as education, housing, and healthcare.
		\item Rural dependency on food expenditure underscores the need for infrastructure and development initiatives.
	\end{itemize}
	\item \textbf{Policy Implications}: 
	\begin{itemize}
		\item Targeted policies can help bridge the economic disparity by diversifying spending in rural areas.
		\item Development in healthcare, education, and transport infrastructure could facilitate balanced urban-rural growth.
	\end{itemize}
\end{itemize}

\begin{figure}[h]
	\centering
	\includegraphics[width=\textwidth]{D:/internship/Competition/Figure_45.png} % Replace the file path with the actual location of your graph
	\caption{Urbanization Impact on Food Share (\%) Across Indian States}
	\label{fig:Food_Share_Impact}
\end{figure}






\newpage
\section*{ Rural Food Expenditure by State in India}

\begin{itemize}
	\item \textbf{Top-Spending State (Assam)}:
	\begin{itemize}
		\item Assam ranks highest in rural food expenditure, potentially due to regional dietary preferences and reliance on agricultural staples.
	\end{itemize}
	\item \textbf{Lowest-Spending State (Kerala)}:
	\begin{itemize}
		\item Kerala records the lowest expenditure, which may indicate more diversified spending patterns or differences in food prices and access.
	\end{itemize}
	\item \textbf{Middle-Tier States}:
	\begin{itemize}
		\item States like Jharkhand, Odisha, and Haryana reflect moderate food expenditures, showcasing varying levels of economic activity and consumption behaviors.
		
	\end{itemize}
	\item \textbf{Regional Trends}:
	\begin{itemize}
		\item Northern states such as Haryana and Rajasthan exhibit relatively lower food expenditures compared to northeastern states like Assam and Bihar, highlighting economic and cultural disparities.
	\end{itemize}
	\item \textbf{Policy Implications}:
	\begin{itemize}
		\item Addressing economic disparity through targeted food subsidy programs in high-spending states.
		\item Promoting infrastructure and economic development in lower-spending states to improve overall expenditure and quality of life.
	\end{itemize}
\end{itemize}

\begin{figure}[h]
	\centering
	\includegraphics[width=\textwidth]{D:/internship/Competition/Figure_46.png} % Replace with the actual file path of the bar chart image
	\caption{Highest Food Expenditure States (Rural)}
	\label{fig:Rural_Food_Expenditure}
\end{figure}


\newpage


\subsection{ Cereal \% vs. Food \% (Rural vs Urban) 
	}

\begin{itemize}
	\item \textbf{Plot Overview}: The scatter plot compares the percentage share of cereal consumption (x-axis) against the percentage share of food consumption (y-axis) in rural and urban areas.
	\item \textbf{Rural vs Urban Distribution}:
	\begin{itemize}
		\item Rural data points (blue dots) are more concentrated toward higher cereal and food share percentages, showcasing a higher dependency on cereals.
		\item Urban data points (red dots) are distributed toward lower cereal and food share percentages, reflecting a diversified diet with lesser emphasis on cereals.
	\end{itemize}
	\item \textbf{Trends and Observations}:
	\begin{itemize}
		\item There is a visible correlation where higher cereal share corresponds to higher food share, especially prominent in rural areas.
		\item Urban areas show less reliance on cereals, with data points skewed toward the lower end of the x-axis.
	\end{itemize}
	\item \textbf{Economic Implications}:
	\begin{itemize}
		\item Higher cereal share in rural areas could indicate limited access to diverse food items or a focus on subsistence farming.
		\item Urban dietary diversity highlights the impact of higher incomes and better access to varied food sources.
	\end{itemize}
	\item \textbf{Policy Insights}:
	\begin{itemize}
		\item Policies targeting food security in rural areas could promote access to diverse and nutritious food.
		\item Urban health policies should address the balance between dietary diversity and nutritional adequacy.
	\end{itemize}
\end{itemize}

\begin{figure}[h]
	\centering
	\includegraphics[width=\textwidth]{D:/internship/Competition/Figure_44.png} % Replace with the actual file path of the scatter plot
	\caption{Cereal \% vs. Food \% (Rural vs Urban)}
	\label{fig:Cereal_Food_ScatterPlot}
\end{figure}


\begin{figure}[h]
	\centering
	\includegraphics[width=\textwidth]{D:/internship/Competition/Figure_47.png} % Replace with the actual file path of the scatter plot
	\caption{Clustering States Based on Food and Cereal Spending}
	\label{fig:Food_Cereal_Clustering}
\end{figure}
\newpage
\subsection{Clustering States Based on Food and Cereal Spending}

\begin{itemize}
	\item \textbf{Visualization Summary}:
	\begin{itemize}
		\item The scatter plot maps states based on their Rural Cereal Share (\%) against Rural Food Share (\%).
		\item States are grouped into three clusters represented by different colors: Cluster 0 (red), Cluster 1 (blue), and Cluster 2 (green).
	\end{itemize}
	\item \textbf{Cluster Analysis}:
	\begin{itemize}
		\item \textbf{Cluster 0 (Red)}:
		\begin{itemize}
			\item Includes states with both low Rural Cereal Share and Rural Food Share.
			\item May represent states where diversified food consumption patterns are observed in rural areas.
		\end{itemize}
		\item \textbf{Cluster 1 (Blue)}:
		\begin{itemize}
			\item Contains states with higher Rural Cereal Share and moderate Rural Food Share percentages.
			\item Reflects a dependency on cereals as a major component of food expenditure.
		\end{itemize}
		\item \textbf{Cluster 2 (Green)}:
		\begin{itemize}
			\item Represents states with the highest Rural Food Share and Rural Cereal Share.
			\item Indicates reliance on agricultural staples and limited dietary diversity.
		\end{itemize}
	\end{itemize}
	\item \textbf{Policy Implications}:
	\begin{itemize}
		\item Cluster 0 states may benefit from policies promoting cereal consumption or access to staple foods.
		\item Cluster 1 states might require initiatives that encourage dietary diversification in rural areas.
		\item Cluster 2 states could benefit from programs enhancing economic growth and access to diverse food items.
	\end{itemize}
\end{itemize}









\begin{figure}[h]
	\centering
	\includegraphics[width=\textwidth]{D:/internship/Competition/Figure_48.png} % Ensure this path leads to your graph
	\caption{Trend in Per Capita Cereal Consumption (1999–2024)}
	\label{fig:Cereal_Consumption_Trend}
\end{figure}

\subsection{Analysis of Per Capita Cereal Consumption Trend (1999–2024)}

\begin{itemize}
	\item \textbf{Decline in Consumption for Both Sectors}: 
	\begin{itemize}
		\item Rural areas show a steeper decline compared to urban areas, indicating changing consumption patterns or reduced reliance on cereals.
		\item Urban areas also witness a gradual drop, though the extent is less pronounced than in rural regions.
	\end{itemize}
	\item \textbf{Comparison Between Rural and Urban Trends}:
	\begin{itemize}
		\item Rural areas start with higher cereal consumption in 1999–00 but converge closer to urban levels by 2024.
		\item The disparity between rural and urban cereal consumption narrows over time, reflecting shifts in dietary habits or economic factors.
	\end{itemize}
	\item \textbf{Year-Wise Observations}:
	\begin{itemize}
		\item 1999–2000: Cereal consumption is at its peak, with rural households consuming significantly more than urban counterparts.
		\item 2011–2012: A steady decline is observed across both rural and urban sectors.
		\item 2022–2024: The lowest recorded cereal consumption is projected, highlighting ongoing changes in food preferences and economic scenarios.
	\end{itemize}
	\item \textbf{Possible Contributing Factors}:
	\begin{itemize}
		\item Dietary Diversification: Increased consumption of other food groups like proteins and processed foods.
		\item Economic Growth: Enhanced income levels leading to diverse eating habits in urban areas.
		\item Lifestyle Changes: Rising preference for ready-to-eat foods, especially in urban households.
	\end{itemize}
	\item \textbf{Policy Implications}:
	\begin{itemize}
		\item Rural areas may require interventions to ensure nutritional adequacy amid declining cereal consumption.
		\item Strategies to promote balanced diets and support for agricultural production might address these trends effectively.
	\end{itemize}
\end{itemize}








\subsection{Analysis of Total Cereal Consumption by State (kg per month)}

\begin{itemize}
	\item \textbf{Highest Consumption (Maharashtra)}:
	\begin{itemize}
		\item Maharashtra records the highest cereal consumption per month among Indian states, reflecting substantial dietary reliance on cereals.
	\end{itemize}
	\item \textbf{Lowest Consumption (Kerala)}:
	\begin{itemize}
		\item Kerala has the lowest cereal consumption, potentially due to dietary preferences or availability of alternative staple foods.
	\end{itemize}
	\item \textbf{Moderate Consumption States}:
	\begin{itemize}
		\item States like Andhra Pradesh, Tamil Nadu, and Gujarat show moderate consumption levels, suggesting diverse diets with a balanced reliance on cereals and other food groups.
	\end{itemize}
	\item \textbf{Regional Patterns}:
	\begin{itemize}
		\item Northeastern states such as Assam exhibit relatively high cereal consumption, indicative of staple food habits in these regions.
		\item Southern states, particularly Kerala, appear to have more diversified consumption patterns with lower dependency on cereals.
	\end{itemize}
	\item \textbf{Policy Implications}:
	\begin{itemize}
		\item High cereal consumption in states like Maharashtra and Assam may require support for sustainable agricultural production and storage.
		\item States with lower cereal consumption, such as Kerala, might benefit from nutritional awareness programs to ensure balanced diets.
	\end{itemize}
\end{itemize}

\begin{figure}[h]
	\centering
	\includegraphics[width=\textwidth]{D:/internship/Competition/Figure_49.png} % Replace the file path with the correct location of your image
	\caption{Total Cereal Consumption by State (kg per month)}
	\label{fig:Cereal_Consumption_By_State}
\end{figure}





\begin{figure}[h]
	\centering
	\includegraphics[width=\textwidth]{D:/internship/Competition/Figure_50.png} % Replace this path with the actual file location of your heatmap
	\caption{Correlation Heatmap for Cereal Consumption and Monthly Per Capita Cereal Consumption (kg)}
	\label{fig:Cereal_Correlation_Heatmap}
\end{figure}

\subsection{Correlation of Cereal Consumption}

\begin{itemize}
	\item \textbf{Correlation Highlights}:
	\begin{itemize}
		\item \textbf{Monthly Per Capita Cereal Consumption (kg)}:
		\begin{itemize}
			\item Positively correlated with \% Share of Rice (0.23), indicating rice consumption contributes slightly to higher overall cereal consumption.
			\item Negatively correlated with \% Share of Wheat (-0.19), \% Share of Coarse Grains (-0.27), and \% Share of Other Cereals (-0.35), suggesting these cereals are less commonly associated with higher total cereal consumption.
		\end{itemize}
		\item \textbf{\% Share of Rice}:
		\begin{itemize}
			\item Strongly negatively correlated with \% Share of Wheat (-0.98), implying that regions with high rice consumption tend to have low wheat consumption.
			\item Moderately negatively correlated with \% Share of Coarse Grains (-0.40) and Other Cereals (-0.46).
		\end{itemize}
		\item \textbf{\% Share of Coarse Grains and Other Cereals}:
		\begin{itemize}
			\item Strong positive correlation (0.94) between these two categories, highlighting co-consumption patterns.
		\end{itemize}
	\end{itemize}
	\item \textbf{Policy Implications}:
	\begin{itemize}
		\item Regions with high reliance on rice or wheat may need diversification programs to include coarse grains and other cereals for balanced nutrition.
		\item Addressing the negative correlations between rice/wheat consumption and coarse grains could promote greater cereal diversity.
	\end{itemize}
\end{itemize}

\begin{figure}[h]
	\centering
	\includegraphics[width=\textwidth]{D:/internship/Competition/Figure_51.png} % Replace this path with the correct file location of your scatter plot image
	\caption{Clustering of States Based on Cereal Consumption Patterns}
	\label{fig:Cereal_Consumption_Clusters}
\end{figure}

\subsection{ Cereal Consumption Clustering}

\begin{itemize}
	\item \textbf{Scatter Plot Overview}:
	\begin{itemize}
		\item The scatter plot compares Rice Consumption (\%) on the x-axis with Wheat Consumption (\%) on the y-axis.
		\item States are clustered into three distinct groups based on their cereal consumption patterns:
		\begin{itemize}
			\item \textbf{Cluster 0 (Red)}: Low Rice Consumption, High Wheat Consumption.
			\item \textbf{Cluster 1 (Blue)}: Balanced Rice and Wheat Consumption.
			\item \textbf{Cluster 2 (Green)}: High Rice Consumption, Low Wheat Consumption.
		\end{itemize}
	\end{itemize}
	\item \textbf{Key Observations}:
	\begin{itemize}
		\item States in Cluster 2 show a strong preference for rice, typical in regions where rice is a staple, e.g., southeastern states.
		\item Cluster 0 highlights states reliant on wheat, commonly seen in northern regions.
		\item Cluster 1 demonstrates balanced cereal consumption, indicating dietary diversity in these states.
	\end{itemize}
	\item \textbf{Policy Implications}:
	\begin{itemize}
		\item Promoting dietary diversity through awareness programs in regions overly reliant on one type of cereal.
		\item Ensuring access to a variety of cereals, particularly in states that display skewed preferences.
		\item Aligning agricultural strategies with regional dietary patterns to optimize resource use.
	\end{itemize}
\end{itemize}






\section{Analysis of Clustering of States Based on Food Expenditure(Rural)}

\begin{itemize}
	\item \textbf{Plot Overview}: 
	\begin{itemize}
		\item The scatter plot visualizes the relationship between Cereal Expenditure (\%) and Processed Food Expenditure (\%) across states in India.
		\item States are categorized into four clusters, each represented by a different color: 
		\begin{itemize}
			\item \textbf{Cluster 0 (Red)}: Odisha, Jharkhand.
			\item \textbf{Cluster 1 (Teal)}: Gujarat, Chhattisgarh, Uttar Pradesh, Madhya Pradesh, Maharashtra, Telangana, West Bengal.
			\item \textbf{Cluster 2 (Green)}: Kerala, Andhra Pradesh, Assam, Karnataka, Tamil Nadu.
			\item \textbf{Cluster 3 (Purple)}: Punjab, Haryana, Rajasthan.
		\end{itemize}
	\end{itemize}
	\item \textbf{Cluster Analysis}:
	\begin{itemize}
		\item \textbf{Cluster 0 (Red)}:
		\begin{itemize}
			\item States in this cluster have relatively high Cereal Expenditure but low Processed Food Expenditure.
			\item Likely represents rural or agrarian economies with traditional dietary patterns.
		\end{itemize}
		\item \textbf{Cluster 1 (Teal)}:
		\begin{itemize}
			\item Includes states with moderate values for both Cereal and Processed Food Expenditure.
			\item Reflects transitional economies moving toward diversified food habits.
		\end{itemize}
		\item \textbf{Cluster 2 (Green)}:
		\begin{itemize}
			\item Features states with low Cereal Expenditure and high Processed Food Expenditure.
			\item Indicates urbanized or economically advanced regions with modernized diets.
		\end{itemize}
		\item \textbf{Cluster 3 (Purple)}:
		\begin{itemize}
			\item States exhibit a balance between Cereal and Processed Food Expenditure.
			\item Suggests regions with varied diets that include both traditional and processed food options.
		\end{itemize}
	\end{itemize}
	\item \textbf{Policy Implications}:
	\begin{itemize}
		\item Cluster 0 states may benefit from initiatives encouraging dietary diversification and improved access to processed or fortified foods.
		\item Cluster 1 states could emphasize infrastructure for food distribution to support economic and dietary transitions.
		\item Cluster 2 states might focus on promoting balanced consumption of traditional staples alongside processed foods for nutritional adequacy.
		\item Cluster 3 regions could leverage their diversity in food habits as a model for balanced regional diets.
	\end{itemize}
\end{itemize}
\begin{figure}[h]
	\centering
	\includegraphics[width=\textwidth]{D:/internship/Competition/Figure_53.png} % Replace with your scatter plot image file location
	\caption{Clustering of States Based on Food Expenditure}
	\label{fig:Food_Expenditure_Clusters}
\end{figure}




\subsection{ Clustering of States Based on Urban Food Expenditure}

\begin{itemize}
	\item \textbf{Scatter Plot Overview}: 
	\begin{itemize}
		\item This scatter plot represents the relationship between Cereal Expenditure (\%) and Processed Food Expenditure (\%) in urban areas across Indian states.
		\item States are grouped into four distinct clusters, visually represented by different colors:
		\begin{itemize}
			\item \textbf{Cluster 0 (Red)}: States such as Tamil Nadu, Telangana, Assam, Karnataka, and Maharashtra.
			\item \textbf{Cluster 1 (Blue)}: Includes states like Kerala, Andhra Pradesh, Gujarat, and Chhattisgarh.
			\item \textbf{Cluster 2 (Green)}: States such as Haryana, Punjab, and Rajasthan.
			\item \textbf{Cluster 3 (Purple)}: Odisha, West Bengal, Bihar, and Jharkhand.
		\end{itemize}
	\end{itemize}
	\item \textbf{Cluster Characteristics}:
	\begin{itemize}
		\item \textbf{Cluster 0 (Red)}:
		\begin{itemize}
			\item Shows higher reliance on processed food expenditure with comparatively lower cereal spending.
			\item Urbanization and lifestyle shifts likely influence these patterns.
		\end{itemize}
		\item \textbf{Cluster 1 (Blue)}:
		\begin{itemize}
			\item Represents a balance between cereal and processed food expenditure.
			\item May include transitioning urban economies with evolving dietary patterns.
		\end{itemize}
		\item \textbf{Cluster 2 (Green)}:
		\begin{itemize}
			\item Exhibits moderate cereal expenditure with diverse urban food habits including processed options.
			\item Indicates relative economic affluence and access to modern food markets.
		\end{itemize}
		\item \textbf{Cluster 3 (Purple)}:
		\begin{itemize}
			\item Skewed heavily towards cereal expenditure with lower processed food costs.
			\item Likely reflects the influence of traditional eating habits in less urbanized states.
		\end{itemize}
	\end{itemize}
	\item \textbf{Policy Implications}:
	\begin{itemize}
		\item For states in Cluster 3, urban food programs could promote dietary diversity and better access to processed food markets.
		\item Cluster 0 and 1 regions might benefit from awareness campaigns on balanced diets to ensure nutritional adequacy alongside lifestyle-driven food consumption.
		\item Cluster 2 urban areas could explore strategies to optimize food production and distribution systems for sustainable urban diets.
	\end{itemize}
\end{itemize}
\begin{figure}[h]
	\centering
	\includegraphics[width=\textwidth]{D:/internship/Competition/Figure_56.png} % Replace this path with your scatter plot image file location
	\caption{Clustering of States Based on Urban Food Expenditure}
	\label{fig:Urban_Food_Expenditure_Clusters}
\end{figure}





\newpage

\subsection{Analysis of State-wise Food Category Distribution}

\begin{itemize}
	\item \textbf{Cereal as a Staple}:
	\begin{itemize}
		\item Cereals dominate expenditure across most states. For example, Bihar, Odisha, and West Bengal allocate a significant percentage to cereals, emphasizing their importance as dietary staples.
	\end{itemize}
	\item \textbf{Milk \& Milk Products}:
	\begin{itemize}
		\item States like Haryana, Punjab, and Gujarat exhibit a larger proportion of expenditure on milk and dairy products, reflecting regional dietary preferences.
	\end{itemize}
	\item \textbf{Fruits \& Vegetables}:
	\begin{itemize}
		\item Kerala and Tamil Nadu show higher expenditure on fruits and vegetables, showcasing dietary diversity and better access to fresh produce.
	\end{itemize}
	\item \textbf{Protein Sources (Egg, Fish \& Meat)}:
	\begin{itemize}
		\item Assam and Odisha highlight notable spending on protein-rich foods, driven by cultural and regional dietary habits.
	\end{itemize}
	\item \textbf{Processed Foods \& Beverages}:
	\begin{itemize}
		\item Urbanized states like Maharashtra and Karnataka reveal increased spending on processed foods and beverages, influenced by lifestyle and convenience.
	\end{itemize}
	\item \textbf{Other Expenditure}:
	\begin{itemize}
		\item States such as Telangana and West Bengal include spending in miscellaneous categories, indicating a variety of food preferences unique to these regions.
	\end{itemize}
\end{itemize}

\section*{Policy Implications}

\begin{itemize}
	\item \textbf{Cereal-centric Diets}: States heavily dependent on cereals could benefit from programs promoting nutritional diversity.
	\item \textbf{Regional Specialization}: Dairy-rich states like Punjab might focus on enhancing dairy production, while fish-consuming states like Odisha could support aquaculture initiatives.
	\item \textbf{Urban vs. Rural Patterns}: Rising processed food consumption in urban areas suggests a need for health awareness campaigns encouraging balanced diets.
\end{itemize}

\begin{figure}[h]
	\centering
	\includegraphics[width=\textwidth]{D:/internship/Competition/Figure_52.png} % Replace this path with your image's actual file location
	\caption{State-wise Food Category Distribution in India}
	\label{fig:Food_Category_Distribution}
\end{figure}

\newpage
\section{Average MPCE by Social Group in Rural Areas}





		
		\begin{longtable}{|p{4cm}|p{3cm}|p{3cm}|p{3cm}|p{2cm}|}
			\hline
			\textbf{State} & \textbf{Scheduled Tribe (ST)} & \textbf{Scheduled Caste (SC)} & \textbf{Other Backward Class (OBC)} & \textbf{Others} \\ \hline
			\endhead
			
			Jammu  Kashmir    & 3948  & 4806  & 3866  & 5092  \\ \hline
			Himachal Pradesh   & 5406  & 5250  & 5815  & 6190  \\ \hline
			Punjab             & 5150  & 5136  & 5921  & 6867  \\ \hline
			Chandigarh         & 20299 & 7643  & 8674  & 9473  \\ \hline
			Uttarakhand        & 4687  & 4213  & 4664  & 5617  \\ \hline
			Haryana            & 5688  & 4756  & 5529  & 5786  \\ \hline
			Delhi              & 6974  & 6862  & 6552  & 8901  \\ \hline
			Rajasthan          & 3384  & 4195  & 4854  & 5238  \\ \hline
			Uttar Pradesh      & 2980  & 3219  & 3480  & 3955  \\ \hline
			Bihar              & 3334  & 3374  & 3733  & 4017  \\ \hline
			Sikkim             & 9318  & 8079  & 9564  & 10170 \\ \hline
			Arunachal Pradesh  & 6048  & 5634  & 7815  & 5568  \\ \hline
			Nagaland           & 5151  & 5487  & 5064  & ---   \\ \hline
			Manipur            & 3997  & 4109  & 4688  & 5102  \\ \hline
			Mizoram            & 5948  & 9733  & ---   & 5901  \\ \hline
			Tripura            & 5803  & 6414  & 6637  & 6678  \\ \hline
			Meghalaya          & 3820  & 3043  & 4227  & 5174  \\ \hline
			Assam              & 3774  & 3758  & 3823  & 3785  \\ \hline
			West Bengal        & 3077  & 3507  & 3630  & 3791  \\ \hline
			Jharkhand          & 2497  & 2827  & 3260  & 3361  \\ \hline
			Odisha             & 2800  & 3235  & 3618  & 3924  \\ \hline
			Chhattisgarh       & 2471  & 3047  & 2902  & 3329  \\ \hline
			Madhya Pradesh     & 3004  & 3416  & 3582  & 3964  \\ \hline
			Gujarat            & 3690  & 3825  & 4008  & 5137  \\ \hline
			D  N.Haveli Daman Diu & 3943  & 5052  & 5147  & 6902  \\ \hline
			Maharashtra        & 3103  & 3986  & 4242  & 4599  \\ \hline
			Andhra Pradesh     & 4162  & 5145  & 5384  & 5715  \\ \hline
			Karnataka          & 4590  & 4620  & 4926  & 5566  \\ \hline
			Goa                & 6006  & 7623  & 7940  & 8638  \\ \hline
			Lakshadweep        & 6263  & 5394  & 16056 & 16624 \\ \hline
			Kerala             & 4908  & 5449  & 6262  & 7738  \\ \hline
			Tamil Nadu         & 4831  & 5340  & 5835  & 6202  \\ \hline
			Puducherry         & ---   & 6749  & 8013  & 16785 \\ \hline
			Andaman Nicobar Islands & 5688  & 4133  & 7649  & 8307  \\ \hline
			Telangana          & 4981  & 5162  & 5482  & 6449  \\ \hline
			Ladakh             & 5010  & ---   & ---   & ---   \\ \hline
			
			\caption{Average MPCE by Social Group in 2023-24 Across Indian States}
			\label{table:mpce_social_groups}
		\end{longtable}

\begin{figure}[h]
	\centering
	\includegraphics[width=\textwidth]{D:/internship/Competition/Figure_59.png} % Replace this path with the actual file location of the bar graph
	\caption{Average MPCE by Social Group in Rural Areas}
	\label{fig:Average_MPCE_Social_Group}
\end{figure}



\begin{itemize}
	\item \textbf{Social Group Comparison}: 
	\begin{itemize}
		\item \textbf{Others (Purple Bar)}: Represents the highest average MPCE, indicating relatively better economic conditions among this group in rural areas.
		\item \textbf{OBC (Green Bar)}: Second highest MPCE, showcasing moderate spending ability, potentially influenced by access to resources and income levels.
		\item \textbf{SC (Red Bar)}: Third position in MPCE, reflecting limited economic expenditure compared to OBC and Others.
		\item \textbf{ST (Blue Bar)}: The lowest average MPCE, highlighting economic constraints faced by this group in rural areas.
	\end{itemize}
	\item \textbf{Observations}:
	\begin{itemize}
		\item The disparity in MPCE across social groups underscores existing socio-economic inequities in rural areas.
		\item Differences in access to income, education, and resources are key contributors to this variation.
	\end{itemize}
	\item \textbf{Policy Implications}:
	\begin{itemize}
		\item \textbf{Focus on ST and SC Groups}: Targeted policies aimed at improving income levels, education, and resource accessibility can help reduce MPCE disparities.
		\item \textbf{Strengthening Rural Economies}: Investment in infrastructure, skills development, and job creation in rural areas can benefit all social groups.
		\item \textbf{Inclusive Growth}: Programs promoting equity and inclusion across social groups can address systemic economic challenges.
	\end{itemize}
\end{itemize}







\newpage
\subsection{Average MPCE by Social Group in Urban Areas}


	
	\begin{longtable}{|p{4cm}|p{3cm}|p{3cm}|p{3cm}|p{2cm}|}
		\hline
		\textbf{State} & \textbf{Scheduled Tribe (ST)} & \textbf{Scheduled Caste (SC)} & \textbf{Other Backward Class (OBC)} & \textbf{Others} \\ \hline
		\endhead
		
		Jammu  Kashmir    & 4872  & 6605  & 6155  & 6372  \\ \hline
		Himachal Pradesh   & 12124 & 7669  & 7835  & 9996  \\ \hline
		Punjab             & 4039  & 6160  & 7054  & 8290  \\ \hline
		Chandigarh         & ---   & 9140  & 11601 & 14740 \\ \hline
		Uttarakhand        & 8513  & 6283  & 6306  & 8329  \\ \hline
		Haryana            & 11188 & 6080  & 7465  & 10113 \\ \hline
		Delhi              & 7116  & 7459  & 7855  & 9475  \\ \hline
		Rajasthan          & 6065  & 5050  & 6051  & 8011  \\ \hline
		Uttar Pradesh      & 5383  & 4918  & 5030  & 6522  \\ \hline
		Bihar              & 4208  & 4174  & 5120  & 5650  \\ \hline
		Sikkim             & 14160 & 11886 & 14298 & 13522 \\ \hline
		Arunachal Pradesh  & 10115 & 10560 & 8879  & 9376  \\ \hline
		Nagaland           & 8161  & 6835  & 6456  & 8638  \\ \hline
		Manipur            & 6406  & 5200  & 5929  & 6113  \\ \hline
		Mizoram            & 8707  & 9944  & 7721  & ---   \\ \hline
		Tripura            & 8714  & 7630  & 7617  & 8461  \\ \hline
		Meghalaya          & 7657  & 7763  & 9612  & 9180  \\ \hline
		Assam              & 6914  & 5918  & 7026  & 6825  \\ \hline
		West Bengal        & 4761  & 4704  & 4976  & 6172  \\ \hline
		Jharkhand          & 4761  & 3969  & 5321  & 6543  \\ \hline
		Odisha             & 4650  & 4611  & 5785  & 6807  \\ \hline
		Chhattisgarh       & 3988  & 4337  & 4596  & 6179  \\ \hline
		Madhya Pradesh     & 4445  & 4711  & 5221  & 6571  \\ \hline
		Gujarat            & 5837  & 5955  & 6415  & 7971  \\ \hline
		D  N. Haveli  Daman  Diu & 5804  & 6517  & 6912  & 7324  \\ \hline
		Maharashtra        & 5377  & 5956  & 6580  & 8250  \\ \hline
		Andhra Pradesh     & 6244  & 6686  & 7102  & 7543  \\ \hline
		Karnataka          & 5908  & 6584  & 7971  & 9614  \\ \hline
		Goa                & 9231  & 6772  & 8465  & 10089 \\ \hline
		Lakshadweep        & 6262  & 7262  & 12688 & 12931 \\ \hline
		Kerala             & 7688  & 6609  & 7165  & 9708  \\ \hline
		Tamil Nadu         & 8050  & 6662  & 8338  & 9139  \\ \hline
		Puducherry         & 14265 & 7140  & 8840  & 9976  \\ \hline
		Andaman  Nicobar Islands & 7218  & 11687 & 12439 & 10011 \\ \hline
		Telangana          & 9065  & 7680  & 8421  & 10893 \\ \hline
		Ladakh             & 7034  & 6157  & ---   & 9752  \\ \hline
		
		\caption{Average MPCE by Social Group Across Indian States (2023-24)}
		\label{table:mpce_social_groups}
	\end{longtable}
	

\begin{figure}[h]
	\centering
	\includegraphics[width=\textwidth]{D:/internship/Competition/Figure_60.png} % Replace with the actual file path of your image
	\caption{Average MPCE by Social Group in Urban Areas}
	\label{fig:Average_MPCE_Social_Group_Urban}
\end{figure}

\begin{itemize}
	\item \textbf{Social Group Comparison}: 
	\begin{itemize}
		\item \textbf{Others (Purple Bar)}: Highest average MPCE, approximately 7500, showcasing better economic conditions and access to resources in urban settings.
		\item \textbf{OBC (Green Bar)}: Slightly above 6000 MPCE, reflecting moderate expenditure levels and economic stability.
		\item \textbf{SC (Red Bar)}: Slightly below 6000 MPCE, highlighting limited economic expenditure compared to OBC and Others.
		\item \textbf{ST (Blue Bar)}: Lowest average MPCE, around 6000, indicating economic disparities faced by Scheduled Tribes in urban areas.
	\end{itemize}
	\newpage
	\item \textbf{Observations}:
	\begin{itemize}
		\item The urban MPCE disparities among social groups reveal ongoing socio-economic inequities despite better access to urban infrastructure.
		\item Factors such as income levels, education, and job opportunities contribute significantly to this variation.
	\end{itemize}
	\item \textbf{Policy Implications}:
	\begin{itemize}
		\item \textbf{Support for ST and SC Groups}: Targeted policies for skill development and resource accessibility to uplift these social groups.
		\item \textbf{Inclusive Urban Growth}: Urban development strategies focused on equitable access to employment, education, and housing.
		\item \textbf{Economic Empowerment Programs}: Initiatives aimed at reducing income disparities across urban social groups.
	\end{itemize}
\end{itemize}


\newpage









\newpage

\section{Analysis of MPCE by Household Type Across Indian States (2023-24)}



	
	\begin{longtable}{|p{2cm}|p{1.9cm}|p{1.95cm}|p{2cm}|p{1.85cm}|p{2cm}|p{2cm}|p{1cm}|}
		\hline
		\textbf{State} & \textbf{Self-Employment in Agriculture} & \textbf{Self-Employment in Non-Agriculture} & \textbf{Regular Wage/Salary in Agriculture} & \textbf{Regular Wage/Salary in Non-Agriculture} & \textbf{Casual Labour in Agriculture} & \textbf{Casual Labour in Non-Agriculture} & \textbf{Others} \\ \hline
		\endhead
		
		Jammu  Kashmir    & 4785  & 5087  & 4502  & 5302  & 3213  & 3635  & 6091  \\ \hline
		Himachal Pradesh   & 5534  & 6156  & 10127 & 6391  & 4430  & 4663  & 7424  \\ \hline
		Punjab             & 6971  & 5772  & 4677  & 5776  & 5020  & 4738  & 7238  \\ \hline
		Chandigarh         & 9940  & 8291  & 6507  & 9234  & 6503  & 7974  & 10519 \\ \hline
		Uttarakhand        & 5124  & 5011  & 5511  & 5407  & 4341  & 4167  & 5352  \\ \hline
		Haryana            & 5829  & 5136  & 5102  & 5397  & 5055  & 4879  & 5841  \\ \hline
		Delhi              & 7467  & 7204  & 8176  & 7545  & 7565  & 6526  & 6823  \\ \hline
		Rajasthan          & 4509  & 4848  & 4645  & 5008  & 3760  & 3908  & 5553  \\ \hline
		Uttar Pradesh      & 3618  & 3513  & 3688  & 3769  & 3161  & 3131  & 3779  \\ \hline
		Bihar              & 3694  & 3954  & 3027  & 3878  & 3327  & 3478  & 4031  \\ \hline
		Sikkim             & 8240  & 9914  & 9378  & 10324 & 5044  & 8191  & 12471 \\ \hline
		Arunachal Pradesh  & 5288  & 6967  & 6085  & 7151  & 4629  & 5351  & 9590  \\ \hline
		Nagaland           & 4796  & 5661  & 5018  & 5919  & 4625  & 4388  & 6419  \\ \hline
		Manipur            & 3995  & 4686  & 4289  & 5265  & 4733  & 4443  & 5557  \\ \hline
		Mizoram            & 5035  & 6482  & 6449  & 6869  & 4344  & 5646  & 5275  \\ \hline
		Tripura            & 5823  & 6463  & 5780  & 7419  & 5690  & 5672  & 8119  \\ \hline
		Meghalaya          & 3684  & 4154  & 3550  & 4608  & 3177  & 3766  & 5482  \\ \hline
		Assam              & 3677  & 3907  & 3236  & 4505  & 3403  & 3587  & 4835  \\ \hline
		West Bengal        & 3522  & 3886  & 3618  & 4235  & 3223  & 3387  & 4156  \\ \hline
		Jharkhand          & 2674  & 3335  & 3058  & 3428  & 2743  & 2830  & 3240  \\ \hline
		Odisha             & 3181  & 3837  & 3844  & 4336  & 2976  & 3008  & 3236  \\ \hline
		Chhattisgarh       & 2622  & 2957  & 3369  & 3435  & 2606  & 2698  & 3066  \\ \hline
		Madhya Pradesh     & 3470  & 3778  & 3456  & 4201  & 3269  & 3072  & 4121  \\ \hline
		Gujarat            & 4207  & 4394  & 4287  & 4589  & 3441  & 3687  & 5725  \\ \hline
		D N Haveli Daman Diu & 3979  & 5353  & 3964  & 4544  & 3111  & 2889  & 5459  \\ \hline
		Maharashtra        & 4099  & 4691  & 4184  & 5261  & 3298  & 3628  & 5013  \\ \hline
		Andhra Pradesh     & 5368  & 5525  & 6339  & 5797  & 4944  & 5097  & 5308  \\ \hline
		Karnataka          & 4670  & 5220  & 4590  & 6899  & 4097  & 4746  & 5822  \\ \hline
		Goa                & 8446  & 8543  & 7962  & 8219  & 5518  & 6596  & 7454  \\ \hline
		Lakshadweep        & 4626  & 6188  & 3596  & 6936  & 10225 & 5724  & 6524  \\ \hline
		Kerala             & 7272  & 6672  & 6072  & 6914  & 5137  & 5712  & 7506  \\ \hline
		Tamil Nadu         & 5521  & 5984  & 6131  & 6381  & 5069  & 5264  & 5860  \\ \hline
		Puducherry         & 7115  & 8241  & 10097 & 8285  & 6716  & 6696  & 8934  \\ \hline
		Andaman Nicobar Islands & 7264  & 7681  & 9080  & 8394  & 4821  & 5461  & 8174  \\ \hline
		Telangana          & 5423  & 5724  & 4638  & 6056  & 4955  & 5082  & 5549  \\ \hline
		Ladakh             & 4687  & 5306  & 5433  & 5429  & 4501  & 6304  & 4492  \\ \hline
		\caption{Average MPCE by Employment Type Across Indian States for Rural Sector (2023-24) Rural}
		\label{table:mpce_employment}
	\end{longtable}
	



	\begin{longtable}{|p{4cm}|p{3cm}|p{3cm}|p{3cm}|p{3cm}|}
		\hline
		\textbf{State} & \textbf{Self-Employment in Agriculture} & \textbf{Self-Employment in Non-Agriculture} & \textbf{Regular Wage/Salary in Agriculture} & \textbf{Others} \\ \hline
		\endhead
		
		Jammu \& Kashmir    & 6142  & 6689  & 4846  & 7459  \\ \hline
		Himachal Pradesh   & 8285  & 9216  & 5836  & 12367 \\ \hline
		Punjab             & 7624  & 7424  & 5618  & 8858  \\ \hline
		Chandigarh         & 12131 & 13703 & 8240  & 23047 \\ \hline
		Uttarakhand        & 7434  & 7764  & 5264  & 7462  \\ \hline
		Haryana            & 9184  & 8748  & 5250  & 10252 \\ \hline
		Delhi              & 8346  & 8573  & 7032  & 10821 \\ \hline
		Rajasthan          & 6204  & 7018  & 4267  & 10543 \\ \hline
		Uttar Pradesh      & 5015  & 6296  & 4080  & 7394  \\ \hline
		Bihar              & 5002  & 5428  & 4198  & 6298  \\ \hline
		Sikkim             & 12693 & 13799 & 12155 & 21921 \\ \hline
		Arunachal Pradesh  & 9035  & 10271 & 8476  & 14299 \\ \hline
		Nagaland           & 7652  & 8314  & 6313  & 9131  \\ \hline
		Manipur            & 5711  & 6388  & 5082  & 7882  \\ \hline
		Mizoram            & 8530  & 8937  & 7918  & 8164  \\ \hline
		Tripura            & 7603  & 8776  & 6195  & 10582 \\ \hline
		Meghalaya          & 7096  & 8548  & 5752  & 9603  \\ \hline
		Assam              & 6267  & 7468  & 4662  & 9433  \\ \hline
		West Bengal        & 5440  & 6348  & 3868  & 7898  \\ \hline
		Jharkhand          & 4883  & 5927  & 3440  & 7658  \\ \hline
		Odisha             & 5379  & 6563  & 3742  & 7193  \\ \hline
		Chhattisgarh       & 4987  & 5204  & 3478  & 5664  \\ \hline
		Madhya Pradesh     & 5380  & 6077  & 3973  & 7667  \\ \hline
		Gujarat            & 7337  & 7296  & 4856  & 8796  \\ \hline
		D N. Haveli Daman Diu  & 6572  & 6899  & 5411  & 10061 \\ \hline
		Maharashtra        & 7117  & 7551  & 4487  & 10170 \\ \hline
		Andhra Pradesh     & 6942  & 7687  & 5752  & 8535  \\ \hline
		Karnataka          & 6905  & 9800  & 5653  & 10108 \\ \hline
		Goa                & 9798  & 10036 & 7334  & 8500  \\ \hline
		Lakshadweep        & 5660  & 6970  & 6565  & 6038  \\ \hline
		Kerala             & 7420  & 8537  & 6357  & 8693  \\ \hline
		Tamil Nadu         & 8181  & 8581  & 6184  & 9693  \\ \hline
		Puducherry         & 9016  & 8815  & 7073  & 8701  \\ \hline
		Andaman  Nicobar Islands & 9687  & 10718 & 9055  & 12630 \\ \hline
		Telangana          & 7965  & 9217  & 6535  & 12619 \\ \hline
		Ladakh             & 7089  & 8356  & 8800  & 6496  \\ \hline
		
		\caption{Average MPCE by Employment Type Across Indian States (2023-24) Urabn}
		\label{table:mpce_employment}
	\end{longtable}
	



\begin{figure}[h]
	\centering
	\includegraphics[width=\textwidth]{D:/internship/Competition/Figure_69.png} % Replace this with the path to your image file
	\caption{Average MPCE by Household Type in Urban Areas (2023-24)}
	\label{fig:Urban_MPCE_Household_Type}
\end{figure}

\begin{itemize}
	\item 
	\begin{itemize}
		\item The chart compares the \textbf{Monthly Per Capita Expenditure (MPCE)} across various household types in urban areas for different Indian states and All-India.
		\item Household types include:
		\begin{itemize}
			\item Self-Employment (blue bar)
			\item Regular Wage/Salary (orange bar)
			\item Casual Labour (green bar)
			\item Others (red bar)
			\item All (purple bar)
		\end{itemize}
	\end{itemize}
	\item \textbf{Key Observations}:
	\begin{itemize}
		\item \textbf{Highest MPCE Category}:
		\begin{itemize}
			\item Regular Wage/Salary households exhibit the highest MPCE across most states, indicating economic advantages tied to steady income sources.
		\end{itemize}
		\item \textbf{Lowest MPCE Category}:
		\begin{itemize}
			\item Casual Labour households have the lowest MPCE, reflecting economic vulnerability and irregular income patterns.
		\end{itemize}
		\item \textbf{All-India Level}:
		\begin{itemize}
			\item Regular Wage/Salary and Self-Employment categories show significantly higher MPCE compared to others.
		\end{itemize}
	\end{itemize}
	\item \textbf{State-Level Observations}:
	\begin{itemize}
		\item \textbf{Kerala}: Exhibits the highest MPCE across all categories, highlighting robust urban economic conditions.
		\item \textbf{Bihar}: Has the lowest MPCE, suggesting economic challenges in urban household expenditures.
		\item \textbf{Gujarat and Maharashtra}: Demonstrate balanced MPCE distributions among household types.
	\end{itemize}
\end{itemize}

\section*{Policy Implications}

\begin{itemize}
	\item \textbf{Support for Casual Labour Households}:
	\begin{itemize}
		\item Targeted policies addressing job security and income stabilization can uplift casual labour households.
	\end{itemize}
	\item \textbf{State-Specific Interventions}:
	\begin{itemize}
		\item States like Bihar need focused development programs to enhance urban income levels.
		\item Kerala's model of balanced urban economic growth could serve as a benchmark for other states.
	\end{itemize}
	\item \textbf{Encouraging Self-Employment}:
	\begin{itemize}
		\item Promoting micro-enterprises and entrepreneurship can boost MPCE for self-employed households.
	\end{itemize}
\end{itemize}





\newpage







\begin{itemize}
	\item \textbf{Household Types Represented}:
	\begin{itemize}
		\item Categories are: \textbf{Self-Employment (blue)}, \textbf{Regular Wage/Salary (light blue)}, \textbf{Casual Labour (gray)}, \textbf{Others (orange)}, and \textbf{All Categories (red)}.
		\item Each household type reflects the MPCE (Monthly Per Capita Expenditure) distribution in urban areas, measured in Indian Rupees.
	\end{itemize}
	\item \textbf{Key Observations by States}:
	\begin{itemize}
		\item Kerala:
		\begin{itemize}
			\item Highest MPCE across most household types, with the \textbf{Regular Wage/Salary} category standing out significantly.
		\end{itemize}
		\item Bihar:
		\begin{itemize}
			\item Exhibits the lowest MPCE across household types, highlighting significant economic disparities compared to other states.
		\end{itemize}
		\item Gujarat and Maharashtra:
		\begin{itemize}
			\item Showcase relatively balanced MPCE distributions among all household categories, suggesting stable urban income patterns.
		\end{itemize}
	\end{itemize}
	\item \textbf{Category Insights}:
	\begin{itemize}
		\item \textbf{Regular Wage/Salary}:
		\begin{itemize}
			\item This category consistently records the highest MPCE in most states, signifying the economic advantages of steady employment.
		\end{itemize}
		\item \textbf{Casual Labour}:
		\begin{itemize}
			\item Lowest MPCE, illustrating economic vulnerability due to irregular income patterns.
		\end{itemize}
		\item \textbf{Self-Employment}:
		\begin{itemize}
			\item Falls between the extremes of Regular Wage/Salary and Casual Labour, reflecting variability in income levels within this group.
		\end{itemize}
	\end{itemize}
\end{itemize}

\newpage

\begin{figure}[h]
	\centering
	\includegraphics[width=\textwidth]{D:/internship/Competition/Figure_70.png} % Replace with the file path of your image
	\caption{Stacked Bar Chart of MPCE by Household Type Across Indian States}
	\label{fig:MPCE_Stacked_Bar_Chart}
\end{figure}

\begin{figure}[h]
	\centering
	\includegraphics[width=\textwidth]{D:/internship/Competition/Figure_71.png} % Replace with the actual file path of the heatmap image
	\caption{Heatmap of Monthly Per Capita Expenditure (MPCE) by Household Type Across Indian States}
	\label{fig:Heatmap_MPCE_Household_Type}
\end{figure}

\newpage
\section*{Policy Implications}

\begin{itemize}
	\item \textbf{Support for Casual Labour Households}:
	\begin{itemize}
		\item Policies aimed at stabilizing income and expanding job opportunities can uplift this vulnerable segment.
	\end{itemize}
	\item \textbf{State-Specific Interventions}:
	\begin{itemize}
		\item Development programs in low-MPCE states like Bihar are crucial for reducing regional disparities.
	\end{itemize}
	\item \textbf{Encouraging Self-Employment}:
	\begin{itemize}
		\item Promoting small businesses and micro-enterprises can significantly boost income in the self-employed category.
	\end{itemize}
	\item \textbf{Role Models}:
	\begin{itemize}
		\item High-performing states like Kerala provide valuable lessons in achieving balanced urban economic growth.
	\end{itemize}
\end{itemize}






\subsection{ MPCE by Household Types}




\begin{itemize}
	\item \textbf{Chart Overview}:
	\begin{itemize}
		\item The box plot presents the \textbf{Monthly Per Capita Expenditure (MPCE)} in Indian Rupees (\textbf{Rs.}) across different household types in urban areas.
		\item The y-axis denotes MPCE ranging from 4000 to 12000 Rs.
		\item Categories analyzed: \textbf{Self-Employment, Regular Wage/Salary, Casual Labour, Others}, and \textbf{All Types}.
	\end{itemize}
	
	\item \textbf{Key Insights}:
	\begin{itemize}
		\item \textbf{Median MPCE Comparison}:
		\begin{itemize}
			\item \textbf{Regular Wage/Salary households} exhibit the highest median MPCE, indicating economic benefits from steady employment.
			\item \textbf{Casual Labour households} have the lowest median MPCE, reflecting income vulnerability and inconsistent earnings.
		\end{itemize}
		\item \textbf{Interquartile Range (IQR)}:
		\begin{itemize}
			\item The widest IQR is observed in the \textbf{Self-Employment category}, indicating high variability in expenditure patterns within this group.
			\item \textbf{Regular Wage/Salary} has a relatively narrow IQR, suggesting consistent expenditure among salaried households.
		\end{itemize}
		\item \textbf{Whiskers and Outliers}:
		\begin{itemize}
			\item \textbf{Casual Labour} and \textbf{Others} exhibit lower whiskers, reflecting minimal MPCE in some households.
			\item Outliers in \textbf{Self-Employment} and \textbf{Regular Wage/Salary} indicate households with exceptionally high expenditure.
		\end{itemize}
	\end{itemize}
\end{itemize}
\begin{figure}[h]
	\centering
	\includegraphics[width=\textwidth]{D:/internship/Competition/Figure_72.png} % Replace this with the correct file path for your box plot image
	\caption{Box Plot of MPCE Across Household Types in Indian Urban Areas}
	\label{fig:MPCE_Boxplot_Household_Types}
\end{figure}

\section*{Policy Implications}

\begin{itemize}
	\item \textbf{Support for Vulnerable Groups}:
	\begin{itemize}
		\item Policies aimed at stabilizing the incomes of \textbf{Casual Labour households} are vital for reducing expenditure disparities.
	\end{itemize}
	\item \textbf{Encourage Entrepreneurship}:
	\begin{itemize}
		\item Supporting \textbf{Self-Employment} through skill training and financial inclusion can enhance economic stability in this group.
	\end{itemize}
	\item \textbf{Inclusive Urban Growth}:
	\begin{itemize}
		\item Leveraging the steady expenditure patterns of salaried households can guide urban economic planning and investment strategies.
	\end{itemize}
\end{itemize}





\newpage





\subsection{Analysis of Top 5 States by MPCE in Self-Employment}

\begin{itemize}
	\item \textbf{Overview of the Chart}:
	\begin{itemize}
		\item The horizontal bar chart compares the Monthly Per Capita Expenditure (MPCE) for self-employment across five states in India.
		\item States listed include \textbf{Haryana}, \textbf{Tamil Nadu}, \textbf{Telangana}, \textbf{Punjab}, and \textbf{Kerala}.
	\end{itemize}
	\item \textbf{Key Observations}:
	\begin{itemize}
		\item \textbf{Highest MPCE (Haryana)}:
		\begin{itemize}
			\item Haryana leads with the highest MPCE value (close to Rs. 9000), indicating significant expenditure capacity among self-employed households.
		\end{itemize}
		\item \textbf{Lowest MPCE (Kerala)}:
		\begin{itemize}
			\item Kerala ranks fifth, with MPCE slightly below Rs. 7000, despite its reputation for strong economic growth and urban infrastructure.
		\end{itemize}
		\item \textbf{Regional Patterns}:
		\begin{itemize}
			\item Southern states (Tamil Nadu, Telangana, Kerala) display relatively uniform MPCE values, suggesting consistent economic conditions across the region.
			\item Northern states (Haryana and Punjab) show higher expenditure levels, which may reflect varying income distributions or economic dynamics.
		\end{itemize}
	\end{itemize}
\end{itemize}

\section*{Policy Implications}

\begin{itemize}
	\item \textbf{Boost Self-Employment in Lower-MPCE States}:
	\begin{itemize}
		\item Kerala could benefit from targeted programs supporting self-employment ventures, aiming to increase income levels and expenditure capacity.
	\end{itemize}
	\item \textbf{Learn from High-MPCE States}:
	\begin{itemize}
		\item Haryana's strong performance in self-employment expenditure can serve as a model for designing supportive policies in other states.
	\end{itemize}
	\item \textbf{Regional Equity}:
	\begin{itemize}
		\item Promoting balanced regional development to address disparities between northern and southern states could enhance overall economic stability.
	\end{itemize}
\end{itemize}
\begin{figure}[h]
	\centering
	\includegraphics[width=\textwidth]{D:/internship/Competition/Figure_73.png} % Update this path with the actual location of your bar chart
	\caption{Top 5 States by MPCE in Self-Employment}
	\label{fig:MPCE_Self_Employment_Top5}
\end{figure}


\newpage

\subsection{Analysis of Bottom 5 States by MPCE in Self-Employment}

\begin{itemize}
	\item \textbf{Overview of the Chart}:
	\begin{itemize}
		\item The horizontal bar chart represents the Monthly Per Capita Expenditure (MPCE) of self-employed households for the bottom five states in India.
		\item States listed from top to bottom are: \textbf{Madhya Pradesh, Chhattisgarh, Uttar Pradesh, Bihar, and Jharkhand}.
		\item The x-axis denotes MPCE values ranging from 0 to 5000 Rs., with bars shaded from dark red (highest MPCE) to light red (lowest MPCE).
	\end{itemize}
	\item \textbf{Key Observations}:
	\begin{itemize}
		\item \textbf{Lowest MPCE State (Jharkhand)}:
		\begin{itemize}
			\item Jharkhand exhibits the lowest MPCE, with values barely crossing the 3000 Rs. mark.
		\end{itemize}
		\item \textbf{Highest MPCE among Bottom 5 (Madhya Pradesh)}:
		\begin{itemize}
			\item Madhya Pradesh leads among the bottom five states, with MPCE values nearing 5000 Rs., indicating relatively better spending capacity.
		\end{itemize}
		\item \textbf{Regional Trends}:
		\begin{itemize}
			\item Eastern states like \textbf{Bihar} and \textbf{Jharkhand} consistently rank lower in MPCE, highlighting economic vulnerabilities in the region.
			\item Central states such as \textbf{Madhya Pradesh} and \textbf{Chhattisgarh} show slightly better MPCE values, suggesting moderate expenditure levels in these areas.
		\end{itemize}
	\end{itemize}
\end{itemize}

\section*{Policy Implications}

\begin{itemize}
	\item \textbf{Economic Support for Vulnerable States}:
	\begin{itemize}
		\item Targeted interventions in Bihar and Jharkhand can help boost incomes and improve expenditure levels among self-employed households.
	\end{itemize}
	\item \textbf{Strengthening Regional Economies}:
	\begin{itemize}
		\item Focused development programs in eastern and central regions, such as infrastructure investments and market access, can help reduce disparities.
	\end{itemize}
	\item \textbf{Promoting Self-Employment Ventures}:
	\begin{itemize}
		\item Incentives such as microfinance schemes and skill development initiatives can empower self-employed individuals and enhance spending capacity.
	\end{itemize}
\end{itemize}


\begin{figure}[h]
	\centering
	\includegraphics[width=\textwidth]{D:/internship/Competition/Figure_74.png} % Replace with the file path of your bar chart image
	\caption{Bottom 5 States by MPCE in Self-Employment}
	\label{fig:Bottom5_MPCE_SelfEmployment}
\end{figure}



\newpage
\subsection{Analysis of Top 5 States by MPCE in Regular Wage/Salary}

\begin{itemize}
	\item \textbf{Chart Overview}:
	\begin{itemize}
		\item The horizontal bar chart depicts the \textbf{Monthly Per Capita Expenditure (MPCE)} for regular wage/salary households in the top five states of India: \textbf{Karnataka}, \textbf{Telangana}, \textbf{Haryana}, \textbf{Tamil Nadu}, and \textbf{Kerala}.
		\item MPCE values range from 0 to 10,000 Rs. on the x-axis, represented by bars in varying shades of blue.
	\end{itemize}
	\item \textbf{Key Observations}:
	\begin{itemize}
		\item \textbf{Highest MPCE State (Karnataka)}:
		\begin{itemize}
			\item Karnataka leads with the highest MPCE, indicating strong economic conditions for salaried households.
		\end{itemize}
		\item \textbf{Lowest MPCE among Top 5 (Kerala)}:
		\begin{itemize}
			\item Kerala, while still among the top five, has the lowest MPCE of the group, suggesting relatively modest expenditure compared to the other states.
		\end{itemize}
		\item \textbf{Regional Insights}:
		\begin{itemize}
			\item Southern states like Karnataka, Tamil Nadu, and Kerala dominate the list, showcasing robust economic activity in the south.
			\item Telangana and Haryana reflect significant urban infrastructure and economic opportunities supporting salaried households.
		\end{itemize}
	\end{itemize}
\end{itemize}

\section*{Policy Implications}

\begin{itemize}
	\item \textbf{Promote Wage/Salaried Employment}:
	\begin{itemize}
		\item States with lower MPCE in wage/salary households could benefit from skill development programs and education initiatives to support steady employment opportunities.
	\end{itemize}
	\item \textbf{Regional Development Focus}:
	\begin{itemize}
		\item Lessons from Karnataka's strong economic performance can be utilized to strengthen salaried employment in other regions.
	\end{itemize}
	\item \textbf{Investment in Urban Infrastructure}:
	\begin{itemize}
		\item Enhancing infrastructure and employment ecosystems in lower-ranking states can bridge regional disparities in expenditure levels.
	\end{itemize}
\end{itemize}

\begin{figure}[h]
	\centering
	\includegraphics[width=\textwidth]{D:/internship/Competition/Figure_75.png} % Replace this path with the correct file path for your bar chart
	\caption{Top 5 States by MPCE in Regular Wage/Salary}
	\label{fig:Top5_MPCE_Wage_Salary}
\end{figure}





\newpage

\subsection{Analysis of Bottom 5 States by MPCE in Regular Wage/Salary}

\begin{itemize}
	\item \textbf{Overview of the Chart}:
	\begin{itemize}
		\item This bar chart compares the Monthly Per Capita Expenditure (MPCE) of regular wage/salary households across the five states with the lowest MPCE: \textbf{Uttar Pradesh, Madhya Pradesh, Jharkhand, Bihar, and Chhattisgarh}.
		\item The x-axis displays the MPCE in Indian Rupees (Rs.), ranging from 0 to 6000, while the y-axis lists the states.
	\end{itemize}
	\item \textbf{Key Observations}:
	\begin{itemize}
		\item \textbf{Highest MPCE among Bottom 5 (Uttar Pradesh)}:
		\begin{itemize}
			\item Uttar Pradesh has the highest MPCE in this group, nearing 6000 Rs.
		\end{itemize}
		\item \textbf{Lowest MPCE (Chhattisgarh)}:
		\begin{itemize}
			\item Chhattisgarh records the lowest MPCE, approximately around 4000 Rs., indicating relatively limited expenditure capacity.
		\end{itemize}
		\item \textbf{Regional Trends}:
		\begin{itemize}
			\item Northern and central states (e.g., Uttar Pradesh and Madhya Pradesh) show slightly better expenditure patterns compared to eastern states like Bihar and Jharkhand.
		\end{itemize}
	\end{itemize}
\end{itemize}

\section*{Policy Implications}

\begin{itemize}
	\item \textbf{Focus on Vulnerable States}:
	\begin{itemize}
		\item Bihar and Chhattisgarh could benefit from targeted economic policies that boost wages and provide better employment opportunities for regular wage/salary households.
	\end{itemize}
	\item \textbf{Support Infrastructure Development}:
	\begin{itemize}
		\item Investment in infrastructure and urban development in these states can help improve MPCE values.
	\end{itemize}
	\item \textbf{Skill Development Programs}:
	\begin{itemize}
		\item Introducing training and education initiatives to increase skill levels among workers may raise wage levels and, consequently, MPCE values.
	\end{itemize}
\end{itemize}

\begin{figure}[h]
	\centering
	\includegraphics[width=\textwidth]{D:/internship/Competition/Figure_76.png} % Replace this path with the correct file path for your bar chart
	\caption{Bottom 5 States by MPCE in Regular Wage/Salary}
	\label{fig:Bottom5_MPCE_Wage_Salary}
\end{figure}





\newpage

\subsection{Analysis of Top 5 States by MPCE in Casual Labour}

\begin{itemize}
	\item \textbf{Chart Overview}:
	\begin{itemize}
		\item The chart presents the Monthly Per Capita Expenditure (MPCE) for casual labor households in the top five states of India: \textbf{Telangana, Kerala, Tamil Nadu, Andhra Pradesh, and Karnataka}.
		\item The x-axis represents MPCE values in Indian Rupees (\textbf{Rs.}), while the y-axis lists the states in descending order based on expenditure.
	\end{itemize}
	\item \textbf{Key Observations}:
	\begin{itemize}
		\item \textbf{Highest MPCE (Telangana)}:
		\begin{itemize}
			\item Telangana leads with the highest MPCE among casual labor households, reflecting better economic conditions and spending capacity.
		\end{itemize}
		\item \textbf{Lowest MPCE (Karnataka)}:
		\begin{itemize}
			\item Karnataka ranks fifth, indicating relatively lower expenditure levels compared to other states in the list.
		\end{itemize}
		\item \textbf{Regional Trends}:
		\begin{itemize}
			\item Southern states dominate the list, showcasing robust economic activity and support for casual laborers in the region.
		\end{itemize}
	\end{itemize}
\end{itemize}

\section*{Policy Implications}

\begin{itemize}
	\item \textbf{Strengthening Casual Labor Support}:
	\begin{itemize}
		\item Targeted policies in Karnataka could aim to increase incomes and improve living conditions for casual labor households.
	\end{itemize}
	\item \textbf{Leveraging High-MPCE States}:
	\begin{itemize}
		\item States like Telangana and Kerala could serve as benchmarks for developing welfare programs in other regions.
	\end{itemize}
	\item \textbf{Encouraging Regional Development}:
	\begin{itemize}
		\item Expanding economic opportunities for casual laborers across states can help bridge disparities in expenditure levels.
	\end{itemize}
\end{itemize}

\begin{figure}[h]
	\centering
	\includegraphics[width=\textwidth]{D:/internship/Competition/Figure_78.png} % Replace this path with the actual file location of your chart image
	\caption{Top 5 States by MPCE in Casual Labour}
	\label{fig:Top5_MPCE_Casual_Labour}
\end{figure}
\newpage




\subsection{Analysis of Bottom 5 States by MPCE in Casual Labour}

\begin{itemize}
	\item \textbf{Overview of the Chart}:
	\begin{itemize}
		\item This horizontal bar chart displays the \textbf{Monthly Per Capita Expenditure (MPCE)} for casual labour households across five states with the lowest MPCE.
		\item States listed from top to bottom are: \textbf{Madhya Pradesh, West Bengal, Odisha, Chhattisgarh, and Jharkhand}.
		\item MPCE values are shown on the x-axis in Indian Rupees (\textbf{Rs.}), ranging from 0 to 4000.
	\end{itemize}
	\item \textbf{Key Observations}:
	\begin{itemize}
		\item \textbf{Highest MPCE among Bottom 5 (Madhya Pradesh)}:
		\begin{itemize}
			\item Madhya Pradesh has the highest MPCE among the bottom 5 states, nearing 3500 Rs.
		\end{itemize}
		\item \textbf{Lowest MPCE (Jharkhand)}:
		\begin{itemize}
			\item Jharkhand records the lowest MPCE, barely exceeding 2500 Rs., indicating significant economic challenges for casual labour households.
		\end{itemize}
		\item \textbf{Regional Trends}:
		\begin{itemize}
			\item Eastern states like \textbf{West Bengal}, \textbf{Odisha}, and \textbf{Jharkhand} consistently rank lower, underscoring economic vulnerabilities in the region.
			\item Central states such as \textbf{Madhya Pradesh} and \textbf{Chhattisgarh} show slightly better MPCE values but still fall below the national average.
		\end{itemize}
	\end{itemize}
\end{itemize}

\section*{Policy Implications}

\begin{itemize}
	\item \textbf{Economic Support for Vulnerable States}:
	\begin{itemize}
		\item Policies targeting income stabilization and welfare for casual labourers in Jharkhand and Odisha can reduce disparities.
	\end{itemize}
	\item \textbf{Infrastructure Development}:
	\begin{itemize}
		\item Investments in rural and urban infrastructure in these states can provide better employment opportunities and improve household expenditures.
	\end{itemize}
	\item \textbf{Skill Development and Education Initiatives}:
	\begin{itemize}
		\item Training programs aimed at enhancing skill levels of casual labourers can help elevate income levels and, consequently, MPCE values.
	\end{itemize}
\end{itemize}
\begin{figure}[h]
	\centering
	\includegraphics[width=\textwidth]{D:/internship/Competition/Figure_80.png} % Replace this path with the actual file path of the chart image
	\caption{Bottom 5 States by MPCE in Casual Labour}
	\label{fig:Bottom5_MPCE_CasualLabour}
\end{figure}





\newpage

\subsection{Analysis of Top 5 States by MPCE in Others}

\begin{itemize}
	\item \textbf{Chart Overview}:
	\begin{itemize}
		\item This bar chart represents the Monthly Per Capita Expenditure (MPCE) for the "Others" category in the top five states of India.
		\item The states listed are: \textbf{Telangana, Rajasthan, Haryana, Maharashtra, and Karnataka}.
		\item The x-axis displays MPCE in Indian Rupees (\textbf{Rs.}), while the y-axis lists the states in descending order of expenditure.
	\end{itemize}
	\item \textbf{Key Observations}:
	\begin{itemize}
		\item \textbf{Highest MPCE (Telangana)}:
		\begin{itemize}
			\item Telangana leads the list with the highest MPCE, showcasing strong expenditure in this category.
		\end{itemize}
		\item \textbf{Lowest MPCE among the Top 5 (Karnataka)}:
		\begin{itemize}
			\item Karnataka ranks fifth but still exhibits significant MPCE, reflecting robust spending capacity.
		\end{itemize}
		\item \textbf{Regional Patterns}:
		\begin{itemize}
			\item The inclusion of both southern states (e.g., Telangana and Karnataka) and northern/western states (e.g., Rajasthan and Haryana) highlights regional diversity in the "Others" category.
		\end{itemize}
	\end{itemize}
\end{itemize}

\section*{Policy Implications}

\begin{itemize}
	\item \textbf{Leveraging High-Performing States}:
	\begin{itemize}
		\item Telangana’s strong performance in this category could be studied for replicable strategies to boost MPCE in other states.
	\end{itemize}
	\item \textbf{Encouraging Balanced Growth}:
	\begin{itemize}
		\item States like Rajasthan and Haryana, which perform well in this category, could focus on translating these strengths into improved overall economic indicators.
	\end{itemize}
	\item \textbf{Strengthening the "Others" Category Nationwide}:
	\begin{itemize}
		\item Investment in niche sectors or consumption categories falling under "Others" could spur economic growth across regions.
	\end{itemize}
\end{itemize}


\begin{figure}[h]
	\centering
	\includegraphics[width=\textwidth]{D:/internship/Competition/Figure_81.png} % Replace this path with the correct location of your chart file
	\caption{Top 5 States by MPCE in Others}
	\label{fig:Top5_MPCE_Others}
\end{figure}


\newpage

\subsection{Analysis of Bottom 5 States by MPCE in Others}

\begin{itemize}
	\item \textbf{Chart Overview}:
	\begin{itemize}
		\item The horizontal bar chart showcases the Monthly Per Capita Expenditure (MPCE) for the bottom five states in the "Others" category.
		\item States are listed from \textbf{Madhya Pradesh} to \textbf{Chhattisgarh}, with the x-axis representing MPCE in Indian Rupees (Rs.) from 0 to 8000.
		\item The y-axis labels the states, while the bar colors range from dark to light red, indicating expenditure levels.
	\end{itemize}
	\item \textbf{Key Observations}:
	\begin{itemize}
		\item \textbf{Highest MPCE among Bottom 5 (Madhya Pradesh)}:
		\begin{itemize}
			\item Madhya Pradesh leads this group with an MPCE close to 8000 Rs., suggesting relatively better expenditure conditions.
		\end{itemize}
		\item \textbf{Lowest MPCE (Chhattisgarh)}:
		\begin{itemize}
			\item Chhattisgarh registers the lowest MPCE, just above 6000 Rs., indicating significant expenditure challenges.
		\end{itemize}
		\item \textbf{Regional Trends}:
		\begin{itemize}
			\item Northern and central states (e.g., Madhya Pradesh and Uttar Pradesh) exhibit higher MPCE within this group.
			\item Eastern states, such as Odisha and Bihar, reflect consistent low MPCE values, highlighting regional disparities.
		\end{itemize}
	\end{itemize}
\end{itemize}

\section*{Policy Implications}

\begin{itemize}
	\item \textbf{Targeted Economic Policies}:
	\begin{itemize}
		\item States like Chhattisgarh and Bihar need focused welfare programs and income stabilization initiatives to uplift MPCE in the "Others" category.
	\end{itemize}
	\item \textbf{Infrastructure Investments}:
	\begin{itemize}
		\item Enhancing urban and rural infrastructure can provide equitable access to opportunities, increasing expenditure capacity.
	\end{itemize}
	\item \textbf{Promoting Regional Equity}:
	\begin{itemize}
		\item Addressing regional disparities through economic development programs in low-MPCE states can ensure balanced growth.
	\end{itemize}
\end{itemize}



\begin{figure}[h]
	\centering
	\includegraphics[width=\textwidth]{D:/internship/Competition/Figure_82.png} % Update this path to match the file location of the bar chart image
	\caption{Bottom 5 States by MPCE in Others}
	\label{fig:Bottom5_MPCE_Others}
\end{figure}



\newpage

\subsection{Analysis of Top 5 States by MPCE in All}

\begin{itemize}
	\item \textbf{Chart Overview}:
	\begin{itemize}
		\item The chart highlights the Monthly Per Capita Expenditure (MPCE) across the top 5 states for all household types combined.
		\item States included are \textbf{Telangana, Haryana, Tamil Nadu, Karnataka, and Kerala}, listed in descending order based on MPCE values.
		\item The x-axis shows MPCE in Indian Rupees (Rs.), while the bars are shaded in varying tones of blue to represent expenditure levels.
	\end{itemize}
	\item \textbf{Key Observations}:
	\begin{itemize}
		\item \textbf{Highest MPCE (Telangana)}:
		\begin{itemize}
			\item Telangana has the highest MPCE, exceeding 8500 Rs., reflecting strong expenditure trends across all household categories.
		\end{itemize}
		\item \textbf{Lowest MPCE (Kerala)}:
		\begin{itemize}
			\item Kerala ranks fifth, with an MPCE below 7500 Rs., indicating stable but slightly lower spending compared to other top-performing states.
		\end{itemize}
		\item \textbf{Regional Insights}:
		\begin{itemize}
			\item Southern states like Telangana, Tamil Nadu, Karnataka, and Kerala dominate the list, highlighting consistent economic advantages in the region.
			\item Haryana represents a strong contender from northern India, showcasing robust expenditure trends.
		\end{itemize}
	\end{itemize}
\end{itemize}

\section*{Policy Implications}

\begin{itemize}
	\item **Encouraging High-MPCE States**:
	\begin{itemize}
		\item States with high MPCE, such as Telangana, could be studied for strategies to enhance economic growth across other regions.
	\end{itemize}
	\item **Fostering Regional Equality**:
	\begin{itemize}
		\item Bridging the gap between high-MPCE states and others can be achieved through targeted economic development and infrastructure investments.
	\end{itemize}
	\item **Strengthening Household Expenditure Nationwide**:
	\begin{itemize}
		\item Promoting balanced expenditure across regions and household types can ensure sustainable economic growth.
	\end{itemize}
\end{itemize}


\begin{figure}[h]
	\centering
	\includegraphics[width=\textwidth]{D:/internship/Competition/Figure_83.png} % Replace this path with the actual file location of your image
	\caption{Top 5 States by MPCE in All}
	\label{fig:Top5_MPCE_All}
\end{figure}





\newpage

\subsection{Analysis of Top 5 States by MPCE in All}

\begin{itemize}
	\item \textbf{Chart Overview}:
	\begin{itemize}
		\item The chart highlights the Monthly Per Capita Expenditure (MPCE) across the top 5 states for all household types combined.
		\item States included are \textbf{Telangana, Haryana, Tamil Nadu, Karnataka, and Kerala}, listed in descending order based on MPCE values.
		\item The x-axis shows MPCE in Indian Rupees (Rs.), while the bars are shaded in varying tones of blue to represent expenditure levels.
	\end{itemize}
	\item \textbf{Key Observations}:
	\begin{itemize}
		\item \textbf{Highest MPCE (Telangana)}:
		\begin{itemize}
			\item Telangana has the highest MPCE, exceeding 8500 Rs., reflecting strong expenditure trends across all household categories.
		\end{itemize}
		\item \textbf{Lowest MPCE (Kerala)}:
		\begin{itemize}
			\item Kerala ranks fifth, with an MPCE below 7500 Rs., indicating stable but slightly lower spending compared to other top-performing states.
		\end{itemize}
		\item \textbf{Regional Insights}:
		\begin{itemize}
			\item Southern states like Telangana, Tamil Nadu, Karnataka, and Kerala dominate the list, highlighting consistent economic advantages in the region.
			\item Haryana represents a strong contender from northern India, showcasing robust expenditure trends.
		\end{itemize}
	\end{itemize}
\end{itemize}

\section*{Policy Implications}

\begin{itemize}
	\item **Encouraging High-MPCE States**:
	\begin{itemize}
		\item States with high MPCE, such as Telangana, could be studied for strategies to enhance economic growth across other regions.
	\end{itemize}
	\item **Fostering Regional Equality**:
	\begin{itemize}
		\item Bridging the gap between high-MPCE states and others can be achieved through targeted economic development and infrastructure investments.
	\end{itemize}
	\item **Strengthening Household Expenditure Nationwide**:
	\begin{itemize}
		\item Promoting balanced expenditure across regions and household types can ensure sustainable economic growth.
	\end{itemize}
\end{itemize}


\begin{figure}[h]
	\centering
	\includegraphics[width=\textwidth]{D:/internship/Competition/Figure_84.png} % Replace this path with the actual file location of your image
	\caption{Top 5 States by MPCE in All}
	\label{fig:Top5_MPCE_All}
\end{figure}

\newpage


\section{Analysis of Rural and Urban Gini Coefficients by State Source(Table 3.15: Gini coefficient of total consumption expenditure in 2011-12, 2022-23 and 2023-24, All-India) (Table	3.16: Gini coefficient of total consumption expenditurein 2023-24, major states)}


	
	
	\begin{itemize}
		\item \textbf{Higher Urban Gini Coefficients:}
		\begin{itemize}
			\item Urban Gini coefficients (orange bars) exceed rural coefficients (purple bars) in most states, suggesting greater income inequality in urban areas.
			\item States like \textbf{Kerala} and \textbf{Punjab} display the largest urban-rural inequality gaps, highlighting significant economic disparities within these regions.
		\end{itemize}
		\item \textbf{Relatively Equal States:}
		\begin{itemize}
			\item \textbf{Chhattisgarh} and \textbf{Rajasthan} exhibit smaller differences between rural and urban Gini coefficients, reflecting more balanced income distributions.
		\end{itemize}
		\item \textbf{All-India Trend:}
		\begin{itemize}
			\item At the national level, urban areas generally report higher inequality compared to rural areas, as reflected in the larger Gini coefficients for urban settings.
		\end{itemize}
	\end{itemize}
	
	\section*{Key Trends and Insights}
	\begin{itemize}
		\item \textbf{Highest Urban Gini Coefficients:}
		\begin{itemize}
			\item States like \textbf{Kerala}, \textbf{Punjab}, and \textbf{Haryana} show the highest urban inequality, reflecting concentrated wealth disparities in urban areas.
		\end{itemize}
		\item \textbf{Lowest Rural Gini Coefficients:}
		\begin{itemize}
			\item \textbf{Chhattisgarh}, \textbf{Assam}, and \textbf{Bihar} exhibit the least income inequality in rural regions, highlighting more equitable rural income distributions.
		\end{itemize}
	\end{itemize}
	
	\section*{Policy Implications}
	\begin{itemize}
		\item \textbf{Urban Challenges:}
		\begin{itemize}
			\item Address urban inequality through equitable economic opportunities, access to education, and affordable housing policies in high-inequality states like \textbf{Kerala} and \textbf{Punjab}.
		\end{itemize}
		\item \textbf{Rural Stability:}
		\begin{itemize}
			\item Despite relatively low rural inequality, investment in rural infrastructure, skills development, and industries can further reduce poverty and improve quality of life.
		\end{itemize}
		\item \textbf{Balanced Growth:}
		\begin{itemize}
			\item States like \textbf{Rajasthan} and \textbf{Chhattisgarh} can serve as models for achieving balanced urban-rural development, promoting inclusive policies nationwide.
		\end{itemize}
	\end{itemize}
	



	\begin{figure}[h]
	\centering
	\includegraphics[width=\textwidth]{D:/internship/Competition/Figure_62.png} % Replace this path with the actual file path of your image
	\caption{Analysis of Rural and Urban Gini Coefficients by State}
	\label{fig:Analysis of Rural and Urban Gini Coefficients by State}
\end{figure}



\section*{Gini Coefficient Formula}

\begin{equation}
	G = \frac{\sum_{i=1}^{n} \sum_{j=1}^{n} |x_i - x_j|}{2n^2\mu}
\end{equation}

Where:
\begin{itemize}
	\item $G$: Gini coefficient.
	\item $n$: Total number of observations.
	\item $x_i$ and $x_j$: Individual income or expenditure observations.
	\item $\mu$: Mean income or expenditure.
\end{itemize}



\section{Analysis of Urban Non-Food Expenditure Distribution Across Indian States}

\begin{itemize}
	\item \textbf{Categories of Non-Food Expenditure:}
	\begin{itemize}
		\item The expenditure is categorized into:
		\begin{itemize}
			\item \textbf{Fuel \& Light}: Spending on energy resources and utilities.
			\item \textbf{Conveyance}: Transportation expenses.
			\item \textbf{Education}: Costs for schooling, higher education, and resources.
			\item \textbf{Medical}: Healthcare expenses including medicines and treatments.
			\item \textbf{Miscellaneous Goods \& Entertainment}: Leisure and general items.
			\item \textbf{Clothing, Bedding \& Footwear}: Apparel and similar purchases.
			\item \textbf{Durable Goods}: Long-term assets like appliances and furniture.
			\item \textbf{Others}: Residual expenditures.
		\end{itemize}
	\end{itemize}
	\begin{figure}[h]
		\centering
		\includegraphics[width=\textwidth]{D:/internship/Competition/Figure_85.png} % Replace this path with the actual file path of your image
		\caption{Urban Non-Food Expenditure Distribution Across Indian States}
		\label{fig:Urban_Non_Food_Expenditure}
	\end{figure}
	
	\item \textbf{Key State-Level Insights:}
	\begin{itemize}
		\item \textbf{High-Spending States:}
		\begin{itemize}
			\item \textbf{Kerala \& Tamil Nadu:} Significant allocation to \textbf{Education} and \textbf{Medical} expenses, reflecting high development indicators.
			\item \textbf{Gujarat \& Haryana:} Higher spending on \textbf{Miscellaneous Goods \& Entertainment}, showcasing lifestyle-driven consumption patterns.
		\end{itemize}
		\item \textbf{Low-Spending States:}
		\begin{itemize}
			\item \textbf{Bihar, Odisha, \& Jharkhand:} Greater focus on \textbf{Fuel \& Light}, highlighting basic living needs.
			\item \textbf{Chhattisgarh \& Madhya Pradesh:} Spending leans toward healthcare and basic survival categories.
		\end{itemize}
	\end{itemize}
	
	\item \textbf{All-India Trends:}
	\begin{itemize}
		\item \textbf{Fuel \& Light:} A substantial share across all states, indicating the importance of energy and utilities.
		\item \textbf{Education:} Higher spending in developed states like Kerala, while lagging in regions like Bihar.
		\item \textbf{Medical:} Rising healthcare costs dominate household budgets universally.
		\item \textbf{Durable Goods:} Significant variations, with advanced states allocating more compared to underdeveloped ones.
	\end{itemize}
\end{itemize}

\section*{Policy Implications}

\begin{itemize}
	\item \textbf{Address Regional Disparities:}
	\begin{itemize}
		\item Eastern states like Bihar and Odisha need targeted development in education and durable goods spending.
	\end{itemize}
	\item \textbf{Foster Human Development:}
	\begin{itemize}
		\item Public investment in \textbf{education} and \textbf{healthcare} can reduce inequalities and improve economic outcomes.
	\end{itemize}
	\item \textbf{Encourage Balanced Growth:}
	\begin{itemize}
		\item Southern states serve as models for balancing essential and lifestyle expenditures.
	\end{itemize}
\end{itemize}




	
	\subsection{ Analysis of Gini Coefficient Disparity (Urban - Rural)}
	
	\begin{itemize}
		\item \textbf{Chart Overview:}
		\begin{itemize}
			\item The bar chart represents the urban-rural Gini gap for various Indian states, quantifying disparities in income inequality.
			\item The x-axis lists states such as Haryana, Telangana, West Bengal, Jharkhand, Uttar Pradesh, Odisha, and others.
			\item The y-axis indicates the Gini coefficient difference, ranging from 0.00 to 0.10.
		\end{itemize}
		
		\item \textbf{Observations:}
		\begin{itemize}
			\item \textbf{Highest Urban-Rural Gap (Haryana):}
			\begin{itemize}
				\item Haryana has the largest Gini disparity (approximately 0.10), highlighting significant inequality in urban areas compared to rural ones.
			\end{itemize}
			
			\item \textbf{Lowest Urban-Rural Gap (Kerala):}
			\begin{itemize}
				\item Kerala shows the smallest Gini gap (slightly above 0.00), indicating relative equality between urban and rural income distributions.
				\item States like Punjab and Gujarat exhibit moderate urban-rural inequality gaps.
			\end{itemize}
			
			\item \textbf{All-India Average:}
			\begin{itemize}
				\item The national average urban-rural gap falls between 0.04 and 0.05, reinforcing the general trend of greater urban inequality.
			\end{itemize}
		\end{itemize}
		
		\item \textbf{Key Trends}
		\begin{itemize}
			\item \textbf{Urban Dominance:}
			\begin{itemize}
				\item States with higher urban Gini coefficients face challenges in managing wealth concentration and access to resources.
			\end{itemize}
			
			\item \textbf{Rural Stability:}
			\begin{itemize}
				\item Rural regions generally display lower inequality due to smaller income ranges and subsistence-based economies.
			\end{itemize}
		\end{itemize}
		
		\item \textbf{Policy Implications}
		\begin{itemize}
			\item \textbf{Targeted Programs for High-Gap States:}
			\begin{itemize}
				\item Initiatives like equitable urban development and skill-building programs could alleviate disparities in states like Haryana and West Bengal.
			\end{itemize}
			
			\item \textbf{Focus on Low-Gap States:}
			\begin{itemize}
				\item States like Kerala can maintain their balanced growth model, serving as a benchmark for equitable development.
			\end{itemize}
			
			\item \textbf{National-Level Strategies:}
			\begin{itemize}
				\item Policies fostering inclusive growth, such as investments in education, health, and infrastructure, can reduce the urban-rural divide across India.
			\end{itemize}
		\end{itemize}
	\end{itemize}
	



\begin{figure}[h]
	\centering
	\includegraphics[width=\textwidth]{D:/internship/Competition/Figure_63.png} % Replace this path with the actual file path of your image
	\caption{Analysis of Rural and Urban Gini Coefficients by State}
	\label{fig:Analysis of Rural and Urban Gini Coefficients by State}
\end{figure}







	
	\subsection{ Rural vs Urban Gini Coefficients by State}
	
	\subsection*{Plot Overview}
	\begin{itemize}
		\item The scatter plot compares Rural Gini Coefficients (x-axis) with Urban Gini Coefficients (y-axis) for Indian states.
		\item The data points are color-coded to indicate different states, with labels for clarity.
	\end{itemize}
	
	\subsection*{General Observations}
	\begin{itemize}
		\item Urban Gini coefficients are consistently higher than rural coefficients, reinforcing the trend that income inequality is greater in urban areas.
		\item States like \textbf{Kerala} and \textbf{Punjab} show higher overall Gini coefficients, indicating significant economic disparity in both rural and urban settings.
		\item \textbf{Chhattisgarh} and \textbf{Bihar} exhibit relatively low Gini coefficients, suggesting more equitable income distribution in these states.
	\end{itemize}
	
	\subsection*{Correlation}
	\begin{itemize}
		\item A positive correlation is observed, meaning states with higher rural inequality also tend to have higher urban inequality.
		\item However, the disparity is more pronounced in urban areas as seen from the wider spread of data along the y-axis.
	\end{itemize}
	
	\section*{Key Insights}
	
	\subsection*{High Inequality States}
	\begin{itemize}
		\item Kerala and Punjab show prominent disparities and require targeted policies to address inequality in both urban and rural areas.
	\end{itemize}
	
	\subsection*{Equitable Distribution States}
	\begin{itemize}
		\item States like Chhattisgarh and Bihar may serve as examples for designing equitable income distribution policies, despite their lower overall economic output.
	\end{itemize}
	
	\subsection*{Urban Focus}
	\begin{itemize}
		\item The consistently higher urban Gini coefficients indicate a need for better income redistribution mechanisms and inclusive urban development strategies.
	\end{itemize}
	

   \section*{Policy Implications}
   
   \begin{itemize}
   	\item \textbf{Focus on High-Inequality States:}
   	\begin{itemize}
   		\item States with high Gini coefficients require targeted interventions, such as investment in education, healthcare, and job creation, to reduce income disparities.
   	\end{itemize}
   	
   	\item \textbf{Urban Development Policies:}
   	\begin{itemize}
   		\item Developing affordable housing, skill training programs, and employment opportunities can help tackle urban inequality.
   	\end{itemize}
   	
   	\item \textbf{Addressing Rural Gaps:}
   	\begin{itemize}
   		\item Promoting sustainable agricultural development and rural infrastructure can reduce rural income disparities.
   	\end{itemize}
   \end{itemize}
   


\begin{figure}[h]
	\centering
	\includegraphics[width=\textwidth]{D:/internship/Competition/Figure_64.png} % Replace this path with the actual file path of your image
	\caption{Rural vs Urban Gini Coefficients by State}
	\label{fig:Rural vs Urban Gini Coefficients by State}
\end{figure}






	
	\subsection{Analysis of Urban Gini Coefficients}
	
	\subsection{Top 5 States with Highest Urban Gini (Left Chart)}
	\begin{itemize}
		\item Jharkhand, Haryana, Karnataka, Maharashtra, and Odisha exhibit the highest Urban Gini coefficients.
		\item These states show significant income inequality in urban areas, reflecting uneven wealth distribution and access to resources.
	\end{itemize}
	
	\subsection{Bottom 5 States with Lowest Urban Gini (Right Chart)}
	\begin{itemize}
		\item Punjab, Bihar, Gujarat, Andhra Pradesh, and Assam have the lowest Urban Gini coefficients.
		\item These states demonstrate more equitable income distribution in urban areas compared to others.
	\end{itemize}
	
	\section{Key Observations}
	\begin{itemize}
		\item \textbf{Inequality in High-Gini States:}
		\begin{itemize}
			\item Urban income disparities in states like Jharkhand and Haryana may stem from uneven development, limited infrastructure, and disparities in job opportunities.
		\end{itemize}
		\item \textbf{Equality in Low-Gini States:}
		\begin{itemize}
			\item States such as Punjab and Assam reflect more balanced economic conditions, possibly due to inclusive policies and evenly distributed urban infrastructure.
		\end{itemize}
	\end{itemize}
	
	\section{Policy Implications}
	\begin{itemize}
		\item \textbf{Address High Inequality:}
		\begin{itemize}
			\item High-Gini states may benefit from programs targeting wealth redistribution, improved access to education, and economic development in urban areas.
		\end{itemize}
		\item \textbf{Encourage Low-Gini States:}
		\begin{itemize}
			\item Low-Gini states could serve as models for fostering inclusive economic policies that promote equity in urban settings.
		\end{itemize}
	\end{itemize}
	


\begin{figure}[h]
	\centering
	\includegraphics[width=\textwidth]{D:/internship/Competition/Figure_67.png} % Replace this path with the actual file path of your image
	\caption{Urban Gini Coefficients in Top 5 and Bottom 5 States}
	\label{fig:Urban Gini Coefficients in Top 5 and Bottom 5 States}
\end{figure}
















\newpage


\section{Electricity Equity in Action: Sectoral Zero-Deprivation Achievements}
	\begin{table}[htbp]
		\centering
		\small
	\textbf{	\caption{State-wise Household Electricity Deprivation (HCES 2023-24) \\  Rural=1 and Urban=2}}
		\begin{tabular}{llSSS}
			\toprule
			\textbf{State} & \textbf{Sector} & \textbf{Total Households} & \textbf{Deprived Households} & \textbf{Deprivation Rate (\%)} \\
			\midrule
			Meghalaya & 1 & 563295 & 24510 & 4.35 \\
			Jharkhand & 1 & 5709801 & 99983 & 1.75 \\
			Uttar Pradesh & 1 & 30398615 & 368267 & 1.21 \\
			Tripura & 1 & 804896 & 9710 & 1.21 \\
			Gujarat & 1 & 7071053 & 80261 & 1.14 \\
			Chhattisgarh & 1 & 4591188 & 49844 & 1.09 \\
			Odisha & 1 & 9384147 & 98502 & 1.05 \\
			Assam & 1 & 6248139 & 62332 & 0.998 \\
		    Rajasthan & 1 & 10936974 & 96940 & 0.886 \\
		    West Bengal & 1 & 16976256 & 144533 & 0.851 \\
		    Maharashtra & 2 & 12562136 & 103050 & 0.820 \\
		    Manipur & 2 & 183691 & 1467 & 0.799 \\
		    Andaman \& N. Island & 1 & 76438 & 600 & 0.785 \\
		    Manipur & 1 & 499367 & 3833 & 0.768 \\
		    Tripura & 2 & 201873 & 1033 & 0.512 \\
		    Andhra Pradesh & 2 & 4515617 & 22659 & 0.502 \\
		    Maharashtra & 1 & 14768694 & 70186 & 0.475 \\
		    Meghalaya & 2 & 110882 & 510 & 0.460 \\
		    Haryana & 2 & 2554196 & 11289 & 0.442 \\
		    Tamilnadu & 2 & 10686204 & 44057 & 0.412 \\
		    Bihar & 2 & 2167306 & 8855 & 0.409 \\
		    Madhya Pradesh & 2 & 4484483 & 16287 & 0.363 \\
		    Tamilnadu & 1 & 11176960 & 39823 & 0.356 \\
		    West Bengal & 2 & 8004246 & 27761 & 0.347 \\
		    Andaman \& N. Island & 2 & 50536 & 156 & 0.309 \\
		    Uttar Pradesh & 2 & 9075344 & 27668 & 0.305 \\
		    D \& N. Haveli \& Daman \& Diu & 2 & 131873 & 381 & 0.289 \\
		    Puduchery & 1 & 106244 & 295 & 0.277 \\
		    Gujarat & 2 & 6551886 & 17418 & 0.266 \\
		    Assam & 2 & 981145 & 2360 & 0.241 \\
		    Punjab & 1 & 3804113 & 7127 & 0.187 \\
		    Jammu \& Kashmir & 1 & 1799259 & 3312 & 0.184 \\
		    Uttarakhand & 1 & 1856434 & 3244 & 0.175 \\
		    Arunachal Pradesh & 1 & 217653 & 377 & 0.173 \\
		    Andhra Pradesh & 1 & 9468709 & 16205 & 0.171 \\
		    Himachal Pradesh & 2 & 243984 & 380 & 0.156 \\
		    Bihar & 1 & 20151145 & 30937 & 0.154 \\
		    Telangana & 2 & 5407536 & 7771 & 0.144 \\
		    Odisha & 2 & 2193468 & 2870 & 0.131 \\
		    Rajasthan & 2 & 4247935 & 5512 & 0.130 \\
		    Madhya Pradesh & 1 & 11266164 & 13161 & 0.117 \\
		    Delhi & 2 & 3643667 & 3332 & 0.091 \\
		    Kerala & 1 & 4886446 & 4000 & 0.082 \\
		    Sikkim & 1 & 120725 & 95 & 0.078 \\
		    Haryana & 1 & 3627657 & 2724 & 0.075 \\
		    Chhattisgarh & 2 & 1332890 & 959 & 0.072 \\
		    Karnataka & 1 & 9300573 & 4904 & 0.053 \\
		    Kerala & 2 & 4925091 & 2560 & 0.052 \\
		    Punjab & 2 & 2744353 & 815 & 0.030 \\
			\bottomrule
		\end{tabular}
	\end{table}
	



\begin{table}[htbp]
	\centering
	\small
	\caption{States and Sectors with Zero Electricity-Deprived Households}
	\begin{tabular}{lrrr}
		\toprule
		State & Sector & Total Households & Deprived Households \\
		\midrule
		Jammu \& Kashmir & 2 & 566668 & 0 \\
		Himachal Pradesh & 1 & 1557010 & 0 \\
		Chandigarh & 1 & 11313 & 0 \\
		Chandigarh & 2 & 242733 & 0 \\
		Uttarakhand & 2 & 754955 & 0 \\
		Delhi & 1 & 89607 & 0 \\
		Sikkim & 2 & 57498 & 0 \\
		Arunachal Pradesh & 2 & 62305 & 0 \\
		Nagaland & 1 & 276945 & 0 \\
		Nagaland & 2 & 120338 & 0 \\
		Mizoram & 1 & 112840 & 0 \\
		Mizoram & 2 & 102746 & 0 \\
		Jharkhand & 2 & 1615380 & 0 \\
		D \& N. Haveli \& Daman \& Diu & 1 & 54438 & 0 \\
		Karnataka & 2 & 6880675 & 0 \\
		Goa & 1 & 134511 & 0 \\
		Goa & 2 & 223872 & 0 \\
		Lakshadweep & 1 & 3017 & 0 \\
		Lakshadweep & 2 & 9344 & 0 \\
		Puduchery & 2 & 240645 & 0 \\
		Telangana & 1 & 5689698 & 0 \\
		Ladakh & 1 & 43519 & 0 \\
		Ladakh & 2 & 9736 & 0 \\
		\bottomrule
	\end{tabular}
\end{table}




India’s commitment to providing universal, inclusive, and reliable electricity access has led to remarkable progress in several states and union territories (UTs). Despite this, electricity deprivation remains a challenge in certain regions, with varying degrees of intensity across rural (sector 1) and urban (sector 2) areas. While many states and UTs have made notable strides, there are still regions where a significant portion of households are deprived of electricity.

\subsection*{Rural Electrification: Reaching the Unreachable}

States such as \textbf{Himachal Pradesh}, \textbf{Nagaland}, \textbf{Mizoram}, \textbf{Telangana}, \textbf{Goa}, \textbf{Ladakh}, and \textbf{Dadra \& Nagar Haveli and Daman \& Diu} have reported no rural households facing electricity deprivation. This achievement reflects substantial progress in overcoming geographical challenges through innovative, decentralized energy solutions, such as solar micro-grids and stand-alone systems.

However, states like \textbf{Meghalaya}, \textbf{Jharkhand}, and \textbf{Uttar Pradesh} continue to experience higher deprivation rates (approximately 4.35\%, 1.75\%, and 1.21\%, respectively). These disparities highlight the need for targeted interventions to address localized barriers, such as remote geography, inadequate infrastructure, and socio-economic exclusion. 

\subsection*{Urban Electrification: High Coverage with Room for Equity}

Urban areas in \textbf{Delhi}, \textbf{Chandigarh}, \textbf{Sikkim}, \textbf{Jharkhand}, \textbf{Karnataka}, and \textbf{Puduchery} have achieved near-total electricity access (100\%), reflecting the success of urban infrastructure development, grid connectivity, and efficient administrative outreach. 

Nonetheless, deprivation persists in certain urban regions, with states like \textbf{Maharashtra}, \textbf{Bihar}, and \textbf{West Bengal} reporting modest deprivation rates (0.82\%, 0.41\%, and 0.85\%, respectively) in some urban sectors. This suggests the presence of intra-urban inequality, particularly in slums, informal settlements, and peri-urban areas. These gaps often stem from issues such as land tenure security, bureaucratic challenges, and insufficient targeted urban policies.

Therefore, there is an urgent need for \textbf{equity-focused urban policies}, addressing the unique needs of underserved slum areas, informal settlements, and marginalized urban sectors. Key strategies should include ensuring formal electricity access, expanding last-mile connectivity, and addressing tenure security in informal housing.

\subsection*{Remote Regions Leading by Example}

Smaller, geographically isolated regions such as \textbf{Lakshadweep}, \textbf{Arunachal Pradesh}, and \textbf{Ladakh} have successfully overcome geographic isolation to achieve near-universal electricity access. These regions serve as examples of how \textbf{tailored electrification policies}, including renewable energy solutions and community-based electrification models, can address access gaps even in the most remote contexts.

However, regions like \textbf{Madhya Pradesh} (0.12\%), \textbf{Uttarakhand} (0.18\%), and \textbf{Punjab} (0.19\%) still report modest deprivation rates in rural areas, indicating the need for sustained and targeted efforts to address remaining gaps in electricity access, particularly in isolated or sparsely populated communities.

\section*{ Policy Implications}

\begin{itemize}
	\item \textbf{Strengthen Rural Electrification in High-Deprivation States:}  
	States such as \textbf{Meghalaya}, \textbf{Uttar Pradesh}, and \textbf{Jharkhand}—with rural deprivation rates exceeding 1\%—should prioritize improving grid infrastructure and last-mile connectivity. Expanding investments in \textbf{solar micro-grids} and other renewable energy solutions suited to local conditions is essential to bridging the access divide.
	
	\item \textbf{Targeted Urban Policies to Address Intra-Urban Deprivation:}  
	Although urban areas like \textbf{Delhi} and \textbf{Chandigarh} have reached near-total coverage, states like \textbf{Maharashtra}, \textbf{Bihar}, and \textbf{West Bengal} still face pockets of electricity deprivation. Addressing intra-urban inequalities will require \textbf{targeted urban policies}, with a focus on slums, informal settlements, and underserved urban communities. Policies must also address tenure security and the need for formal electricity access in these marginalized areas.
	
	\item \textbf{Leverage Renewable Energy Solutions for Remote Areas:}  
	The successful electrification of remote regions like \textbf{Arunachal Pradesh}, \textbf{Lakshadweep}, and \textbf{Ladakh} highlights the potential of \textbf{decentralized renewable energy systems}. These models should be scaled up and adapted to other underserved, remote regions to ensure \textbf{universal access} to electricity, even in the most geographically challenging areas.
	
	\item \textbf{Ensure Equity in Electrification:}  
	While coverage rates have improved, \textbf{equity} remains a crucial consideration. It is imperative that the government focus on reducing electricity deprivation rates in both rural and urban sectors, ensuring that \textbf{no household}—particularly in economically disadvantaged areas—is left behind. Achieving universal electricity access should be framed not only in terms of coverage but also in terms of ensuring \textbf{equal access} for all citizens, regardless of their socio-economic status or geographic location.
\end{itemize}




\newpage

\subsection{Electricity Deprivation: Mapping State- and Sector-Level Disparities}

	
	\begin{center}
		\small
		\begin{longtable}{l l r r r}
			\caption{State-wise and Sector-wise Distribution of Households with No Lighting Arrangement} \\
			\toprule
			\textbf{State Name} & \textbf{Sector} & \textbf{Estimated Total HHs} & \textbf{Estimated No Lighting  HHs} & \textbf{No Lighting Rate (\%)} \\
			\midrule
			\endfirsthead
			
			\toprule
			\textbf{State Name} & \textbf{Sector} & \textbf{Total Weighted HHs} & \textbf{No Lighting Weighted HHs} & \textbf{No Lighting Rate (\%)} \\
			\midrule
			\endhead
			
			\bottomrule
			\endfoot
			
			Jharkhand & 1 & 5,709,801 & 76,700 & 1.34 \\
			Rajasthan & 1 & 10,936,974 & 55,760 & 0.510 \\
			Manipur & 2 & 183,691 & 770 & 0.419 \\
			Gujarat & 1 & 7,071,053 & 25,120 & 0.355 \\
			Bihar & 2 & 2,167,306 & 7,229 & 0.334 \\
			Haryana & 2 & 2,554,196 & 7,486 & 0.293 \\
			Meghalaya & 2 & 110,882 & 316 & 0.285 \\
			Odisha & 1 & 9,384,147 & 25,217 & 0.269 \\
			Tamilnadu & 2 & 10,686,204 & 20,760 & 0.194 \\
			Punjab & 1 & 3,804,113 & 7,127 & 0.187 \\
			West Bengal & 2 & 8,004,246 & 13,590 & 0.170 \\
			Uttar Pradesh & 1 & 30,398,615 & 49,911 & 0.164 \\
			Maharashtra & 1 & 14,768,694 & 22,016 & 0.149 \\
			Chhattisgarh & 1 & 4,591,188 & 6,371 & 0.139 \\
			Gujarat & 2 & 6,551,886 & 8,894 & 0.136 \\
			Uttar Pradesh & 2 & 9,075,344 & 10,266 & 0.113 \\
			Madhya Pradesh & 2 & 4,484,483 & 5,015 & 0.112 \\
			Jammu \& Kashmir & 1 & 1,799,259 & 1,656 & 0.0920 \\
			Tripura & 2 & 201,873 & 178 & 0.0884 \\
			Andhra Pradesh & 2 & 4,515,617 & 3,817 & 0.0845 \\
			Andhra Pradesh & 1 & 9,468,709 & 7,537 & 0.0796 \\
			Sikkim & 1 & 120,725 & 94.7 & 0.0784 \\
			Maharashtra & 2 & 12,562,136 & 9,126 & 0.0726 \\
			Meghalaya & 1 & 563,295 & 357 & 0.0634 \\
			Tripura & 1 & 804,896 & 492 & 0.0611 \\
			Rajasthan & 2 & 4,247,935 & 1,968 & 0.0463 \\
			Karnataka & 1 & 9,300,573 & 3,078 & 0.0331 \\
			West Bengal & 1 & 16,976,256 & 5,594 & 0.0330 \\
			Delhi & 2 & 3,643,667 & 1,096 & 0.0301 \\
			Tamilnadu & 1 & 11,176,960 & 2,596 & 0.0232 \\
			Kerala & 1 & 4,886,446 & 993 & 0.0203 \\
			Bihar & 1 & 20,151,145 & 2,574 & 0.0128 \\
			Telangana & 2 & 5,407,536 & 633 & 0.0117 \\
			Madhya Pradesh & 1 & 11,266,164 & 527 & 0.00468 \\
		
		\end{longtable}
	\end{center}
	
Despite India's remarkable achievements in expanding electricity access, a small yet significant number of households continue to experience extreme energy deprivation, as evidenced by the absence of any lighting arrangement. The classification of households with \textit{``no lighting arrangement''} represents the most severe form of electricity exclusion—indicating a complete lack of access to both conventional and alternative sources of lighting. This indicator serves as a critical benchmark for identifying the deepest pockets of energy poverty, often concentrated in geographically remote or socio-economically marginalized regions. Addressing this residual deprivation is essential for realizing the goal of universal, equitable, and inclusive electrification.


\subsection*{Persistent Deprivation in Rural Areas of Key States}

As shown in Table~\ref{tab:no_lighting_arrangement}, rural sectors (Sector 1) in several states still report considerable shares of households with no lighting arrangement. Jharkhand stands out with a rural deprivation rate of 1.34\%, followed by Rajasthan (0.51\%), Gujarat (0.36\%), and Odisha (0.27\%). These figures indicate persistent structural gaps in infrastructure access, particularly in remote hamlets, forested regions, and tribal belts.

While national programs such as Saubhagya have expanded grid connectivity, these findings suggest that grid access alone does not guarantee service delivery. Barriers such as household-level affordability constraints, poor maintenance of infrastructure, or limited awareness may continue to obstruct meaningful access.

\subsection*{Urban Deprivation: Low but Unequal}

Though overall urban deprivation is significantly lower, some urban areas still report measurable levels of electricity exclusion. For example, Bihar (urban deprivation: 0.33\%), Haryana (0.29\%), and West Bengal (0.17\%) reflect urban inequality, possibly rooted in the presence of slums, unrecognized settlements, or undocumented migrant households without formal grid access. Maharashtra, despite its industrialized status, reports a rural deprivation rate of 0.15\% and urban deprivation of 0.07\%, highlighting intra-state disparity.

The persistence of electricity deprivation in urban areas calls for more inclusive planning frameworks that target tenure insecurity, bureaucratic bottlenecks in utility provisioning, and infrastructure deficits in informal settlements.

\subsection*{Zero Deprivation Regions: Models of Success}

Several states and UTs---such as Himachal Pradesh, Chandigarh, Kerala, Sikkim, Nagaland, Arunachal Pradesh, and the Andaman \& Nicobar Islands---report zero deprivation in both rural and urban sectors. These regions demonstrate that universal electrification is attainable even in geographically challenging contexts, when supported by localized energy planning and sustained public investment.

Their success often stems from a combination of factors: decentralized renewable solutions (e.g., solar micro-grids), strong local governance, and community-based models of energy service delivery.

\subsection*{Key Policy Insights}

\begin{itemize}
	\item \textbf{Targeted Rural Investment:} States with rural deprivation rates above 0.25\%---notably Jharkhand, Rajasthan, Gujarat, and Odisha---require granular, household-level interventions. Last-mile electrification efforts should integrate renewable energy micro-solutions and leverage Gram Panchayat-level data to identify pockets of exclusion.
	
	\item \textbf{Addressing Urban Informality:} Urban electricity deprivation, though numerically smaller, often affects marginalized populations with weak formal entitlements. Policy focus must shift toward integrating slum areas into formal electricity networks, using innovative delivery mechanisms such as prepaid meters and community billing systems.
	
	\item \textbf{Replication of Success Models:} Zero-deprivation states provide replicable governance and technology models that can be adapted in high-deprivation regions. Cross-state knowledge transfers and central coordination could accelerate this process.
	
	\item \textbf{Equity as the New Metric:} Electrification success should not be measured merely by connectivity rates but by equitable and inclusive access. Households with no lighting arrangement represent the most severe form of energy poverty and must be prioritized in all future energy policies.
\end{itemize}



\newpage

\section{Evaluating the Pradhan Mantri Garib Kalyan Yojana (PMGKY): Coverage Rates and Regional Implications}

\begin{table}[htbp]
	\centering
	\resizebox{\textwidth}{!}{
		\begin{tabular}{lllll}
			\hline
			State Name & Sector & Estimated Total Households &Estimated Total PMGKY Households & PMGKY\_Coverage\_Rate \\
			\hline
			\textbf{Top 5} & & & & \\
			Nagaland      & 1  & 276944.77  & 265184.1   & 95.75 \\
			Andhra Pradesh & 1  & 9468708.6  & 8570151.7  & 90.51 \\
			Telangana     & 1  & 5689698.13 & 5080737.15 & 89.30 \\
			Chhattisgarh  & 1  & 4591188.06 & 4000220.03 & 87.13 \\
			Assam         & 1  & 6248139.01 & 4794848.88 & 76.74 \\
			\hline
			\textbf{Bottom 5} & & & & \\
			Delhi          & 1  & 89607.17   & 0          & 0.00  \\
			Chandigarh     & 2  & 242732.94  & 960.18     & 0.40  \\
			Chandigarh     & 1  & 11312.84   & 56.10      & 0.50  \\
			Delhi          & 2  & 3643666.77 & 33303.68   & 0.91  \\
			Andaman \& N. Island & 2 & 50536.00 & 583.12     & 1.15  \\
		\end{tabular}
	}
	\caption{Top 5 and Bottom 5 PMGKY Coverage Rates by State}
\end{table}



The table presents the \textbf{Top 5} and \textbf{Bottom 5} states in terms of the \textbf{PMGKY (Pradhan Mantri Garib Kalyan Yojana) Coverage Rate}, which reflects the proportion of households benefiting from the PMGKY scheme relative to the total weighted households in each state. The coverage rate is a critical indicator of the effectiveness of the scheme in reaching the intended beneficiaries, particularly in the context of welfare distribution.

\subsection*{Top 5 States by PMGKY Coverage Rate}
\begin{itemize}
	\item \textbf{Nagaland} stands out with the highest coverage rate of \textbf{95.75\%}, indicating that almost the entire targeted population in the state has received the benefits of the PMGKY scheme. This suggests a strong implementation and effective reach in this region.
	\item \textbf{Andhra Pradesh} follows closely with a coverage rate of \textbf{90.51\%}, reflecting a highly successful implementation of the scheme in this state.
	\item \textbf{Telangana} and \textbf{Chhattisgarh} also show high coverage rates of \textbf{89.30\%} and \textbf{87.13\%}, respectively, demonstrating the substantial outreach of the PMGKY program in these states.
	\item \textbf{Assam}, with a coverage rate of \textbf{76.74\%}, rounds out the top five, indicating that a significant portion of the population in Assam has benefited from the program, although there is still room for further improvements.
\end{itemize}

These states reflect successful implementation and high outreach, with most of the eligible households being included in the scheme.

\subsection*{Bottom 5 States by PMGKY Coverage Rate}
\begin{itemize}
	\item \textbf{Delhi} records the lowest coverage rate at \textbf{0.00\%}. This likely points to either incomplete or absent implementation of the PMGKY scheme for a significant portion of Delhi's population, which warrants further investigation into the barriers to coverage in the national capital.
	\item \textbf{Chandigarh} (both Sector 1 and Sector 2) follows with coverage rates of \textbf{0.50\%} and \textbf{0.40\%} respectively, suggesting very limited reach of the scheme in the region. The stark contrast between Chandigarh’s performance and other states highlights a need for targeted intervention to improve coverage in the union territory.
	\item \textbf{Andaman \& Nicobar Islands}, with a coverage rate of \textbf{1.15\%}, represents another region where PMGKY’s reach appears minimal. Although slightly higher than Delhi and Chandigarh, it still suggests a lack of widespread access.
	\item The relatively lower coverage rates in these states (Delhi, Chandigarh, Andaman \& Nicobar Islands) signal the necessity for enhancing program implementation and outreach, particularly in urban and remote regions where logistical challenges might be more prominent.
\end{itemize}

\subsection*{Policy Implications}
The stark contrast in the \textbf{PMGKY Coverage Rates} across states highlights significant regional disparities in the implementation of the scheme. The states with higher coverage rates have likely benefitted from better infrastructure, governance, and administrative efficiency, whereas those with lower rates might be facing logistical challenges, such as insufficient distribution mechanisms or barriers to beneficiary identification.

To ensure more equitable distribution of welfare benefits, it is crucial for policymakers to focus on:
\begin{itemize}
	\item \textbf{Strengthening outreach efforts} in states with lower coverage, particularly Delhi, Chandigarh, and Andaman \& Nicobar Islands.
	\item \textbf{Improving administrative mechanisms}, such as targeting vulnerable populations and overcoming barriers to scheme accessibility.
	\item \textbf{Monitoring and evaluation} to assess the effectiveness of the program, ensuring that the benefits are reaching the intended recipients across all states.
\end{itemize}




\subsection{ No Cooking Arrangement Rates by State}
%\subsection*{Top 5 States with the Highest No Cooking Arrangement Rate}



\begin{table}[h!]
	\centering
	\caption{ No Cooking Arrangement Rate for Major States}
	\centering% Reduces font size
	\resizebox{\textwidth}{!}{ % Resizes the table to the text width
		\begin{tabular}{|l|l|l|r|r|r|}
			\hline
			
			\textbf{State Name} & \textbf{Sector} & \textbf{Est Total Households} & \textbf{Est PMGKY Households} & \textbf{No Cooking Arrangement Rate(\%)} \\
			\hline
			
			\textbf{Top 5} & & & & \\
			
			Telangana            & 2  & 5407536  & 1453080 & 26.87 \\
			Odisha               & 2  & 2193468  & 390149  & 17.79 \\
			Himachal Pradesh     & 2  & 243984   & 36306   & 14.88 \\
			Karnataka            & 2  & 6880675  & 990654  & 14.4  \\
			Tamilnadu            & 2  & 10686204 & 1176535 & 11.01 \\
			\hline
		
			\textbf{Bottom 5} & & & & \\
			
			Jammu \& Kashmir     & 1  & 1799259  & 0       & 0     \\
			Chandigarh           & 1  & 11313    & 0       & 0     \\
			Haryana              & 1  & 3627657  & 0       & 0     \\
			Delhi                & 1  & 89607    & 0       & 0     \\
			Arunachal Pradesh    & 1  & 217653   & 0       & 0     \\
			\hline
		\end{tabular}
	}
\end{table}



\begin{itemize}
	\item \textbf{Telangana} (26.87\%)
	\item \textbf{Odisha} (17.79\%)
	\item \textbf{Himachal Pradesh} (14.88\%)
	\item \textbf{Karnataka} (14.4\%)
	\item \textbf{Tamil Nadu} (11.01\%)
\end{itemize}

The high No Cooking Arrangement Rates in these states may indicate a significant portion of households either lack or are not utilizing cooking arrangements, which could be reflective of underlying socio-economic conditions. These findings warrant further investigation into factors like access to cooking fuel, affordability, and the effectiveness of government schemes in these regions.

For instance, \textbf{Telangana} stands out with the highest rate of 26.87\%, suggesting a particular area of concern that may require targeted intervention to ensure households have access to basic cooking facilities. Similarly, \textbf{Odisha} and \textbf{Himachal Pradesh}, though with lower rates, show relatively high proportions of households without cooking arrangements, pointing to potential gaps in service delivery or availability of essential resources.

\subsection*{Bottom 5 States with the Lowest No Cooking Arrangement Rate}

On the other hand, the states with the lowest No Cooking Arrangement Rates, indicating relatively better access to cooking arrangements, are:

\begin{itemize}
	\item \textbf{Jammu \& Kashmir} (0\%)
	\item \textbf{Chandigarh} (0\%)
	\item \textbf{Haryana} (0\%)
	\item \textbf{Delhi} (0\%)
	\item \textbf{Arunachal Pradesh} (0\%)
\end{itemize}

Notably, these states have recorded a 0\% rate, suggesting that, at least within the context of the data analyzed, there is no significant population without access to cooking arrangements. This could point to better infrastructure, targeted policy interventions, or more favorable socio-economic conditions in these regions. However, it is essential to verify if the zero rates are a reflection of accurate reporting or potential underreporting. Further investigation into local data collection methods and government schemes may be necessary to confirm these findings.




\subsection{State-wise Rural Dependence on Firewood and Chips as Primary Cooking Energy Source in India}

\begin{table}[h!]
	\centering
	\caption{State-wise Rural Usage of Firewood and Chips as Primary Cooking Energy Source}
	
		\begin{tabular}{|l|c|r|r|r|l|}
			\hline
			\textbf{State Name} & \textbf{Sector} & \textbf{Estimated  Households } & \textbf{Estimated Total  Households} & \textbf{Percentage Share} \\
			\hline
			Chhattisgarh        & 1 & 3,814,579   & 4,591,188   & 83.1   \\
			Meghalaya           & 1 &   438,670   &   563,295   & 77.9   \\
			West Bengal         & 1 &13,108,669   &16,976,256   & 77.2   \\
			Rajasthan           & 1 & 8,290,335   &10,936,974   & 75.8   \\
			Odisha              & 1 & 7,105,336   & 9,384,147   & 75.7   \\
			Jharkhand           & 1 & 4,317,514   & 5,709,801   & 75.6   \\
			Tripura             & 1 &   578,654   &   804,896   & 71.9  \\
			Madhya Pradesh      & 1 & 7,964,373   &11,266,164   & 70.7   \\
			Assam               & 1 & 3,977,712   & 6,248,139   & 63.7   \\
			Nagaland            & 1 &   175,068   &   276,945   & 63.2  \\
			Himachal Pradesh    & 1 &   918,412   & 1,557,010   & 59.0   \\
			Uttar Pradesh       & 1 &14,802,041   &30,398,615   & 48.7   \\
			\hline
		\end{tabular}
\end{table}

A substantial proportion of rural households in several Indian states continue to rely on traditional biomass specifically firewood and chips as their primary source of cooking energy. States such as \textbf{Chhattisgarh (83.1\%)}, \textbf{Meghalaya (77.9\%)}, \textbf{West Bengal (77.2\%)}, and \textbf{Rajasthan (75.8\%)} report some of the highest usage rates, with over three-fourths of households in these regions dependent on firewood-based cooking.

Such reliance suggests enduring challenges in clean energy accessibility, despite the roll-out of national initiatives like the \textbf{\textit{Pradhan Mantri Ujjwala Yojana (PMUY)}}. Particularly in central and eastern India, the data points to regional disparities in the availability and affordability of LPG connections, as well as infrastructural limitations in last-mile delivery.

The use of firewood and chips is not only an indicator of energy poverty but also carries significant implications for public health, environmental sustainability, and gender equity. Exposure to indoor air pollution from biomass fuels increases the burden of respiratory diseases, while the time spent collecting firewood disproportionately affects women and children, limiting opportunities for education and employment.

These findings underscore the need for targeted, region-specific policy responses. State-level interventions such as enhanced LPG subsidy mechanisms, community-based awareness campaigns, and decentralized energy solutions could serve as effective levers in transitioning rural households toward cleaner cooking alternatives.




\subsection{State-wise Estimation of LPG Coverage: Top and Bottom Performers in Household Access}
\begin{table}[htbp]
	\centering
	\caption{Top and Bottom 5 States by Estimated LPG Coverage Rates}
	\resizebox{\textwidth}{!}{
		\begin{tabular}{|l|c|r|r|r|}
			\hline
			\textbf{State Name} & \textbf{Sector} & \textbf{Total Households} & \textbf{LPG Households} & \textbf{Coverage (\%)} \\
			\hline
			\multicolumn{5}{|c|}{\textbf{Top 5 States}} \\
			\hline
			Chandigarh           & 1 & 11312.84    & 11312.84    & 100.00 \\
			Mizoram              & 2 & 102745.61   & 102745.61   & 100.00 \\
			Ladakh               & 2 & 9736.46     & 9736.46     & 100.00 \\
			Ladakh               & 1 & 43307.49    & 43519.41    & 99.51 \\
			Goa                  & 2 & 222301.59   & 223872.07   & 99.30 \\
			\hline
			\multicolumn{5}{|c|}{\textbf{Bottom 5 States}} \\
			\hline
			Jharkhand            & 1 & 979300.63   & 5709801.22  & 17.15 \\
			Odisha               & 1 & 1580619.08  & 9384147.19  & 16.84 \\
			Chhattisgarh         & 1 & 701863.28   & 4591188.06  & 15.29 \\
			West Bengal          & 1 & 3477749.54  & 16976255.74 & 20.49 \\
			Rajasthan            & 1 & 2505953.94  & 10936973.91 & 22.91 \\
			\hline
		\end{tabular}
	}
\end{table}



The table presents the \textbf{Estimated LPG Coverage Rates} for the top and bottom five states in India, highlighting the proportion of households with access to liquefied petroleum gas (LPG) across different sectors (rural and urban). The \textbf{top five states}, namely \textbf{Chandigarh, Mizoram, Ladakh, Goa, and Delhi}, exhibit near-perfect LPG coverage rates, with all but Ladakh achieving 100\% LPG coverage. This suggests that these regions have effectively implemented policies or programs that ensure universal access to LPG, aligning with national objectives to promote cleaner cooking fuels.

Conversely, the \textbf{bottom five states}—\textbf{Jharkhand, Odisha, Chhattisgarh, West Bengal, and Rajasthan}—display significantly lower LPG coverage, ranging from 15.29\% to 22.91\%. This disparity underscores potential challenges in these regions, such as infrastructure limitations, affordability issues, or insufficient program outreach. These states may require targeted interventions to improve LPG penetration, especially in rural areas, in line with the \textbf{Pradhan Mantri Ujjwala Yojana} and other government initiatives aimed at increasing access to clean cooking fuel.

By highlighting the extremes in coverage, this analysis emphasizes the need for region-specific policy measures that focus on accessibility, affordability, and sustainability of LPG distribution systems. Further research could explore the socioeconomic and demographic factors that contribute to these disparities and propose strategies for enhancing coverage in lagging states.

\vspace{0.3cm}

\begin{itemize}
	\item \textbf{Top performers}: 100\% LPG coverage indicates successful implementation and access.
	\item \textbf{Bottom performers}: Highlight gaps in infrastructure and access that need attention.
	\item \textbf{Policy Implication}: Need for tailored interventions in lagging states.
	\item \textbf{National Impact}: Aligns with government efforts to provide cleaner cooking fuels for all.
\end{itemize}

Provides a critical insight into India's progress in ensuring equitable access to LPG, relevant for understanding national policy and the socio-economic implications of clean energy access.



\newpage
\subsection{Analysis of Ration Card Distribution and Coverage Gaps Sector and State-Level Examination}
The data from HCES 2023–24 reveals a critical blind spot in the coverage of India’s food security infrastructure. An analysis of the percentage of households not possessing any form of ration card — including Antyodaya Anna Yojana (AAY), Below Poverty Line (BPL), Priority Households (PHH), State Food Security Scheme (SFSS), and others — highlights states and sectors where systemic exclusion from the Public Distribution System (PDS) remains alarmingly high.

The top five state–sector combinations with the highest percentage of uncarded households are

\begin{table}[htbp]
	\centering
	\caption{Top 5 States by Proportion of Households Without a Health Card}
	\begin{tabular}{llr}
		\toprule
		\textbf{State Name} & \textbf{Sector} & \textbf{(\%) Households Without Ration Card
			} \\
		\midrule
		Chandigarh & Rural & 81.1\% \\
		Chandigarh & Urban & 79.8\% \\
		Delhi      & Urban & 74.2\% \\
		Meghalaya  & Urban & 66.4\% \\
		Delhi      & Rural & 65.0\% \\
		\bottomrule
	\end{tabular}
\end{table}

Chandigarh emerges as the most affected region, indicating near-total exclusion from the Public Distribution System (PDS), irrespective of settlement type. Key figures include
\begin{itemize}
	\item 81.1\% of rural households without any ration card
	\item 79.8\% of urban households without any ration card
\end{itemize}
These alarming statistics highlight serious lapses in identification, outreach, and enrolment processes.

In Delhi, the situation is comparably dire, despite being the National Capital Territory with expectedly robust administrative mechanisms
\begin{itemize}
	\item 74.2\% of urban households report having no ration card
	\item 65.0\% of rural households report having no ration card
\end{itemize}
This systemic undercoverage may be attributed to factors such as large-scale in-migration, urban informality, and documentation gaps among slum-dwellers and tenants.

Meghalaya’s urban sector also presents significant concerns, with 66.4\% of households lacking ration cards. Potential impediments include topographical barriers, limited bureaucratic penetration, and digital connectivity constraints that hinder effective beneficiary registration and service delivery.

\noindent \textbf{Note:} A detailed table presenting the percentage of households without any ration card across all States and Union Territories, disaggregated by rural and urban sectors, is provided in the Appendix.


\subsection{State-Level Dependence on State Food Security Schemes (SFSS)}

The table presents the top five states with the highest Percentage of households dependent on State Food Security Schemes (SFSS), based on weighted estimates from the HCES 2023–24.

\begin{table}[h]
	\centering
	\caption{Top 5 States by Household Dependence on State Food Security Schemes (SFSS)}
	\begin{tabular}{lrrr}
		\toprule
		\textbf{State Name} & \textbf{Estimated SFSS Households} & \textbf{Estimated Total Households} & \textbf{SFSS Share (\%)} \\
		\midrule
		West Bengal   & 6,637,859  & 24,980,502 & 26.60\% \\
		Telangana     & 1,679,019  & 11,097,234 & 15.10\% \\
		Punjab        &   826,673  &  6,548,466 & 12.60\% \\
		Uttarakhand   &   251,691  &  2,611,389 & 9.64\% \\
		Odisha        &   417,648  & 11,577,615 & 3.61\% \\
		\bottomrule
	\end{tabular}
\end{table}

\begin{itemize}[leftmargin=1.5em]
	
	\item \textbf{West Bengal} records the highest household dependence on State Food Security Schemes (SFSS), with over \textbf{26.6\%} of households accessing rations through state-funded provisions. This substantial figure translating to more than \textbf{6.6 million households} indicates that the state government is playing a dominant role in ensuring food security, shouldering a responsibility that ideally should be more equitably distributed between central and state institutions.
	
	\item \textbf{Telangana} and \textbf{Punjab} follow with \textbf{15.1\%} and \textbf{12.6\%} coverage, respectively.These levels of dependence suggest that large sections of the population in these states may be underserved by centrally provided entitlements, compelling state governments to fill critical welfare gaps through independent schemes.
	
	\item In \textbf{Uttarakhand}, where \textbf{9.64\%} of households are enrolled in SFSS, the burden on the state may stem from geographic isolation, logistical challenges, and the need for customised distribution networks, especially in rural and hilly regions.
	
	\item \textbf{Odisha’s} relatively lower SFSS coverage of \textbf{3.61\%} still amounts to a significant absolute burden when viewed in the context of the state's broader welfare portfolio and fiscal constraints.
	
\end{itemize}





	\section{Insights and Policy Implications}
	The findings from the analysis reveal significant variations in household non-food consumption patterns across Indian states, shaped by diverse economic and social factors. Key determinants such as income disparities, regional economic development, government interventions, and access to essential services significantly influence expenditure trends.
	
	For high-performing states like Kerala and Tamil Nadu, the analysis emphasizes robust investments in education and healthcare, reflecting these states' commitment to fostering human capital and infrastructure. Their expenditure distribution indicates a focus on long-term socio-economic well-being. Conversely, states like Bihar, Odisha, and Jharkhand exhibit a larger allocation towards basic necessities such as fuel and light, which underscores the economic challenges faced by households in these regions.
	
	Nationally, the rising healthcare expenditure and variations in spending on durable goods and education highlight the dual nature of household priorities: survival-driven spending versus aspirational consumption. Developed states appear to allocate resources toward lifestyle-oriented categories, while less-developed regions prioritize basic sustenance.
	
	These insights provide a foundation for targeted policy recommendations aimed at improving economic welfare and household resilience. Addressing regional disparities through enhanced infrastructure, skill development, and inclusive growth strategies can uplift underperforming states. Moreover, sustained investment in education and healthcare across all regions can build resilience and foster equitable urban growth.
	
	By addressing these factors holistically, policymakers can balance household consumption patterns, ensuring both immediate needs and long-term aspirations are met.
	
	% Conclusion and Recommendations
	\section{Conclusion and Recommendations}
	
	\textbf{Conclusion:}  
	The analysis underscores the diverse expenditure patterns across rural and urban households, reflecting economic disparities and regional inequalities. Rural areas tend to allocate a larger share of their income to basic necessities like fuel and light, while urban households show a broader distribution, with significant spending on education, healthcare, and durable goods in more developed states. Inflation emerges as a critical factor influencing consumption behavior, especially in low-income groups where rising costs constrain spending on non-essential categories. Additionally, the role of government interventions in shaping household financial stability cannot be overstated. Well-designed policies have the potential to alleviate disparities and promote inclusive growth.
	
	\textbf{Recommendations:}  
	\begin{itemize}
		\item \textbf{Address Rural-Urban Disparities:} 
		\begin{itemize}
			\item Implement targeted welfare programs aimed at improving access to education and healthcare in rural areas.
			\item Invest in rural infrastructure to reduce the dependency on basic necessity expenditure, freeing resources for long-term investments.
		\end{itemize}
		
		\item \textbf{Mitigate the Effects of Inflation:}
		\begin{itemize}
			\item Strengthen price regulation mechanisms to protect lower-income groups from the adverse impacts of rising costs.
			\item Introduce subsidies for essential goods and services to reduce the financial burden on vulnerable households.
		\end{itemize}
		
		\item \textbf{Promote Balanced Regional Development:}
		\begin{itemize}
			\item Encourage investments in education and skill development in underperforming states like Bihar, Jharkhand, and Odisha.
			\item Foster urbanization in a sustainable manner to ensure equitable access to opportunities and services.
		\end{itemize}
		
		\item \textbf{Strengthen Government Policies:}
		\begin{itemize}
			\item Design policies that prioritize inclusive growth by addressing structural barriers to economic participation.
			\item Regularly assess and adapt policies to respond to changing economic conditions and household needs.
		\end{itemize}
	\end{itemize}
	
	By implementing these recommendations, policymakers can bridge the gap between rural and urban expenditure patterns, enhance household resilience to inflation, and create a framework for equitable and sustainable economic growth.
	
\section{Policy Recommendations}

The findings from the analysis provide valuable insights into the socio-economic dynamics shaping household expenditure patterns across India. Based on these observations, the following policy recommendations are proposed to address key challenges such as income disparities, inflation impacts, and infrastructure gaps, while promoting inclusivity and sustainable development.

\subsection{Strengthening Welfare Programs}
The disparities in household expenditure patterns, particularly between rural and urban areas, highlight the need for targeted welfare initiatives to support lower-income households. Policies can focus on:
\begin{itemize}
	\item \textbf{Direct Financial Assistance:} 
	Cash transfer schemes can ensure immediate relief to vulnerable households, enabling them to meet essential expenditure needs such as food, healthcare, and education.
	\item \textbf{Subsidized Essential Services:} 
	Providing subsidies for energy, education, and healthcare can alleviate economic burdens and enhance access to critical resources.
	\item \textbf{Job Creation Programs:} 
	Initiatives focusing on skill-building, vocational training, and entrepreneurship development can empower low-income groups to secure sustainable livelihoods.
\end{itemize}
By strengthening welfare programs, policymakers can bridge income gaps and promote financial resilience across diverse demographics.

\subsection{Monitoring Inflation Impacts}
Inflation is a critical factor influencing household consumption behavior, particularly in lower-income groups where rising costs can disproportionately affect spending on non-essential categories. Key measures include:
\begin{itemize}
	\item \textbf{Price Regulation Mechanisms:} 
	Establishing monitoring systems to track inflation trends in essential commodity markets can help implement effective price controls.
	\item \textbf{Affordable Pricing Frameworks:} 
	Regulatory measures ensuring stable prices for food, fuel, and medical supplies can safeguard vulnerable households against economic shocks.
	\item \textbf{Public Awareness Campaigns:} 
	Educating citizens about inflation-related developments and providing guidance on budgeting and financial planning can enhance household resilience.
\end{itemize}
By addressing the impacts of inflation, policymakers can create a more stable economic environment conducive to balanced consumption patterns.

\subsection{Enhancing Digital Payment Infrastructure}
Digital payment systems have become integral to modern economies, offering convenience, transparency, and security in financial transactions. Expanding this infrastructure can promote financial inclusion and efficiency. Recommended actions include:
\begin{itemize}
	\item \textbf{Expanding Access to Digital Platforms:} 
	Investment in digital payment systems across urban and rural areas can ensure ease of access for all households, reducing dependency on cash transactions.
	\item \textbf{Financial Literacy Programs:} 
	Providing educational resources on using digital payment platforms can empower households to transition to a cashless economy.
	\item \textbf{Collaborations with Financial Institutions:} 
	Partnering with banks and technology providers to improve the reliability, security, and user experience of digital payment systems can accelerate adoption rates.
\end{itemize}
Enhanced digital infrastructure can simplify financial transactions, reduce transaction costs, and foster greater economic participation.

\subsection{Addressing Regional Disparities}
The stark differences in expenditure patterns between developed and underdeveloped states necessitate targeted interventions to promote balanced regional growth. Policymakers can focus on:
\begin{itemize}
	\item \textbf{Investments in Infrastructure:} 
	Enhancing rural and urban infrastructure in underperforming states like Bihar and Odisha can improve access to education, healthcare, and employment opportunities.
	\item \textbf{Skill Development and Education:} 
	Introducing skill-building programs and expanding educational facilities in economically disadvantaged regions can foster upward mobility.
	\item \textbf{Promoting Entrepreneurship:} 
	Creating a supportive ecosystem for micro-enterprises and small businesses can boost local economies and reduce dependency on external aid.
\end{itemize}
Addressing regional disparities through these initiatives can bridge the gap between high-performing and underperforming states.

\subsection{Fostering Sustainable Urban Development}
Urban areas often exhibit higher expenditure levels in categories such as healthcare, education, and lifestyle-driven goods. To ensure this growth is inclusive and sustainable, policymakers can implement:
\begin{itemize}
	\item \textbf{Affordable Housing Programs:} 
	Expanding access to affordable housing can mitigate the cost of living in urban centers and improve financial stability for low-income households.
	\item \textbf{Enhanced Public Services:} 
	Investments in urban infrastructure, including public transportation, healthcare facilities, and educational institutions, can support balanced urban growth.
	\item \textbf{Smart Urban Planning:} 
	Leveraging technology and data-driven approaches to plan urban development can ensure equitable access to resources and services.
\end{itemize}
Sustainable urban development strategies can enhance the quality of life for all urban residents while minimizing environmental and social risks.

\subsection{Strengthening Government Policies}
Effective government policies play a pivotal role in shaping household financial stability and consumption patterns. To maximize their impact, policymakers should consider:
\begin{itemize}
	\item \textbf{Inclusive Growth Strategies:} 
	Designing policies that address structural barriers to economic participation, such as discrimination and unequal access to resources, can foster inclusivity.
	\item \textbf{Periodic Policy Assessments:} 
	Regular evaluations of policy effectiveness can help adapt interventions to evolving economic conditions and household needs.
	\item \textbf{Collaborative Policy Design:} 
	Engaging stakeholders, including local governments, NGOs, and community representatives, in policy formulation can ensure initiatives are aligned with ground realities.
\end{itemize}
Strong government policies can lay the foundation for long-term economic stability and social well-being.
\newpage

\section{Appendix}

\subsection{Estimated Total Households PMGKY Coverage Rates by State}
\begin{table}[h!]
	\centering
	\caption{Estimated Total Households PMGKY Coverage Rates by State}
	\centering% Reduces font size
	\resizebox{\textwidth}{!}{ % Resizes the table to the text width
		\begin{tabular}{|l|l|l|r|r|r|}
			\hline
			\textbf{State Name} & \textbf{Sector} & \textbf{Estimated Total Households} & \textbf{Estimated PMGKY  Households} & \textbf{PMGKY Coverage Rate} \\
			\hline
			Delhi                         & 1  & 89607.17   & 0        & 0.00  \\
			Chandigarh                    & 2  & 242732.94  & 960.18   & 0.40  \\
			Chandigarh                    & 1  & 11312.84   & 56.10    & 0.50  \\
			Delhi                         & 2  & 3643666.77 & 33303.68 & 0.91  \\
			Andaman \& N. Island          & 2  & 50536.00   & 583.12   & 1.15  \\
			Himachal Pradesh              & 2  & 243984.11  & 6504.44  & 2.67  \\
			Mizoram                       & 2  & 102745.61  & 4237.84  & 4.12  \\
			Mizoram                       & 1  & 112840.32  & 5784.21  & 5.13  \\
			Himachal Pradesh              & 1  & 1557009.68 & 163361.80 & 10.49 \\
			Punjab                        & 2  & 2744353.4  & 304035.77 & 11.08 \\
			Andaman \& N. Island          & 1  & 76437.63   & 8504.69  & 11.13 \\
			Meghalaya                     & 2  & 110882.15  & 13117.21 & 11.83 \\
			Jammu \& Kashmir              & 2  & 566667.93  & 68092.14 & 12.02 \\
			D \& N. Haveli \& Daman \& Diu & 2  & 131872.64  & 20967.97 & 15.90 \\
			Haryana                       & 2  & 2554195.86 & 415768.63 & 16.28 \\
			Ladakh                        & 2  & 9736.46    & 1728.85  & 17.76 \\
			Maharashtra                   & 2  & 12562136.12 & 2377529.01 & 18.93 \\
			Gujarat                        & 2  & 6551886.4  & 1288341.34 & 19.66 \\
			Arunachal Pradesh             & 2  & 62305.06   & 14115.17  & 22.65 \\
			Sikkim                        & 2  & 57497.87   & 13695.65  & 23.82 \\
			Uttarakhand                   & 2  & 754954.96  & 182106.8  & 24.12 \\
			Punjab                        & 1  & 3804112.58 & 943878.26 & 24.81 \\
			Rajasthan                     & 2  & 4247934.51 & 1119608.87 & 26.36 \\
			Ladakh                        & 1  & 43519.41   & 11600.28  & 26.66 \\
			Goa                           & 2  & 223872.07  & 63983.15  & 28.58 \\
			Puduchery                     & 2  & 240644.51  & 71882.81  & 29.87 \\
			Madhya Pradesh                & 2  & 4484482.51 & 1427289.05 & 31.83 \\
			Jharkhand                     & 2  & 1615379.51 & 517381.91 & 32.03 \\
			Assam                         & 2  & 981145.46  & 314833.53 & 32.09 \\
			Jammu \& Kashmir              & 1  & 1799259.39 & 603183.97 & 33.52 \\
			Karnataka                     & 2  & 6880674.79 & 2338623.25 & 33.99 \\
			Manipur                       & 2  & 183690.58  & 63298.12  & 34.46 \\
			Haryana                       & 1  & 3627657.33 & 1251326.07 & 34.49 \\
			Kerala                        & 2  & 4925091.17 & 1770124.57 & 35.94 \\
			Tamilnadu                     & 2  & 10686203.56 & 3865701.57 & 36.17 \\
			Odisha                        & 2  & 2193467.79 & 798064.55 & 36.38 \\
			Lakshadweep                   & 1  & 3016.77    & 1139.68   & 37.78 \\
			West Bengal                   & 2  & 8004246.47 & 3036925.05 & 37.94 \\
			Lakshadweep                   & 2  & 9343.77    & 3823.12   & 40.92 \\
			Puduchery                     & 1  & 106244.13  & 43590.3   & 41.03 \\
			Uttar Pradesh                 & 2  & 9075343.67 & 3908988.67 & 43.07 \\
			Tamilnadu                     & 1  & 11176960.22 & 4857665.36 & 43.46 \\
			Kerala                        & 1  & 4886446.27 & 2198356.7  & 44.99 \\
			Bihar                         & 2  & 2167306.01 & 986182.09  & 45.50 \\
			Manipur                       & 1  & 499367.07  & 228367.02  & 45.73 \\
			Meghalaya                     & 1  & 563295.34  & 263967.62  & 46.86 \\
			Tripura                       & 2  & 201872.51  & 98071.52   & 48.58 \\
			Telangana                     & 2  & 5407535.9  & 2722447.07 & 50.35 \\
			Uttarakhand                   & 1  & 1856433.92 & 998905.64  & 53.81 \\
			Goa                           & 1  & 134511.13  & 72384.36   & 53.81 \\
			Chhattisgarh                  & 2  & 1332889.98 & 776042.61  & 58.22 \\
			Arunachal Pradesh             & 1  & 217652.92  & 129196.71  & 59.36 \\
			Gujarat                        & 1  & 7071052.71 & 4333705.83 & 61.29 \\
			Madhya Pradesh                & 1  & 11266164.09 & 7111959.84 & 63.13 \\
			Maharashtra                   & 1  & 14768693.59 & 9347707.35 & 63.29 \\
			West Bengal                   & 1  & 16976255.74 & 10866517.28 & 64.01 \\
			Odisha                        & 1  & 9384147.19 & 6047324.5  & 64.44 \\
			Rajasthan                     & 1  & 10936973.91 & 7277146.14 & 66.54 \\
			Jharkhand                     & 1  & 5709801.22 & 3827899.12 & 67.04 \\
			Sikkim                        & 1  & 120725.43   & 82332.13   & 68.20 \\
			D \& N. Haveli \& Daman \& Diu & 1  & 54438.21   & 37310.99   & 68.54 \\
			Tripura                       & 1  & 804895.95  & 565920.69  & 70.31 \\
			Karnataka                     & 1  & 9300572.68 & 6612606.94 & 71.10 \\
			Bihar                         & 1  & 20151145.33 & 14419793.53 & 71.56 \\
			Andhra Pradesh                & 2  & 4515617.41 & 3243661.63 & 71.83 \\
			Uttar Pradesh                 & 1  & 30398614.51 & 22013180.67 & 72.42 \\
			Nagaland                      & 2  & 120338.14   & 87964.07   & 73.10 \\
			Assam                         & 1  & 6248139.01  & 4794848.88 & 76.74 \\
			Chhattisgarh                  & 1  & 4591188.06  & 4000220.03 & 87.13 \\
			Telangana                     & 1  & 5689698.13  & 5080737.15 & 89.30 \\
			Andhra Pradesh                & 1  & 9468708.6   & 8570151.7  & 90.51 \\
			Nagaland                      & 1  & 276944.77   & 265184.1   & 95.75 \\
			\hline
		\end{tabular}
	}
\end{table}
\subsection{Estimated Total Households No Cooking Arrangement Rate by State}
\begin{table}[h!]
\centering
\caption{Estimated Total Households No Cooking Arrangement Rate by State}
\centering% Reduces font size
\resizebox{\textwidth}{!}{ % Resizes the table to the text width
	\begin{tabular}{|l|l|l|r|r|r|}
		\hline
		\textbf{State Name} & \textbf{Sector} & \textbf{Est Total Households} & \textbf{Est PMGKY  Households} & \textbf{No Cooking Arrangement Rate} \\
		\hline
		Telangana                      & 2 & 5407536  & 1453080 & 26.87 \\
		Odisha                         & 2 & 2193468  & 390149  & 17.79 \\
		Himachal Pradesh               & 2 & 243984   & 36306   & 14.88 \\
		Karnataka                      & 2 & 6880675  & 990654  & 14.4  \\
		Tamilnadu                      & 2 & 10686204 & 1176535 & 11.01 \\
		Andaman \& N. Island           & 1 & 76438    & 6238    & 8.16  \\
		Andhra Pradesh                 & 2 & 4515617  & 364725  & 8.08  \\
		Puduchery                      & 2 & 240645   & 17957   & 7.46  \\
		Sikkim                         & 2 & 57498    & 4113    & 7.15  \\
		Lakshadweep                    & 2 & 9344     & 636     & 6.81  \\
		Odisha                         & 1 & 9384147  & 599548  & 6.39  \\
		D \& N. Haveli \& Daman \& Diu & 2 & 131873   & 7727    & 5.86  \\
		Rajasthan                      & 2 & 4247935  & 244859  & 5.76  \\
		Kerala                         & 2 & 4925091  & 274733  & 5.58  \\
		Maharashtra                    & 2 & 12562136 & 642970  & 5.12  \\
		Gujarat                        & 2 & 6551886  & 322761  & 4.93  \\
		Chhattisgarh                   & 2 & 1332890  & 54675   & 4.1   \\
		Lakshadweep                    & 1 & 3017     & 108     & 3.58  \\
		D \& N. Haveli \& Daman \& Diu & 1 & 54438    & 1725    & 3.17  \\
		Andhra Pradesh                 & 1 & 9468709  & 294600  & 3.11  \\
		West Bengal                    & 2 & 8004246  & 239068  & 2.99  \\
		Jharkhand                      & 2 & 1615380  & 44436   & 2.75  \\
		Telangana                      & 1 & 5689698  & 145289  & 2.55  \\
		Meghalaya                      & 2 & 110882   & 2417    & 2.18  \\
		Delhi                          & 2 & 3643667  & 73018   & 2     \\
		Manipur                        & 2 & 183691   & 3614    & 1.97  \\
		Haryana                        & 2 & 2554196  & 49370   & 1.93  \\
		Tamilnadu                      & 1 & 11176960 & 205642  & 1.84  \\
		Kerala                         & 1 & 4886446  & 88277   & 1.81  \\
		Maharashtra                    & 1 & 14768694 & 238181  & 1.61  \\
		Chandigarh                     & 2 & 242733   & 3879    & 1.6   \\
		Uttarakhand                    & 2 & 754955   & 11608   & 1.54  \\
		Jammu \& Kashmir               & 2 & 566668   & 8251    & 1.46  \\
		Uttar Pradesh                  & 2 & 9075344  & 131205  & 1.45  \\
		Karnataka                      & 1 & 9300573  & 134053  & 1.44  \\
		Assam                          & 2 & 981145   & 12011   & 1.22  \\
		Chhattisgarh                   & 1 & 4591188  & 55232   & 1.2   \\
		Punjab                         & 2 & 2744353  & 32355   & 1.18  \\
		Bihar                          & 2 & 2167306  & 23368   & 1.08  \\
		Madhya Pradesh                 & 2 & 4484483  & 47335   & 1.06  \\
		Puduchery                      & 1 & 106244   & 1079    & 1.02  \\
		Uttarakhand                    & 1 & 1856434  & 15881   & 0.86  \\
		Rajasthan                      & 1 & 10936974 & 78872   & 0.72  \\
		Himachal Pradesh               & 1 & 1557010  & 10890   & 0.7   \\
		Jharkhand                      & 1 & 5709801  & 32299   & 0.57  \\
		West Bengal                    & 1 & 16976256 & 95077   & 0.56  \\
		Gujarat                        & 1 & 7071053  & 39115   & 0.55  \\
		Goa                            & 2 & 223872   & 1178    & 0.53  \\
		Punjab                         & 1 & 3804113  & 19087   & 0.5   \\
		Sikkim                         & 1 & 120725   & 350     & 0.29  \\
		Arunachal Pradesh              & 2 & 62305    & 171     & 0.27  \\
		Uttar Pradesh                  & 1 & 30398615 & 39317   & 0.13  \\
		Goa                            & 1 & 134511   & 179     & 0.13  \\
		Madhya Pradesh                 & 1 & 11266164 & 3439    & 0.03  \\
		Assam                          & 1 & 6248139  & 1213    & 0.02  \\
		Bihar                          & 1 & 20151145 & 2564    & 0.01  \\
		Jammu \& Kashmir               & 1 & 1799259  & 0       & 0     \\
		Chandigarh                     & 1 & 11313    & 0       & 0     \\
		Haryana                        & 1 & 3627657  & 0       & 0     \\
		Delhi                          & 1 & 89607    & 0       & 0     \\
		Arunachal Pradesh              & 1 & 217653   & 0       & 0     \\
		Nagaland                       & 1 & 276945   & 0       & 0     \\
		Nagaland                       & 2 & 120338   & 0       & 0     \\
		Manipur                        & 1 & 499367   & 0       & 0     \\
		Mizoram                        & 1 & 112840   & 0       & 0     \\
		Mizoram                        & 2 & 102746   & 0       & 0     \\
		Tripura                        & 1 & 804896   & 0       & 0     \\
		Tripura                        & 2 & 201873   & 0       & 0     \\
		Meghalaya                      & 1 & 563295   & 0       & 0     \\
		Andaman \& N. Island           & 2 & 50536    & 0       & 0     \\
		Ladakh                         & 1 & 43519    & 0       & 0     \\
		Ladakh                         & 2 & 9736     & 0       & 0    \\
		\hline
	\end{tabular}
}
\end{table}

\subsection{Estimated Total Households LPG Coverage Rates by State}
\begin{table}[h!]
	\centering
	\caption{Estimated Total Households LPG Coverage Rates by State}
	\centering% Reduces font size
	\resizebox{0.9\textwidth}{!}{ % Resizes the table to the text width
		\begin{tabular}{|l|l|l|r|r|r|}
			\hline
		State\_Name                    & Sector & Estimated total Households & Estimated LPG Households & Percentage\_Share \\
		\hline
		Chandigarh           & 1 & 11312.84    & 11312.84    & 100   \\
		Mizoram              & 2 & 102745.61   & 102745.61   & 100   \\
		Ladakh               & 2 & 9736.46     & 9736.46     & 100   \\
		Ladakh               & 1 & 43307.49    & 43519.41    & 99.51 \\
		Goa                  & 2 & 222301.59   & 223872.07   & 99.3  \\
		Delhi                & 1 & 88914.49    & 89607.17    & 99.23 \\
		Chandigarh           & 2 & 237689.28   & 242732.94   & 97.92 \\
		Andaman \& N. Island & 2 & 49321.68    & 50536       & 97.6  \\
		Goa                  & 1 & 130719.98   & 134511.13   & 97.18 \\
		Punjab               & 2 & 2664312.6   & 2744353.4   & 97.08 \\
		Arunachal Pradesh    & 2 & 60360.17    & 62305.06    & 96.88 \\
		Manipur              & 2 & 177480.82   & 183690.58   & 96.62 \\
		Telangana            & 1 & 5421699.89  & 5689698.13  & 95.29 \\
		Haryana              & 2 & 2419726.68  & 2554195.86  & 94.74 \\
		Uttarakhand          & 2 & 713697.93   & 754954.96   & 94.54 \\
		Assam                & 2 & 925233.65   & 981145.46   & 94.3  \\
		Jammu \& Kashmir     & 2 & 532860.7    & 566667.93   & 94.03 \\
		Nagaland             & 2 & 112028.46   & 120338.14   & 93.09 \\
		Mizoram              & 1 & 105018.27   & 112840.32   & 93.07 \\
		Sikkim               & 2 & 53385.21    & 57497.87    & 92.85 \\
		Puduchery            & 2 & 221220.4    & 240644.51   & 91.93 \\
		Puduchery            & 1 & 96544.03    & 106244.13   & 90.87 \\
		Madhya Pradesh       & 2 & 4071363.87  & 4484482.51  & 90.79 \\
		Uttar Pradesh        & 2 & 8204801.52  & 9075343.67  & 90.41 \\
		Andhra Pradesh       & 2 & 4050122.78  & 4515617.41  & 89.69 \\
		D \& N. Haveli \& Daman \& Diu & 2      & 114753.56            & 131872.64                   & 87.02             \\
		Andhra Pradesh       & 1 & 8228560.55  & 9468708.6   & 86.9  \\
		Bihar                & 2 & 1868683.16  & 2167306.01  & 86.22 \\
		Tamilnadu            & 2 & 9172547.74  & 10686203.56 & 85.84 \\
		Rajasthan            & 2 & 3636890.5   & 4247934.51  & 85.62 \\
		Maharashtra          & 2 & 10749947.44 & 12562136.12 & 85.57 \\
		Delhi                & 2 & 3097320.29  & 3643666.77  & 85.01 \\
		Punjab               & 1 & 3199396.36  & 3804112.58  & 84.1  \\
		Karnataka            & 2 & 5785882.63  & 6880674.79  & 84.09 \\
		Andaman \& N. Island & 1 & 64151.39    & 76437.63    & 83.93 \\
		West Bengal          & 2 & 6638326.08  & 8004246.47  & 82.94 \\
		Himachal Pradesh     & 2 & 200040.96   & 243984.11   & 81.99 \\
		Karnataka            & 1 & 7581340.27  & 9300572.68  & 81.51 \\
		Tamilnadu            & 1 & 9096379.58  & 11176960.22 & 81.39 \\
		Chhattisgarh         & 2 & 1083113.67  & 1332889.98  & 81.26 \\
		Sikkim               & 1 & 96583.73    & 120725.43   & 80    \\
		Meghalaya            & 2 & 86748.69    & 110882.15   & 78.24 \\
		Manipur              & 1 & 388489.92   & 499367.07   & 77.8  \\
		Jharkhand            & 2 & 1252909.7   & 1615379.51  & 77.56 \\
		Lakshadweep          & 2 & 7155.13     & 9343.77     & 76.58 \\
		Kerala               & 2 & 3647419.71  & 4925091.17  & 74.06 \\
		Telangana            & 2 & 3945185.09  & 5407535.9   & 72.96 \\
		Tripura              & 2 & 139581.87   & 201872.51   & 69.14 \\
		Maharashtra          & 1 & 10146540.03 & 14768693.59 & 68.7  \\
		Jammu \& Kashmir     & 1 & 1188112.32  & 1799259.39  & 66.03 \\
		Gujarat              & 2 & 4289477.94  & 6551886.4   & 65.47 \\
		Odisha               & 2 & 1392734.45  & 2193467.79  & 63.49 \\
		Gujarat              & 1 & 4327019.73  & 7071052.71  & 61.19 \\
		Haryana              & 1 & 2156350.8   & 3627657.33  & 59.44 \\
		Kerala               & 1 & 2856582.56  & 4886446.27  & 58.46 \\
		Arunachal Pradesh    & 1 & 126248.95   & 217652.92   & 58    \\
		Uttarakhand          & 1 & 1065581.84  & 1856433.92  & 57.4  \\
		Lakshadweep          & 1 & 1705.35     & 3016.77     & 56.53 \\
		D \& N. Haveli \& Daman \& Diu & 1      & 27384                & 54438.21                    & 50.3              \\
		Uttar Pradesh        & 1 & 13504784.46 & 30398614.51 & 44.43 \\
		Himachal Pradesh     & 1 & 621849.45   & 1557009.68  & 39.94 \\
		Bihar                & 1 & 7595293.67  & 20151145.33 & 37.69 \\
		Nagaland             & 1 & 101876.52   & 276944.77   & 36.79 \\
		Assam                & 1 & 2269213.79  & 6248139.01  & 36.32 \\
		Madhya Pradesh       & 1 & 3204506.69  & 11266164.09 & 28.44 \\
		Tripura              & 1 & 226241.59   & 804895.95   & 28.11 \\
		Rajasthan            & 1 & 2505953.94  & 10936973.91 & 22.91 \\
		West Bengal          & 1 & 3477749.54  & 16976255.74 & 20.49 \\
		Meghalaya            & 1 & 110745.93   & 563295.34   & 19.66 \\
		Jharkhand            & 1 & 979300.63   & 5709801.22  & 17.15 \\
		Odisha               & 1 & 1580619.08  & 9384147.19  & 16.84 \\
		Chhattisgarh         & 1 & 701863.28   & 4591188.06  & 15.29\\
		\hline
	\end{tabular}
}
\end{table}



\begin{longtable}{llr}
	\caption{Percentage of Households Without Any type  Ration Card by State and Sector} \\
	\toprule
	\textbf{State Name} & \textbf{Sector} & \textbf{Percentage (\%)} \\
	\midrule
	\endfirsthead
	
	\toprule
	\textbf{State Name} & \textbf{Sector} & \textbf{Precentage (\%)} \\
	\midrule
	\endhead
	
	\bottomrule
	\endfoot
	Chandigarh & 1 & 81.1 \\
	Chandigarh & 2 & 79.8 \\
	Delhi & 2 & 74.2 \\
	Meghalaya & 2 & 66.4 \\
	Delhi & 1 & 65.0 \\
	Sikkim & 2 & 59.1 \\
	Punjab & 2 & 55.5 \\
	Assam & 2 & 55.4 \\
	Odisha & 2 & 53.2 \\
	D \& N. Haveli \& Daman \& Diu & 2 & 50.4 \\
	Madhya Pradesh & 2 & 50.3 \\
	Jharkhand & 2 & 49.6 \\
	Arunachal Pradesh & 2 & 48.7 \\
	Haryana & 2 & 47.5 \\
	Uttarakhand & 2 & 47.2 \\
	Bihar & 2 & 47.2 \\
	Uttar Pradesh & 2 & 45.9 \\
	Himachal Pradesh & 2 & 45.8 \\
	Karnataka & 2 & 45.7 \\
	Telangana & 2 & 39.8 \\
	Maharashtra & 2 & 32.7 \\
	Goa & 2 & 28.6 \\
	Meghalaya & 1 & 27.8 \\
	Manipur & 2 & 27.2 \\
	Punjab & 1 & 26.7 \\
	Nagaland & 2 & 23.2 \\
	Gujarat & 2 & 23.1 \\
	Manipur & 1 & 22.3 \\
	Andhra Pradesh & 2 & 21.9 \\
	Bihar & 1 & 21.8 \\
	Madhya Pradesh & 1 & 19.5 \\
	Chhattisgarh & 2 & 19.0 \\
	Tamilnadu & 2 & 18.8 \\
	Arunachal Pradesh & 1 & 18.7 \\
	Uttar Pradesh & 1 & 17.0 \\
	D \& N. Haveli \& Daman \& Diu & 1 & 16.4 \\
	Haryana & 1 & 15.8 \\
	Rajasthan & 2 & 15.4 \\
	Odisha & 1 & 15.2 \\
	West Bengal & 2 & 14.5 \\
	Goa & 1 & 13.9 \\
	Uttarakhand & 1 & 13.1 \\
	Assam & 1 & 12.1 \\
	Jharkhand & 1 & 10.5 \\
	Maharashtra & 1 & 10.3 \\
	Sikkim & 1 & 10.2 \\
	Ladakh & 2 & 8.86 \\
	Lakshadweep & 2 & 8.36 \\
	Karnataka & 1 & 7.4 \\
	Jammu \& Kashmir & 2 & 7.12 \\
	Kerala & 2 & 6.61 \\
	Tamilnadu & 1 & 5.58 \\
	Telangana & 1 & 5.07 \\
	Himachal Pradesh & 1 & 4.32 \\
	Kerala & 1 & 3.46 \\
	Andhra Pradesh & 1 & 3.37 \\
	Andaman \& N. Island & 2 & 3.04 \\
	Puduchery & 2 & 2.95 \\
	Gujarat & 1 & 2.91 \\
	Nagaland & 1 & 2.87 \\
	Puduchery & 1 & 2.60 \\
	Tripura & 1 & 2.28 \\
	Rajasthan & 1 & 2.23 \\
	Andaman \& N. Island & 1 & 2.19 \\
	Chhattisgarh & 1 & 2.11 \\
	Tripura & 2 & 1.78 \\
	Lakshadweep & 1 & 1.69 \\
	Mizoram & 2 & 1.64 \\
	Jammu \& Kashmir & 1 & 1.61 \\
	Mizoram & 1 & 1.11 \\
	West Bengal & 1 & 0.65 \\
	Ladakh & 1 & 0.32 \\
	
\end{longtable}
\newpage
\section{References}

\begin{itemize}
	\item \textbf{Household Consumer Expenditure Survey (HCES):} Ministry of Statistics and Programme Implementation (MOSPI). Official page available at: \url{https://mospi.gov.in/}
	\item \textbf{HCES Unit-Level Data:} Access and download the unit-level data for the Household Consumer Expenditure Survey (HCES) provided by MOSPI at: \url{https://microdata.gov.in/NADA/index.php/catalog/226}
\end{itemize}

\appendix
\clearpage


	
\end{document}